\documentclass[11pt]{article}

\usepackage[margin=0.9in]{geometry}
\usepackage[T1]{fontenc}
\usepackage{lmodern}
\usepackage{microtype}
\usepackage{amsmath,amssymb,amsthm}
\usepackage{booktabs}
\usepackage{array}
\usepackage{longtable}
\usepackage{algorithm}
\usepackage{algpseudocode}
\usepackage[dvipsnames]{xcolor}
\usepackage{graphicx}
% hyperref is loaded BEFORE tikz on purpose.  hyperref pulls in url, which
% defines \path; tikz defines \path too, and whichever loads last wins.  TikZ
% needs its own (\draw and \node are built on it), so we capture url's verbatim
% behaviour as \filepath first and then let tikz take \path back.
\usepackage{hyperref}
\DeclareUrlCommand\filepath{\urlstyle{tt}}
\usepackage{tikz}
\usetikzlibrary{arrows.meta,positioning,calc}
\usepackage{enumitem}
\usepackage{fancyhdr}
\usepackage[font=small,labelfont=bf,skip=6pt]{caption}

\definecolor{fishgreen}{HTML}{0B5D46}
\definecolor{fishlight}{HTML}{EAF4EF}
\definecolor{fishgray}{HTML}{4B5954}
\hypersetup{colorlinks=true, linkcolor=fishgreen, citecolor=fishgreen, urlcolor=fishgreen}
\setlist{nosep,leftmargin=*}
\setlength{\parindent}{1.2em}
\setlength{\tabcolsep}{4.5pt}
\setlength{\parskip}{0pt plus 1pt}
\pagestyle{fancy}
\fancyhf{}
\lhead{\small FishBot v0.6}
\rhead{\small Tie-breaking, a repaired optimiser, and guarded determinized search}
\cfoot{\thepage}

\newtheorem{theorem}{Theorem}
\newtheorem{proposition}[theorem]{Proposition}
\newtheorem{corollary}[theorem]{Corollary}
\theoremstyle{definition}
\newtheorem{definition}[theorem]{Definition}
\newtheorem{observation}[theorem]{Observation}

% Consistent names for the policies compared.  v0.6 is the policy of this
% report; v0.5, v0.4, v0.3 and v0.2 are the prior generations it is measured
% against.  v0.4-Fast is the policy the v0.4 study shipped and evaluated, and
% v0.4-Block is the same fitted policy with the exact reference belief
% substituted -- the pair the belief-substitution result of \S\ref{sec:belief}
% re-measures.
% These names appear inside captions and section titles, which are moving
% arguments, so the commands must be robust; and they end up in PDF bookmarks,
% so \texorpdfstring supplies a plain spelling for the bookmark text.
\DeclareRobustCommand{\vsixsearch}{\texorpdfstring{\mbox{FishBot v0.6-Search}}{FishBot v0.6-Search}}
\DeclareRobustCommand{\vsix}{\texorpdfstring{\mbox{FishBot v0.6}}{FishBot v0.6}}
\DeclareRobustCommand{\vfive}{\texorpdfstring{\mbox{FishBot v0.5}}{FishBot v0.5}}
\DeclareRobustCommand{\vfour}{\texorpdfstring{\mbox{FishBot v0.4}}{FishBot v0.4}}
\DeclareRobustCommand{\vthree}{\texorpdfstring{\mbox{FishBot v0.3}}{FishBot v0.3}}
\DeclareRobustCommand{\vtwo}{\texorpdfstring{\mbox{FishBot v0.2}}{FishBot v0.2}}
\DeclareRobustCommand{\vfast}{\texorpdfstring{\mbox{v0.4-Fast}}{v0.4-Fast}}
\DeclareRobustCommand{\vblock}{\texorpdfstring{\mbox{v0.4-Block}}{v0.4-Block}}

% SCAFFOLD ONLY.  \todonote marks a heading whose body has not been written and
% whose brief is the \verb|% TODO(v0.6):| comment above it in the section
% source.  It also gives the typesetter breakable material between otherwise
% adjacent headings, which a heading-only document does not have.  Deleting the
% definition is the last step of drafting: once every \todonote is gone the
% command is unused, and if one survives to a release build it is loud on the
% page rather than silent.  Do not use it to stand in for a number -- numbers
% come from numbers_v06.tex.
\newcommand{\todonote}{%
  \par\noindent
  {\small\color{fishgray}\textit{[Scaffold: to be written. The brief for this
   heading is the \texttt{TODO(v0.6)} note in the section source.]}}\par
  \medskip}

\title{\textbf{FishBot v0.6: Breaking the Exact Tie, Repairing the Optimiser,\\
and Guarding a Determinized Search}\\[5pt]
\large Why the exact posterior loses to a policy-aware approximation,\\
what an unguarded search argmax costs, and a measurement protocol\\
that can tell the difference}
\author{Dylan Nguyen\\\small FishLab Research Project}
\date{%
  \small Technical report --- \today\\[2pt]
  \small Repository: \filepath{https://github.com/dylann29/fish-optimization}\\[1pt]
  % TODO(v0.6): replace \commitplaceholder with the short SHA the v0.6
  % experiment battery actually ran at, stated per experiment group if the
  % groups differ (the v0.4 title block states two groups; see
  % research/v06/notes/R5-state-of-knowledge.md e.1).  Add the pointer to
  % research/v06/results/MANIFEST.json once engine/build_manifest.py v06 has
  % been run.
  \small Experiments \texttt{E0--E15}, the follow-ups \texttt{F0--F10} and the
  exploitability probe \texttt{X1} were produced from the \vsix{} working tree
  whose baseline is commit \texttt{bd812fe} (``v0.5''); the \vsix{} sources
  themselves were uncommitted when the battery ran, so the artifacts rather than
  a commit are the record. Digests of all of them are in
  \filepath{research/v06/results/MANIFEST.json}.%
}

%% --- begin numbers_v06 ---
%% numbers_v06.tex -- every number quoted in the v0.6 manuscript.
%%
%% PLACEHOLDER SCAFFOLD.  Nothing in this file has been measured.  Every value
%% is a deliberately impossible literal (XX.XX, XX, X,XXX, XXXXXXX) so that a
%% missing or forgotten regeneration is visible on the page rather than silently
%% plausible.  Do not "fill one in by hand" -- the pipeline is:
%%
%%   engine/experiments_v06.sh  ->  research/v06/results/<artifact>
%%   engine/build_tables_v06.py ->  paper/numbers_v06_generated.tex
%%                                  (\providecommand{}{}\renewcommand{}{value},
%%                                   one source comment per macro)
%%   fishbot_v06.tex            ->  inputs THIS file, then conditionally inputs
%%                                  numbers_v06_generated on top of it
%%
%% (The two file names above are deliberately written without their input
%% commands: paper/inline.py expands any \input it sees, including one inside a
%% comment, so a comment that names its own file recurses to the depth limit and
%% pastes this file into the flattened Overleaf output six times over.)
%%
%% so the generated file always wins and these defaults only reach the PDF when
%% the battery has not been run.  This is the two-file scheme v0.5 introduced
%% (research/v06/notes/R5-state-of-knowledge.md e.6), and it is deliberate that
%% the fallback is ugly.
%%
%% Discipline: no number appears as a literal anywhere in sections_v06/.  If a
%% section needs a figure, it gets a macro here and the macro cites its source
%% artifact in the comment above it.  paper/check.py (a --v06 profile still has
%% to be added at check.py:160-171) enforces the first half of that rule;
%% paper/check_provenance.py enforces the second.
%%
%% Naming convention: \vsix... for a quantity this study measured, \num... for a
%% quantity carried forward from the v0.3/v0.4/v0.5 record or fixed by the
%% protocol.  \providecommand throughout so a partially generated file can be
%% \input first without a clash.
%%
%% Amendment to the header above, made when the background sections were
%% written.  The two classes are no longer both placeholders, and the difference
%% is deliberate.  Every \vsix... macro is a quantity the v0.6 battery has to
%% produce and every one of them is still an impossible literal, so a forgotten
%% regeneration is loud.  A \num... macro carried forward from the v0.4 or v0.5
%% record is a TRANSCRIBED number in the sense of paper/check_provenance.py: the
%% battery will never generate it, it was read out of an artifact that already
%% exists on disk, and leaving it as XX.XX would make the sentence quoting it
%% unreadable rather than visibly unfinished.  Those macros therefore carry their
%% real values, each under a comment naming the artifact it was read from.  A
%% \num... macro whose source is a v0.6 measurement (the comparison rows of
%% sec:corrections, for instance) stays a placeholder like any \vsix... one.

%% ---------------------------------------------------------------------------
%% Title block
%% ---------------------------------------------------------------------------
%% The commit each experiment group ran at.  Fill from `git rev-parse --short HEAD`
%% at the moment experiments_v06.sh is launched; split into two macros if the
%% groups diverge, as the v0.4 title block does.
\providecommand{\commitplaceholder}{XXXXXXX}

%% ===========================================================================
%% sec:rules and app:dialect -- the dialect studied, and the two places the
%% engine is more restrictive than the rules of record allow.  Both channels
%% were measured by the v0.5 study and neither figure is regenerated here.
%% ===========================================================================
%% Turn transfer: research/v05/results/P8-coordination.md sections 1.3 and 1.4,
%% cross-checked in research/v05/results/P8-verify-turn-transfer.md.  A genuine
%% multi-candidate transfer decision arises \numTransferEvents{} times per game;
%% giving one team a ground-truth chooser and pairing on common deals is worth
%% \numTransferDelta{} sets/game.
\providecommand{\numTransferEvents}{0.148}
\providecommand{\numTransferDelta}{-0.0007}
\providecommand{\numTransferCI}{0.0024}
\providecommand{\numTransferGames}{6{,}000}
\providecommand{\numTransferDiverged}{0.48}
%% Declaration arbitration: research/v05/results/P6-verify-arbitration-cost.md
%% sections 3 and 4.  Confidence ranking against lowest-seat arbitration, pooled
%% over five seeds; and the clairvoyant arbitrator that bounds the channel.
\providecommand{\numArbCost}{+0.37}
\providecommand{\numArbRange}{+0.28 to +0.45}
\providecommand{\numArbOracle}{+0.35}
\providecommand{\numArbGames}{30{,}000}
%% The action-cap incidence under the v0.6 battery.  Not yet measured.
\providecommand{\vsixLimitRate}{X.XX}

%% ===========================================================================
%% sec:engine -- verification, information safety, and the v0.6 engine repairs
%% research/v06/results/E1-verify.txt, and the repair regression suite
%% ===========================================================================
\providecommand{\vsixVerifyChecks}{X{,}XXX{,}XXX}
\providecommand{\vsixVerifyViolations}{XX}
\providecommand{\vsixVerifyGames}{XXX}
\providecommand{\vsixVerifyDeterminism}{XXXX}
\providecommand{\vsixVerifyPool}{XX}
\providecommand{\vsixVerifyPairs}{XX}
%% BlockDP thread_local table aliasing (R5 d.3, R0 sec.7 row 1): the mismatch
%% count before the repair and after it, on the same probe.
\providecommand{\vsixAliasChecksBefore}{XXX}
\providecommand{\vsixAliasMismatchBefore}{XXX}
\providecommand{\vsixAliasChecksAfter}{XXX}
\providecommand{\vsixAliasMismatchAfter}{XX}
%% Duplicate-certificate accumulation in Knowledge::onEvent (belief.hpp:168-171)
\providecommand{\vsixCertStoreBefore}{XX}
\providecommand{\vsixCertStoreAfter}{XX}
\providecommand{\vsixBlockCostBefore}{XX.XX}
\providecommand{\vsixBlockCostAfter}{XX.XX}
%% BruteForce::enumerate uninitialised stack at nU==0 (oracle.hpp:102,108-113)
\providecommand{\vsixQzeroShare}{XX.XX}
%% Oracle coverage: the exact oracle only verifies states small enough to
%% brute-force, and the skipped ones are the high-entropy early game
\providecommand{\vsixOracleEnumerated}{XX{,}XXX}
\providecommand{\vsixOracleSkipped}{XX{,}XXX}
\providecommand{\vsixOracleCoverage}{XX.X}
\providecommand{\vsixThroughput}{XXX}
\providecommand{\vsixThroughputVfive}{XXX}

%% ===========================================================================
%% sec:belief -- the observer-conditioned posterior, and the result that the
%% exact posterior is a WORSE predictor than the policy-aware approximation
%% research/v06/results/E?-belief.txt; replicates R6 sec.2b at two seed banks
%% ===========================================================================
\providecommand{\vsixBlockWin}{XX.XX}
\providecommand{\vsixBlockCI}{XX.XX--XX.XX}
\providecommand{\vsixBlockWinB}{XX.XX}
\providecommand{\vsixBlockCIB}{XX.XX--XX.XX}
\providecommand{\vsixBlockSeedA}{XXXXXX}
\providecommand{\vsixBlockSeedB}{XXXXXX}
\providecommand{\vsixBlockGames}{X{,}XXX}
\providecommand{\vsixBlockDeficit}{XX.X}
%% resolution vs calibration: the exact posterior breaks more ties and predicts
%% hits worse.  Both halves must be reported together (R0 sec.6 item 1).
\providecommand{\vsixTieRateFast}{XX.XX}
\providecommand{\vsixTieRateBlock}{XX.XX}
\providecommand{\vsixTieBroken}{XX.X}
\providecommand{\vsixHitFast}{XX.XX}
\providecommand{\vsixHitBlock}{XX.XX}
\providecommand{\vsixDecisionsFast}{X{,}XXX}
\providecommand{\vsixDecisionsBlock}{X{,}XXX}
%% the policy prior itself
\providecommand{\vsixPriorTheta}{X.XXXXX}
\providecommand{\vsixPriorPhi}{X.XXXXX}
\providecommand{\vsixPriorThetaEff}{X.XXXXX}
\providecommand{\vsixPriorDeleteCost}{X.XX}
\providecommand{\vsixPriorDeleteCI}{X.XX, X.XX}
%% the exact count law
\providecommand{\vsixCountStates}{XX{,}XXX}
\providecommand{\vsixCountCost}{X.XXXX}
\providecommand{\vsixSinkhornCost}{XX.X}
\providecommand{\vsixCollapseEvents}{XX.X}

%% ===========================================================================
%% sec:diagnosis -- what v0.5 does wrong, and the two negative results
%% research/v06/results/D?-*.txt (the diagnostic battery)
%% ===========================================================================
%% -- the tie census (R10 sec.4, probe/margin.cpp) ---------------------------
\providecommand{\vsixTieDecisions}{XX{,}XXX}
\providecommand{\vsixTieShare}{XX.XX}
\providecommand{\vsixTieForcedShare}{X.XX}
\providecommand{\vsixTieCandidates}{XX.X}
\providecommand{\vsixTieSpread}{X.XX}
\providecommand{\vsixTieGapMedian}{X.XXXXXX}
\providecommand{\vsixTieSameSuit}{XX.XX}
\providecommand{\vsixTieGroundTruthDiffer}{XX.XX}
\providecommand{\vsixTieDegenerate}{X.XX}
\providecommand{\vsixTieMaxPHit}{XX.XX}
%% -- negative result 1: exchangeable ties are IRREDUCIBLE by rollout --------
\providecommand{\vsixTieOnlyWinLow}{XX.XX}
\providecommand{\vsixTieOnlyCILow}{XX.XX--XX.XX}
\providecommand{\vsixTieOnlyWinHigh}{XX.XX}
\providecommand{\vsixTieOnlyCIHigh}{XX.XX--XX.XX}
\providecommand{\vsixTieOnlyDetLow}{XX}
\providecommand{\vsixTieOnlyDetHigh}{XX}
\providecommand{\vsixTieOnlyN}{XXX}
\providecommand{\vsixTieSymmetricShare}{XX.X}
%% -- negative result 2: v0.5's fit was sampling noise (R4) ------------------
\providecommand{\vsixFitTraceGens}{XX}
\providecommand{\vsixFitIncumbentStart}{X.XXXX}
\providecommand{\vsixFitIncumbentEnd}{X.XXXX}
\providecommand{\vsixFitIncumbentSd}{X.XXXX}
\providecommand{\vsixFitSlope}{X.XXXXX}
\providecommand{\vsixFitSlopeSE}{X.XXXXX}
\providecommand{\vsixFitSlopeT}{X.XX}
\providecommand{\vsixFitBestGap}{X.XXXX}
\providecommand{\vsixFitOrderStat}{X.XXXX}
\providecommand{\vsixFitPop}{XX}
\providecommand{\vsixClipWorst}{XX.X}
\providecommand{\vsixClipWorstKnob}{XXXXXXXX}
\providecommand{\vsixClipFrozenKnob}{XXXXXXXX}
%% -- the loss channels v0.5 still leaks through ----------------------------
\providecommand{\vsixOwnLocked}{X.XX}
\providecommand{\vsixOwnLockedShareOfMiss}{XX.XX}
\providecommand{\vsixAllocDelay}{X.XXX}
\providecommand{\vsixDeclWrong}{X.XX}
\providecommand{\vsixOracleBound}{X.XX}
\providecommand{\vsixOracleBoundWin}{XX.XX}
\providecommand{\vsixTurnValueMid}{X.XX}
\providecommand{\vsixTurnValueLate}{X.XX}

%% ===========================================================================
%% sec:mechanisms -- the v0.6 mechanism set, each keyed to a measured defect
%% research/v06/results/E5-ablations.jsonl (frozen-policy, paired)
%% ===========================================================================
\providecommand{\vsixMechN}{XX}
\providecommand{\vsixAblDeals}{XXX}
\providecommand{\vsixAblSeed}{XXXXXX}
\providecommand{\vsixAblGames}{X{,}XXX}
\providecommand{\vsixAblControl}{XX.XX}
\providecommand{\vsixAblControlCI}{XX.XX--XX.XX}
%% M1 -- per-card / per-target ask dimension (breaks the exact tie)
\providecommand{\vsixMoneDelta}{+X.XX}
\providecommand{\vsixMoneCI}{+X.XX, +X.XX}
\providecommand{\vsixMoneTieAfter}{X.XX}
%% M2 -- harness repair (measured in sec:fitting, gated here)
\providecommand{\vsixMtwoClipAfter}{XX.X}
%% M3 -- shared common-knowledge object plus per-seat refinement
\providecommand{\vsixMthreeDelta}{+X.XX}
\providecommand{\vsixMthreeCI}{+X.XX, +X.XX}
\providecommand{\vsixMthreeOwnLocked}{X.XX}
\providecommand{\vsixMthreeUnsoundFires}{XX}
\providecommand{\vsixMthreeUnsoundChecks}{X{,}XXX{,}XXX}
%% M4 -- memoised proposeDeclaration
\providecommand{\vsixMfourSpeedup}{X.XX}
%% M5 -- local best response / exploitability auditor
\providecommand{\vsixLbrWin}{XX.XX}
\providecommand{\vsixLbrCI}{XX.XX--XX.XX}
\providecommand{\vsixLbrVfive}{XX.XX}
\providecommand{\vsixLbrVfour}{XX.XX}
%% M6 -- multi-valued leaf states
\providecommand{\vsixMsixDelta}{+X.XX}
%% M8 -- recalibrated pAlloc and declaration thresholds
\providecommand{\vsixMeightDelta}{+X.XX}
\providecommand{\vsixMeightECEBefore}{X.XXXXX}
\providecommand{\vsixMeightECEAfter}{X.XXXXX}
%% M9 -- joint allocation naming
\providecommand{\vsixMnineDelta}{+X.XX}
%% M10 -- the three sign-broken ask weights
\providecommand{\vsixSignBroken}{X}
\providecommand{\vsixMtenDelta}{+X.XX}
%% M11 -- offensive void value
\providecommand{\vsixMelevenDelta}{+X.XX}
%% M12 -- per-seat policy specs and the partner-regime split
\providecommand{\vsixMtwelveBotDelta}{+X.XX}
\providecommand{\vsixMtwelveHumanDelta}{+X.XX}

%% ===========================================================================
%% sec:search -- determinized information-set search, the optimizer's curse,
%% and the paired deviation rule
%% research/v06/results/E?-search.jsonl (extends R11)
%% ===========================================================================
\providecommand{\vsixSearchUnguarded}{XX.XX}
\providecommand{\vsixSearchUnguardedCI}{XX.XX--XX.XX}
\providecommand{\vsixSearchControl}{XX.XX}
\providecommand{\vsixSearchControlCI}{XX.XX--XX.XX}
\providecommand{\vsixCurseGap}{XX}
\providecommand{\vsixLeafSd}{X.X}
\providecommand{\vsixKappaSweep}{X, X, X.X, X, X}
\providecommand{\vsixKappaStar}{X.X}
\providecommand{\vsixSearchGuarded}{XX.XX}
\providecommand{\vsixSearchGuardedCI}{XX.XX--XX.XX}
\providecommand{\vsixSearchDet}{XX}
\providecommand{\vsixSearchCand}{X}
\providecommand{\vsixSearchN}{XXX}
\providecommand{\vsixSearchSeed}{XXXXX}
\providecommand{\vsixRolloutCost}{X.XX}
\providecommand{\vsixEventCost}{X.XX}
\providecommand{\vsixDecisionBudget}{XX}
\providecommand{\vsixBranching}{XX.X}
\providecommand{\vsixLiveBranching}{XX.X}
\providecommand{\vsixSearchThroughput}{XX.X}

%% ===========================================================================
%% sec:fitting -- the repaired optimiser
%% research/v06/runs/fit-round?.jsonl, research/v06/runs/roundtrip-assert.txt
%% ===========================================================================
\providecommand{\vsixFitDims}{XX}
\providecommand{\vsixFitAskFeats}{XX}
\providecommand{\vsixFitKnobs}{XX}
\providecommand{\vsixSigmaRule}{X.XX}
\providecommand{\vsixSigmaFloor}{X.XX}
\providecommand{\vsixClipGateThresh}{XX}
\providecommand{\vsixPairedWidth}{X.XXX}
\providecommand{\vsixUnpairedWidth}{X.XXX}
\providecommand{\vsixPairedRatio}{X.X}
\providecommand{\vsixPairedVarianceRatio}{XX}
\providecommand{\vsixPairedSelfDelta}{X.XXXX}
\providecommand{\vsixRegretObjective}{XX.XX}
\providecommand{\vsixSoftMinRatio}{X.XX}
\providecommand{\vsixSoftMinBias}{XX.X}
\providecommand{\vsixFitRounds}{X}
\providecommand{\vsixFitGens}{XX}
\providecommand{\vsixFitBestGen}{XX}
\providecommand{\vsixFitScore}{X.XXXX}
\providecommand{\vsixFitSeed}{XXXXXX}
\providecommand{\vsixFitDeals}{XXX}
\providecommand{\vsixFitPanelN}{X}
\providecommand{\vsixFitPanel}{XXXX, XXXX, XXXX, XXXX, XXXX, XXXX}
\providecommand{\vsixFitTotalGames}{X{,}XXX{,}XXX}
\providecommand{\vsixFitWorst}{XX.XX}
\providecommand{\vsixFitSanityRecovered}{XX}
\providecommand{\vsixFitCurse}{X.XXXX}
\providecommand{\vsixRoundTrip}{XX.XXXX}
\providecommand{\vsixRoundTripPass}{XXXXXX}
%% commit-gate KPIs (the pathology gate; `fish pathology --gate`)
\providecommand{\vsixGateDeadAsk}{X.XX}
\providecommand{\vsixGateDeadRun}{X}
\providecommand{\vsixGateRepeatAsk}{X.XX}
\providecommand{\vsixGateDeclWrong}{X.XX}
\providecommand{\vsixGateOwnLocked}{X.XX}
\providecommand{\vsixGateStarved}{X.XX}
\providecommand{\vsixGateN}{X}

%% ===========================================================================
%% sec:protocol -- six-rotation duplicate blocks, deal-clustered bootstrap,
%% seed separation, paired estimators
%% ===========================================================================
\providecommand{\vsixRotations}{X}
\providecommand{\vsixBootReps}{XX{,}XXX}
\providecommand{\vsixBootAlpha}{XX}
\providecommand{\vsixHeldOutBanks}{X}
\providecommand{\vsixFitBanks}{X}
\providecommand{\vsixReservedSeeds}{XX}
\providecommand{\vsixMirrorDuplication}{XX}

%% ===========================================================================
%% sec:results
%% ===========================================================================
\providecommand{\vsixMirrorWin}{XX.XX}
\providecommand{\vsixMirrorCI}{XX.XX--XX.XX}
\providecommand{\vsixMirrorGames}{X{,}XXX}
\providecommand{\vsixHeadWin}{XX.XX}
\providecommand{\vsixHeadCI}{XX.XX--XX.XX}
\providecommand{\vsixHeadRange}{XX.XX--XX.XX}
\providecommand{\vsixHeadSeeds}{X}
\providecommand{\vsixHeadDeals}{XXX}
\providecommand{\vsixHeadGames}{X{,}XXX}
\providecommand{\vsixHeadPaired}{+X.XX}
\providecommand{\vsixHeadPairedCI}{+X.XX, +X.XX}
\providecommand{\vsixHeadSetsDelta}{+X.XX}
%% the twelve-style panel (E4's nine styles + E10's three deception archetypes),
%% reported as one table with an explicit worst case and minimax regret
%% (R5 f.1: never headline an aggregate)
\providecommand{\vsixStyleN}{XX}
\providecommand{\vsixStyleDeals}{XXX}
\providecommand{\vsixStyleSeed}{XXXXXX}
\providecommand{\vsixStyleGames}{X{,}XXX}
\providecommand{\vsixWorst}{XX.XX}
\providecommand{\vsixWorstOpp}{XXXXXX}
\providecommand{\vsixRegret}{X.XX}
\providecommand{\vsixRegretOpp}{XXXXXX}
\providecommand{\vsixMean}{XX.XX}
\providecommand{\numFiveWorst}{XX.XX}
\providecommand{\numFiveWorstOpp}{XXXXXX}
\providecommand{\numFiveRegret}{X.XX}
\providecommand{\numFiveMean}{XX.XX}
\providecommand{\numFourWorst}{XX.XX}
\providecommand{\numFourRegret}{X.XX}
\providecommand{\numFourMean}{XX.XX}
%% calibration
\providecommand{\vsixCalibAskECE}{X.XXXXX}
\providecommand{\vsixCalibAskBrier}{X.XXXXX}
\providecommand{\vsixCalibAskN}{XX{,}XXX}
\providecommand{\vsixCalibDeclECE}{X.XXXXX}
\providecommand{\vsixCalibDeclBrier}{X.XXXXX}
\providecommand{\vsixCalibDeclN}{X{,}XXX}
%% rules-dialect robustness
\providecommand{\vsixDialectN}{X}
\providecommand{\vsixDialectSpan}{XX.XX--XX.XX}
\providecommand{\vsixDialectSeed}{XXXXXX}
\providecommand{\vsixDialectDeals}{XXX}

%% ===========================================================================
%% sec:corrections -- corrections to the v0.4 and v0.5 records
%% research/v06/notes/R5-state-of-knowledge.md a.4 (U1-U12) and
%% research/v06/notes/R0-OPPORTUNITY-REGISTER.md sec.6
%% ===========================================================================
\providecommand{\vsixCorrN}{XX}
%% U4: the fitted adversary DID reach 50% against v0.4
\providecommand{\numVfourLbrWin}{XX.XX}
\providecommand{\numVfourLbrCI}{XX.XX--XX.XX}
%% U5: value-based declaration was NOT worth measurable win rate
\providecommand{\numVfourVdeclDelta}{+X.XX}
\providecommand{\numVfourVdeclCI}{-X.XX, +X.XX}
%% U8/U9: residue adjudication and the forced-endgame ladder.  Transcribed from
%% research/v05/results/C1-v04-corrections.md entries C7 and C9, which take them
%% in turn from research/v05/results/P4-verify-forcing-horizon.md section 5 and
%% research/v05/results/P2-forced-endgame.md section 5.  Neither figure was
%% printed by the v0.5 paper.
\providecommand{\numResidueHeld}{4.08}
\providecommand{\numLadderLiveSets}{1}
\providecommand{\numLadderOccasions}{566}
\providecommand{\numLadderSeed}{31}
%% C4/defect J: the soft minimum is a weighted mean
\providecommand{\numSoftMinBetaFour}{XX}
\providecommand{\numSoftMinBetaFive}{XX}
\providecommand{\numSoftMinRatioFour}{1.88}
\providecommand{\numSoftMinBelowMin}{11.2}
%% M7's calibration point drifted between v0.4 and v0.5
\providecommand{\numThetaEffFour}{X.XXXXX}
\providecommand{\numThetaEffFive}{X.XXXXX}
\providecommand{\numThetaEffDrift}{XX}
%% gateaudit is dead code for v0.5
\providecommand{\numGateAuditOpps}{X}
\providecommand{\numGateAuditVfourFalseNeg}{XXX}
%% defect H: stale computeAggregates is real but inert
\providecommand{\numDefectHDiffers}{XX.XX}
\providecommand{\numDefectHFlips}{X.XX}
%% carried forward from the v0.4 and v0.5 records for the comparison rows
\providecommand{\numVfiveHeadWin}{XX.XX}
\providecommand{\numVfiveHeadRange}{XX.XX--XX.XX}
\providecommand{\numVfourVsVthree}{XX.XX}

%% ===========================================================================
%% sec:limitations and app:repro -- provenance of the manuscript's own figures,
%% computed by engine/build_tables_v06.py over paper/sections_v06/ so the
%% paragraph that quotes them cannot itself go stale.  Audited by
%% paper/check_provenance.py.
%% ===========================================================================
\providecommand{\vsixProvUsed}{XXX}
\providecommand{\vsixProvGenerated}{XXX}
\providecommand{\vsixProvTranscribed}{XX}
\providecommand{\vsixManifestFiles}{XX}
\providecommand{\vsixWallClock}{XX}
\providecommand{\vsixCores}{XX}

%% ===========================================================================
%% ADDENDUM, added while sec:engine, sec:belief, sec:fitting, sec:protocol,
%% app:inference, app:params and app:repro were drafted.  Same two-class rule as
%% the header states: \vsix... is a quantity engine/build_tables_v06.py emits
%% from a v0.6 artifact and is an impossible literal until the battery has run;
%% \num... is a transcribed quantity read out of a file that already exists on
%% disk, and carries its real value under a comment naming that file.
%% ===========================================================================

%% ---------------------------------------------------------------------------
%% E0 identity controls and the BlockDP aliasing repair (engine/experiments_v06.sh
%% step E0 -> research/v06/results/E0-identity.txt).  vsixIdentity is the md5
%% comparison of the full `fish pathology` output between v05 and v06:legacy=1;
%% the two alias macros separate the QUERY check (must be 0) from the raw
%% shared-pool field read (documented, not fixed).
%% ---------------------------------------------------------------------------
\providecommand{\vsixIdentity}{XXXX}
\providecommand{\vsixAliasQueryMismatch}{XX}
\providecommand{\vsixAliasRawMismatch}{XXX}
%% The pre-repair measurement of the same defect is a v0.5-record number that the
%% v0.6 battery will never regenerate: research/v05/results/P2-forced-endgame.md
%% section 6, `fish blockalias --a=v04 --games=40 --seed=31`.
\providecommand{\numAliasChecks}{294}
\providecommand{\numAliasMismatch}{285}

%% ---------------------------------------------------------------------------
%% E1 verify (research/v06/results/E1-verify.txt).  The generator emits the
%% audit counters and the overall verdict under these names; the pool size and
%% the ordered-pair count are fixed by engine/src/main.cpp and are transcribed.
%% ---------------------------------------------------------------------------
\providecommand{\vsixAuditViolations}{XX}
\providecommand{\vsixAuditChecks}{X{,}XXX{,}XXX}
\providecommand{\vsixVerify}{XXXX}
%% engine/src/main.cpp, the `verify` command's policy pool
\providecommand{\numVerifyPool}{nine}
\providecommand{\numVerifyPairs}{81}

%% ---------------------------------------------------------------------------
%% E10 declaration pre-gate audit, enabled for v0.5 and v0.6 by this study
%% (research/v06/results/E10-gateaudit.txt)
%% ---------------------------------------------------------------------------
\providecommand{\vsixGateOpportunities}{XX{,}XXX}
\providecommand{\vsixGateRejected}{X{,}XXX}
\providecommand{\vsixGateFalseNeg}{XX}
\providecommand{\vsixGateVerdict}{XXXX}

%% ---------------------------------------------------------------------------
%% E9 throughput (research/v06/results/E9-throughput.txt), games/s
%% ---------------------------------------------------------------------------
\providecommand{\vsixThroughputSix}{XXX.X}
\providecommand{\vsixThroughputFive}{XXX.X}
\providecommand{\vsixThroughputFour}{XXX.X}

%% ---------------------------------------------------------------------------
%% E8 belief: each posterior scored as a PREDICTOR on the same states, by
%% predictive log loss over the unresolved cards and by the realised hit rate of
%% its own argmax ask (research/v06/results/E8-belief.txt, rendered as
%% paper/tables_v06/belief.tex).
%% ---------------------------------------------------------------------------
\providecommand{\vsixBeliefExactNLL}{X.XXXXX}
\providecommand{\vsixBeliefExactHit}{XX.XX}
\providecommand{\vsixBeliefNoPriorNLL}{X.XXXXX}
\providecommand{\vsixBeliefNoPriorHit}{XX.XX}
\providecommand{\vsixBeliefShippedNLL}{X.XXXXX}
\providecommand{\vsixBeliefShippedHit}{XX.XX}
%% E8 ties, full live candidate set rather than the tie subset
%% (research/v06/results/E8-ties.txt, the v0.5-mirror block)
\providecommand{\vsixFullFastFive}{XX.XX}
\providecommand{\vsixFullExactFive}{XX.XX}
\providecommand{\vsixFullBestFive}{XX.XX}
\providecommand{\vsixDisagreeFive}{XX.XX}

%% ---------------------------------------------------------------------------
%% The policy prior's two operating points, transcribed from
%% research/v06/notes/R12-mechanism-trials.md section 3.  The battery re-measures
%% the predictive column but not these play-side costs, so they are \num...
%% ---------------------------------------------------------------------------
\providecommand{\numThetaShipped}{0.44458}
\providecommand{\numThetaPredictiveOpt}{0.60--0.80}
\providecommand{\numThetaPriorHitGain}{2.16}
\providecommand{\numThetaPriorNatGain}{0.011}
\providecommand{\numThetaSixtyCost}{4.4}
\providecommand{\numThetaElevenCost}{7.1}
\providecommand{\numThetaTrialN}{1{,}200}

%% ---------------------------------------------------------------------------
%% Parameter-vector shape.  Code constants, transcribed from the headers named.
%% ---------------------------------------------------------------------------
%% engine/src/v04.hpp, NFEAT
\providecommand{\numAskFeats}{20}
%% engine/src/main.cpp, the fourteen knobs appended by --full
\providecommand{\numVfiveKnobs}{14}
%% engine/src/v06.hpp, the three v0.6 ask terms
\providecommand{\numVsixAskTerms}{3}
%% engine/src/tuner.hpp and engine/src/main.cpp: the v0.5 vector
\providecommand{\numVfiveParams}{34}
%% engine/src/v06.hpp, NV6PARAM = NFEAT + 14 + 3
\providecommand{\numVsixParams}{37}

%% ---------------------------------------------------------------------------
%% The four fitting defects, transcribed from
%% research/v06/notes/R4-harness-and-fitting.md sections 2 and 3.  The battery
%% does not re-derive them: they are properties of the v0.5 fit trace.
%% ---------------------------------------------------------------------------
\providecommand{\numSigmaScalar}{0.6}
\providecommand{\numSigmaRangeMin}{0.10}
\providecommand{\numSigmaRangeMax}{40}
\providecommand{\numClipDeclMargin}{93.4}
\providecommand{\numValueWeightStep}{1.5}
\providecommand{\numFitTraceGensVfive}{40}

%% ---------------------------------------------------------------------------
%% The harness acceptance tests.  The self-comparison is an invariant, not a
%% measurement.  The handicapped-base recovery run is an acceptance criterion
%% stated in research/v06/DESIGN.md (E4) and research/v06/notes/R4-harness-and-fitting.md
%% (H2); it is run by hand from the command in app:repro-commands and is NOT one
%% of engine/experiments_v06.sh's blocks, so engine/build_tables_v06.py does not
%% regenerate these and they are transcribed rather than generated.
%% ---------------------------------------------------------------------------
\providecommand{\numFitSanityGens}{10}
\providecommand{\numFitSanityBase}{45.89}
\providecommand{\numFitSanityAfter}{48.33}
\providecommand{\numFitSanityVthree}{76.83}

%% ---------------------------------------------------------------------------
%% The resolution requirement (sec:protocol).  Five mechanisms that cleared a
%% 95% paired interval at 400-1,600 games per cell and returned exactly zero at
%% the larger budget on two disjoint banks.  All transcribed from
%% research/v06/notes/R12-mechanism-trials.md section 6.1.  Sign convention is
%% the corpus's: delta = reference - variant, so a NEGATIVE delta means the
%% variant is better.
%% ---------------------------------------------------------------------------
\providecommand{\numPlateauMechs}{Five}
\providecommand{\numResVoidSmall}{-0.0292}
\providecommand{\numResVoidSmallCI}{-0.0537, -0.0037}
\providecommand{\numResVoidSmallN}{400}
\providecommand{\numResVoidLargeA}{+0.0007}
\providecommand{\numResVoidLargeACI}{-0.0213, +0.0227}
\providecommand{\numResVoidLargeB}{+0.0200}
\providecommand{\numResVoidLargeBCI}{-0.0017, +0.0413}
\providecommand{\numResVoidLargeN}{1{,}000}
\providecommand{\numResVoidThreeSmall}{-0.0313}
\providecommand{\numResVoidThreeSmallCI}{-0.0593, -0.0033}
\providecommand{\numResVoidThreeSmallN}{500}
\providecommand{\numResPerseverSmall}{-0.0200}
\providecommand{\numResPerseverSmallCI}{-0.0356, -0.0044}
\providecommand{\numResPerseverSmallN}{800}
\providecommand{\numResPerseverLargeA}{-0.0012}
\providecommand{\numResPerseverLargeACI}{-0.0137, +0.0112}
\providecommand{\numResPerseverLargeB}{+0.0035}
\providecommand{\numResPerseverLargeBCI}{-0.0092, +0.0158}
\providecommand{\numResPerseverLargeN}{1{,}000}
\providecommand{\numResDeclSmall}{+1.75}
\providecommand{\numResDeclSmallN}{1{,}200}
\providecommand{\numResDeclLarge}{+0.0133}
\providecommand{\numResDeclLargeCI}{+0.0013, +0.0247}
\providecommand{\numResVoiSmall}{-0.0053}
\providecommand{\numResVoiSmallCI}{-0.0244, +0.0141}
\providecommand{\numResVoiSmallN}{800}
\providecommand{\numResLargeBudget}{3{,}000}

%% ---------------------------------------------------------------------------
%% Fit traces, research/v06/runs/fitA.jsonl and research/v06/runs/fitC.jsonl.
%% engine/build_tables_v06.py reads the header record and the generation records
%% and \renewcommand's every one of these; the placeholders exist so the
%% manuscript compiles when the runs directory is absent.
%% ---------------------------------------------------------------------------
\providecommand{\vsixFitAGens}{XX}
\providecommand{\vsixFitAPanel}{XXXX+XXXX+XXXX+XXXX}
\providecommand{\vsixFitASeed}{XXXXXXXX}
\providecommand{\vsixFitAPop}{XX}
\providecommand{\vsixFitADeals}{XXXxX}
\providecommand{\vsixFitASigmaRel}{X.XXX}
\providecommand{\vsixFitAClipGenZero}{XX.XX}
\providecommand{\vsixFitCGens}{XX}
\providecommand{\vsixFitCPanel}{XXXX+XXXX+XXXX+XXXX}
\providecommand{\vsixFitCSeed}{XXXXXXXX}
\providecommand{\vsixFitCPop}{XX}
\providecommand{\vsixFitCDeals}{XXXxX}
\providecommand{\vsixFitCSigmaRel}{X.XXX}
\providecommand{\vsixFitCClipGenZero}{XX.XX}

%% ---------------------------------------------------------------------------
%% Macros quoted by sec:diagnosis and sec:search that had no placeholder.  Added
%% here so the manuscript compiles; the generator emits the E8-, E2-, E12- and
%% E15-derived ones under exactly these names.
%% ---------------------------------------------------------------------------
\providecommand{\vsixAblateRef}{XX.XX}
\providecommand{\vsixTieNFive}{X{,}XXX}
\providecommand{\vsixTieShareFive}{XX.XX}
\providecommand{\vsixTieSameCardFive}{X.XX}
\providecommand{\vsixTieSameSuitFive}{XX.XX}
\providecommand{\vsixTieSpreadFive}{X.XXXXX}
\providecommand{\vsixTieGapFive}{X.XXXXX}
\providecommand{\vsixTieExactSepPctFive}{X.XX}
\providecommand{\vsixTieShipMovedFive}{XX.XX}
\providecommand{\vsixTieHitArrayFive}{XX.XX}
\providecommand{\vsixTieHitFastFive}{XX.XX}
\providecommand{\vsixTieHitShipFive}{XX.XX}
\providecommand{\vsixTieHitExactFive}{XX.XX}
\providecommand{\vsixTieHitBestFive}{XX.XX}
\providecommand{\vsixPathFiveContested}{XX.XX}
\providecommand{\vsixPathFiveDeclwrong}{X.XX}
\providecommand{\vsixPathFiveOwnlocked}{X.XX}
\providecommand{\vsixSearchDead}{XX.XX}
\providecommand{\vsixSearchSpread}{XX.XX}
\providecommand{\vsixSearchDevTwoFive}{XX.XX}
\providecommand{\vsixSearchGps}{XX.X}
\providecommand{\vsixDeadWin}{XX.XX}
\providecommand{\vsixDeadWinLo}{XX.XX}
\providecommand{\vsixDeadWinHi}{XX.XX}
\providecommand{\vsixDeadPathDead}{XX.XX}
\providecommand{\vsixDeadPathLongest}{XXX}
\providecommand{\vsixDeadPathLimit}{X.XX}
%% The two swept theta settings quoted in sec:belief-exactworse, from the same
%% source (research/v06/notes/R12-mechanism-trials.md section 3).
\providecommand{\numThetaTrialLow}{0.60}
\providecommand{\numThetaTrialHigh}{1.10}
%% The small-budget range the five non-replicating mechanisms were first run at
%% (research/v06/notes/R12-mechanism-trials.md, headline and section 6.1).
\providecommand{\numResSmallBudgetLo}{400}
\providecommand{\numResSmallBudgetHi}{1{,}600}

%% ---------------------------------------------------------------------------
%% Compile guard, not a number.  sec:mechanisms writes \textmu, a TS1 symbol
%% that this preamble does not load textcomp for, so it is undefined and halts
%% the build.  \providecommand supplies a math fallback and does nothing if a
%% later preamble change loads textcomp properly.
%% ---------------------------------------------------------------------------
\providecommand{\textmu}{\ensuremath{\mu}}

%% ---------------------------------------------------------------------------
%% Further placeholders for macros quoted by sec:mechanisms and sec:results.
%% \providecommand is idempotent, so these are harmless if the same names are
%% also declared above or by a later edit; they exist so that the manuscript
%% compiles before engine/build_tables_v06.py has run.
%% E2-pathology.txt, per arm; E4-perstyle.jsonl; E3-headtohead.jsonl.
%% ---------------------------------------------------------------------------
\providecommand{\vsixExactBuildUsFive}{XX.X}
\providecommand{\vsixPathSixEvents}{XXX.XX}
\providecommand{\vsixPathSixHit}{XX.XX}
\providecommand{\vsixPathSixDead}{X.XX}
\providecommand{\vsixPathSixLongest}{XX}
\providecommand{\vsixPathSixRunSix}{X.XX}
\providecommand{\vsixPathSixRepeatACT}{X.XX}
\providecommand{\vsixPathSixDeclwrong}{X.XX}
\providecommand{\vsixPathSixLimit}{X.XX}
\providecommand{\vsixPathFiveEvents}{XXX.XX}
\providecommand{\vsixPathFiveHit}{XX.XX}
\providecommand{\vsixPathFiveDead}{X.XX}
\providecommand{\vsixPathFiveLongest}{XX}
\providecommand{\vsixPathFiveRunSix}{X.XX}
\providecommand{\vsixPathFiveRepeatACT}{X.XX}
\providecommand{\vsixPathFiveLimit}{X.XX}
\providecommand{\vsixPathFourEvents}{XXX.XX}
\providecommand{\vsixPathFourHit}{XX.XX}
\providecommand{\vsixPathFourDead}{X.XX}
\providecommand{\vsixPathFourLongest}{XX}
\providecommand{\vsixPathFourRunSix}{X.XX}
\providecommand{\vsixPathFourRepeatACT}{X.XX}
\providecommand{\vsixPathFourDeclwrong}{X.XX}
\providecommand{\vsixPathFourLimit}{X.XX}
\providecommand{\vsixStyleWorstSix}{XX.XX}
\providecommand{\vsixStyleWorstFive}{XX.XX}
\providecommand{\vsixStyleWorstFour}{XX.XX}
\providecommand{\vsixStyleMeanSix}{XX.XX}
\providecommand{\vsixStyleMeanFive}{XX.XX}
\providecommand{\vsixStyleMeanFour}{XX.XX}
\providecommand{\vsixStyleWorstOppSix}{XXXXXX}
\providecommand{\vsixStyleWorstOppFive}{XXXXXX}
\providecommand{\vsixStyleWorstOppFour}{XXXXXX}
\providecommand{\vsixRegretSix}{X.XX}
\providecommand{\vsixRegretFive}{X.XX}
\providecommand{\vsixRegretFour}{X.XX}
\providecommand{\vsixHeadFiveMean}{XX.XX}
\providecommand{\vsixHeadFiveMin}{XX.XX}
\providecommand{\vsixHeadFiveMax}{XX.XX}
\providecommand{\vsixHeadFiveBanks}{X}
\providecommand{\vsixHeadFiveN}{X{,}XXX}
\providecommand{\vsixHeadFourMean}{XX.XX}
\providecommand{\vsixHeadFourMin}{XX.XX}
\providecommand{\vsixHeadFourMax}{XX.XX}

%% ---------------------------------------------------------------------------
%% app:inference -- the v0.6-era exactness checks and the three latent
%% approximations.  All transcribed from research/v06/notes/R2-inference-stack.md
%% section 8; the v0.6 battery does not re-run them, so they are \num...
%% ---------------------------------------------------------------------------
\providecommand{\numCountUniformQueries}{673}
\providecommand{\numCountGeneralQueries}{1{,}633}
\providecommand{\numCountAgree}{1.7\times10^{-15}}
\providecommand{\numAllocSurvivalChecks}{35{,}161}
\providecommand{\numAllocSurvivalCert}{10{,}374}
\providecommand{\numAllocSurvivalFailures}{0}
\providecommand{\numAllocSurvivalDeals}{40}
\providecommand{\numAllocSurvivalSeed}{4242}
%% engine/src/blockdp.hpp: the three capacity constants behind the truncation
\providecommand{\numMaxEnt}{512}
\providecommand{\numMaxCountVectors}{462}
\providecommand{\numCertCap}{16}
\providecommand{\numNodeGuard}{2{,}000{,}000}
%% --- end numbers_v06 ---
% Generated from the experiment artifacts by engine/build_tables_v06.py, on the
% two-file scheme v0.5 introduced: the generator \renewcommand's every macro the
% artifacts settle, so the placeholder values in numbers_v06.tex act only as a
% compile-before-the-battery-runs fallback and can never silently reach the PDF
% once the battery has been run.  The file does not exist yet; \InputIfFileExists
% keeps the manuscript compiling until it does.
%% --- begin numbers_v06_generated (InputIfFileExists) ---
% GENERATED by engine/build_tables_v06.py -- do not edit.
% Each macro names the artifact and field it came from.

% E0-identity.txt
\providecommand{\vsixIdentity}{}\renewcommand{\vsixIdentity}{PASS}
% E0-identity.txt: fish blockalias
\providecommand{\vsixAliasQueryMismatch}{}\renewcommand{\vsixAliasQueryMismatch}{0}
% E0-identity.txt: fish blockalias
\providecommand{\vsixAliasRawMismatch}{}\renewcommand{\vsixAliasRawMismatch}{175}
% E1-verify.txt
\providecommand{\vsixAuditViolations}{}\renewcommand{\vsixAuditViolations}{0}
% E1-verify.txt
\providecommand{\vsixAuditChecks}{}\renewcommand{\vsixAuditChecks}{23,594,580}
% E1-verify.txt
\providecommand{\vsixVerify}{}\renewcommand{\vsixVerify}{PASS}
% E2-pathology.txt "v0.6 mirror" field events
\providecommand{\vsixPathSixEvents}{}\renewcommand{\vsixPathSixEvents}{94.59}
% E2-pathology.txt "v0.6 mirror" field pNinety
\providecommand{\vsixPathSixPNinety}{}\renewcommand{\vsixPathSixPNinety}{110}
% E2-pathology.txt "v0.6 mirror" field pNineNine
\providecommand{\vsixPathSixPNineNine}{}\renewcommand{\vsixPathSixPNineNine}{120}
% E2-pathology.txt "v0.6 mirror" field hit
\providecommand{\vsixPathSixHit}{}\renewcommand{\vsixPathSixHit}{53.29}
% E2-pathology.txt "v0.6 mirror" field dead
\providecommand{\vsixPathSixDead}{}\renewcommand{\vsixPathSixDead}{0.01}
% E2-pathology.txt "v0.6 mirror" field longest
\providecommand{\vsixPathSixLongest}{}\renewcommand{\vsixPathSixLongest}{1}
% E2-pathology.txt "v0.6 mirror" field runSix
\providecommand{\vsixPathSixRunSix}{}\renewcommand{\vsixPathSixRunSix}{0.00}
% E2-pathology.txt "v0.6 mirror" field repeatACT
\providecommand{\vsixPathSixRepeatACT}{}\renewcommand{\vsixPathSixRepeatACT}{2.22}
% E2-pathology.txt "v0.6 mirror" field ownlocked
\providecommand{\vsixPathSixOwnlocked}{}\renewcommand{\vsixPathSixOwnlocked}{11.69}
% E2-pathology.txt "v0.6 mirror" field declwrong
\providecommand{\vsixPathSixDeclwrong}{}\renewcommand{\vsixPathSixDeclwrong}{2.37}
% E2-pathology.txt "v0.6 mirror" field limit
\providecommand{\vsixPathSixLimit}{}\renewcommand{\vsixPathSixLimit}{0.00}
% E2-pathology.txt "v0.6 mirror" field forced
\providecommand{\vsixPathSixForced}{}\renewcommand{\vsixPathSixForced}{0}
% E2-pathology.txt "v0.6 mirror" field forcedwrong
\providecommand{\vsixPathSixForcedwrong}{}\renewcommand{\vsixPathSixForcedwrong}{0}
% E2-pathology.txt "v0.5 mirror" field events
\providecommand{\vsixPathFiveEvents}{}\renewcommand{\vsixPathFiveEvents}{96.72}
% E2-pathology.txt "v0.5 mirror" field pNinety
\providecommand{\vsixPathFivePNinety}{}\renewcommand{\vsixPathFivePNinety}{112}
% E2-pathology.txt "v0.5 mirror" field pNineNine
\providecommand{\vsixPathFivePNineNine}{}\renewcommand{\vsixPathFivePNineNine}{124}
% E2-pathology.txt "v0.5 mirror" field hit
\providecommand{\vsixPathFiveHit}{}\renewcommand{\vsixPathFiveHit}{55.63}
% E2-pathology.txt "v0.5 mirror" field dead
\providecommand{\vsixPathFiveDead}{}\renewcommand{\vsixPathFiveDead}{0.01}
% E2-pathology.txt "v0.5 mirror" field longest
\providecommand{\vsixPathFiveLongest}{}\renewcommand{\vsixPathFiveLongest}{1}
% E2-pathology.txt "v0.5 mirror" field runSix
\providecommand{\vsixPathFiveRunSix}{}\renewcommand{\vsixPathFiveRunSix}{0.00}
% E2-pathology.txt "v0.5 mirror" field repeatACT
\providecommand{\vsixPathFiveRepeatACT}{}\renewcommand{\vsixPathFiveRepeatACT}{2.61}
% E2-pathology.txt "v0.5 mirror" field ownlocked
\providecommand{\vsixPathFiveOwnlocked}{}\renewcommand{\vsixPathFiveOwnlocked}{9.75}
% E2-pathology.txt "v0.5 mirror" field declwrong
\providecommand{\vsixPathFiveDeclwrong}{}\renewcommand{\vsixPathFiveDeclwrong}{2.07}
% E2-pathology.txt "v0.5 mirror" field limit
\providecommand{\vsixPathFiveLimit}{}\renewcommand{\vsixPathFiveLimit}{0.00}
% E2-pathology.txt "v0.5 mirror" field forced
\providecommand{\vsixPathFiveForced}{}\renewcommand{\vsixPathFiveForced}{6}
% E2-pathology.txt "v0.5 mirror" field forcedwrong
\providecommand{\vsixPathFiveForcedwrong}{}\renewcommand{\vsixPathFiveForcedwrong}{6}
% E2-pathology.txt "v0.4 mirror" field events
\providecommand{\vsixPathFourEvents}{}\renewcommand{\vsixPathFourEvents}{143.60}
% E2-pathology.txt "v0.4 mirror" field pNinety
\providecommand{\vsixPathFourPNinety}{}\renewcommand{\vsixPathFourPNinety}{312}
% E2-pathology.txt "v0.4 mirror" field pNineNine
\providecommand{\vsixPathFourPNineNine}{}\renewcommand{\vsixPathFourPNineNine}{321}
% E2-pathology.txt "v0.4 mirror" field hit
\providecommand{\vsixPathFourHit}{}\renewcommand{\vsixPathFourHit}{34.25}
% E2-pathology.txt "v0.4 mirror" field dead
\providecommand{\vsixPathFourDead}{}\renewcommand{\vsixPathFourDead}{39.04}
% E2-pathology.txt "v0.4 mirror" field longest
\providecommand{\vsixPathFourLongest}{}\renewcommand{\vsixPathFourLongest}{286}
% E2-pathology.txt "v0.4 mirror" field runSix
\providecommand{\vsixPathFourRunSix}{}\renewcommand{\vsixPathFourRunSix}{34.33}
% E2-pathology.txt "v0.4 mirror" field repeatACT
\providecommand{\vsixPathFourRepeatACT}{}\renewcommand{\vsixPathFourRepeatACT}{40.03}
% E2-pathology.txt "v0.4 mirror" field ownlocked
\providecommand{\vsixPathFourOwnlocked}{}\renewcommand{\vsixPathFourOwnlocked}{16.46}
% E2-pathology.txt "v0.4 mirror" field declwrong
\providecommand{\vsixPathFourDeclwrong}{}\renewcommand{\vsixPathFourDeclwrong}{10.44}
% E2-pathology.txt "v0.4 mirror" field limit
\providecommand{\vsixPathFourLimit}{}\renewcommand{\vsixPathFourLimit}{0.00}
% E2-pathology.txt "v0.4 mirror" field forced
\providecommand{\vsixPathFourForced}{}\renewcommand{\vsixPathFourForced}{28}
% E2-pathology.txt "v0.4 mirror" field forcedwrong
\providecommand{\vsixPathFourForcedwrong}{}\renewcommand{\vsixPathFourForcedwrong}{28}
% E2-pathology.txt "v0.6 vs v0.5" field events
\providecommand{\vsixPathSixVFiveEvents}{}\renewcommand{\vsixPathSixVFiveEvents}{95.69}
% E2-pathology.txt "v0.6 vs v0.5" field pNinety
\providecommand{\vsixPathSixVFivePNinety}{}\renewcommand{\vsixPathSixVFivePNinety}{110}
% E2-pathology.txt "v0.6 vs v0.5" field pNineNine
\providecommand{\vsixPathSixVFivePNineNine}{}\renewcommand{\vsixPathSixVFivePNineNine}{122}
% E2-pathology.txt "v0.6 vs v0.5" field hit
\providecommand{\vsixPathSixVFiveHit}{}\renewcommand{\vsixPathSixVFiveHit}{55.59}
% E2-pathology.txt "v0.6 vs v0.5" field dead
\providecommand{\vsixPathSixVFiveDead}{}\renewcommand{\vsixPathSixVFiveDead}{0.02}
% E2-pathology.txt "v0.6 vs v0.5" field longest
\providecommand{\vsixPathSixVFiveLongest}{}\renewcommand{\vsixPathSixVFiveLongest}{1}
% E2-pathology.txt "v0.6 vs v0.5" field runSix
\providecommand{\vsixPathSixVFiveRunSix}{}\renewcommand{\vsixPathSixVFiveRunSix}{0.00}
% E2-pathology.txt "v0.6 vs v0.5" field repeatACT
\providecommand{\vsixPathSixVFiveRepeatACT}{}\renewcommand{\vsixPathSixVFiveRepeatACT}{2.63}
% E2-pathology.txt "v0.6 vs v0.5" field ownlocked
\providecommand{\vsixPathSixVFiveOwnlocked}{}\renewcommand{\vsixPathSixVFiveOwnlocked}{9.01}
% E2-pathology.txt "v0.6 vs v0.5" field declwrong
\providecommand{\vsixPathSixVFiveDeclwrong}{}\renewcommand{\vsixPathSixVFiveDeclwrong}{2.04}
% E2-pathology.txt "v0.6 vs v0.5" field limit
\providecommand{\vsixPathSixVFiveLimit}{}\renewcommand{\vsixPathSixVFiveLimit}{0.00}
% E2-pathology.txt "v0.6 vs v0.5" field forced
\providecommand{\vsixPathSixVFiveForced}{}\renewcommand{\vsixPathSixVFiveForced}{10}
% E2-pathology.txt "v0.6 vs v0.5" field forcedwrong
\providecommand{\vsixPathSixVFiveForcedwrong}{}\renewcommand{\vsixPathSixVFiveForcedwrong}{8}
% E2-pathology.txt v0.6 mirror (alias of the KPI row)
\providecommand{\vsixGateDeadAsk}{}\renewcommand{\vsixGateDeadAsk}{0.01}
% E2-pathology.txt v0.6 mirror (alias of the KPI row)
\providecommand{\vsixGateDeadRun}{}\renewcommand{\vsixGateDeadRun}{1}
% E2-pathology.txt v0.6 mirror (alias of the KPI row)
\providecommand{\vsixGateRepeatAsk}{}\renewcommand{\vsixGateRepeatAsk}{2.22}
% E2-pathology.txt v0.6 mirror (alias of the KPI row)
\providecommand{\vsixGateOwnLocked}{}\renewcommand{\vsixGateOwnLocked}{11.69}
% E2-pathology.txt v0.6 mirror (alias of the KPI row)
\providecommand{\vsixGateDeclWrong}{}\renewcommand{\vsixGateDeclWrong}{2.37}
% E2-pathology.txt v0.6 mirror (alias of the KPI row)
\providecommand{\vsixLimitRate}{}\renewcommand{\vsixLimitRate}{0.00}
% E3-headtohead.jsonl: mean over 5 banks vs v05
\providecommand{\vsixHeadFiveMean}{}\renewcommand{\vsixHeadFiveMean}{51.03}
% E3-headtohead.jsonl: min bank vs v05
\providecommand{\vsixHeadFiveMin}{}\renewcommand{\vsixHeadFiveMin}{50.11}
% E3-headtohead.jsonl: max bank vs v05
\providecommand{\vsixHeadFiveMax}{}\renewcommand{\vsixHeadFiveMax}{51.83}
% E3-headtohead.jsonl: bank count vs v05
\providecommand{\vsixHeadFiveBanks}{}\renewcommand{\vsixHeadFiveBanks}{5}
% E3-headtohead.jsonl: total games
\providecommand{\vsixHeadFiveN}{}\renewcommand{\vsixHeadFiveN}{9,000}
% E3-headtohead.jsonl: mean over 5 banks vs v04
\providecommand{\vsixHeadFourMean}{}\renewcommand{\vsixHeadFourMean}{50.86}
% E3-headtohead.jsonl: min bank vs v04
\providecommand{\vsixHeadFourMin}{}\renewcommand{\vsixHeadFourMin}{50.28}
% E3-headtohead.jsonl: max bank vs v04
\providecommand{\vsixHeadFourMax}{}\renewcommand{\vsixHeadFourMax}{52.22}
% E3-headtohead.jsonl: bank count vs v04
\providecommand{\vsixHeadFourBanks}{}\renewcommand{\vsixHeadFourBanks}{5}
% E3-headtohead.jsonl: total games
\providecommand{\vsixHeadFourN}{}\renewcommand{\vsixHeadFourN}{9,000}
% E4-perstyle.jsonl: min over 13 styles for v06
\providecommand{\vsixStyleWorstSix}{}\renewcommand{\vsixStyleWorstSix}{48.67}
% E4-perstyle.jsonl: mean over 13 styles for v06
\providecommand{\vsixStyleMeanSix}{}\renewcommand{\vsixStyleMeanSix}{78.39}
% E4-perstyle.jsonl: argmin style for v06
\providecommand{\vsixStyleWorstOppSix}{}\renewcommand{\vsixStyleWorstOppSix}{v04}
% E4-perstyle.jsonl: min over 13 styles for v05
\providecommand{\vsixStyleWorstFive}{}\renewcommand{\vsixStyleWorstFive}{50.00}
% E4-perstyle.jsonl: mean over 13 styles for v05
\providecommand{\vsixStyleMeanFive}{}\renewcommand{\vsixStyleMeanFive}{77.57}
% E4-perstyle.jsonl: argmin style for v05
\providecommand{\vsixStyleWorstOppFive}{}\renewcommand{\vsixStyleWorstOppFive}{v05}
% E4-perstyle.jsonl: min over 13 styles for v04
\providecommand{\vsixStyleWorstFour}{}\renewcommand{\vsixStyleWorstFour}{48.39}
% E4-perstyle.jsonl: mean over 13 styles for v04
\providecommand{\vsixStyleMeanFour}{}\renewcommand{\vsixStyleMeanFour}{76.77}
% E4-perstyle.jsonl: argmin style for v04
\providecommand{\vsixStyleWorstOppFour}{}\renewcommand{\vsixStyleWorstOppFour}{v05}
% E4-perstyle.jsonl: max regret over the style set for v06
\providecommand{\vsixRegretSix}{}\renewcommand{\vsixRegretSix}{3.06}
% E4-perstyle.jsonl: max regret over the style set for v05
\providecommand{\vsixRegretFive}{}\renewcommand{\vsixRegretFive}{9.06}
% E4-perstyle.jsonl: max regret over the style set for v04
\providecommand{\vsixRegretFour}{}\renewcommand{\vsixRegretFour}{11.56}
% E5-ablations.json: reference pooled win rate
\providecommand{\vsixAblateRef}{}\renewcommand{\vsixAblateRef}{68.42}
% E8-ties.txt tie structure and every tie-break rule, v0.6 mirror
\providecommand{\vsixTieNSix}{}\renewcommand{\vsixTieNSix}{3,613}
% E8-ties.txt tie structure and every tie-break rule, v0.6 mirror
\providecommand{\vsixTieShareSix}{}\renewcommand{\vsixTieShareSix}{53.80}
% E8-ties.txt tie structure and every tie-break rule, v0.6 mirror
\providecommand{\vsixPathSixContested}{}\renewcommand{\vsixPathSixContested}{98.91}
% E8-ties.txt tie structure and every tie-break rule, v0.6 mirror
\providecommand{\vsixTieGapSix}{}\renewcommand{\vsixTieGapSix}{1.463}
% E8-ties.txt tie structure and every tie-break rule, v0.6 mirror
\providecommand{\vsixTieSpreadSix}{}\renewcommand{\vsixTieSpreadSix}{12.195}
% E8-ties.txt tie structure and every tie-break rule, v0.6 mirror
\providecommand{\vsixTieSameCardSix}{}\renewcommand{\vsixTieSameCardSix}{4.35}
% E8-ties.txt tie structure and every tie-break rule, v0.6 mirror
\providecommand{\vsixExactBuildUsSix}{}\renewcommand{\vsixExactBuildUsSix}{410.9}
% E8-ties.txt tie structure and every tie-break rule, v0.6 mirror
\providecommand{\vsixTieSameSuitSix}{}\renewcommand{\vsixTieSameSuitSix}{93.14}
% E8-ties.txt tie structure and every tie-break rule, v0.6 mirror
\providecommand{\vsixTieExactSepSix}{}\renewcommand{\vsixTieExactSepSix}{0}
% E8-ties.txt tie structure and every tie-break rule, v0.6 mirror
\providecommand{\vsixTieShipMovedSix}{}\renewcommand{\vsixTieShipMovedSix}{0.72}
% E8-ties.txt tie structure and every tie-break rule, v0.6 mirror
\providecommand{\vsixTieExactSepPctSix}{}\renewcommand{\vsixTieExactSepPctSix}{0.00}
% E8-ties.txt tie structure and every tie-break rule, v0.6 mirror
\providecommand{\vsixTieHitArraySix}{}\renewcommand{\vsixTieHitArraySix}{42.18}
% E8-ties.txt tie structure and every tie-break rule, v0.6 mirror
\providecommand{\vsixTieHitFastSix}{}\renewcommand{\vsixTieHitFastSix}{42.18}
% E8-ties.txt tie structure and every tie-break rule, v0.6 mirror
\providecommand{\vsixTieHitShipSix}{}\renewcommand{\vsixTieHitShipSix}{42.24}
% E8-ties.txt tie structure and every tie-break rule, v0.6 mirror
\providecommand{\vsixTieHitExactSix}{}\renewcommand{\vsixTieHitExactSix}{42.18}
% E8-ties.txt tie structure and every tie-break rule, v0.6 mirror
\providecommand{\vsixTieHitBestSix}{}\renewcommand{\vsixTieHitBestSix}{70.77}
% E8-ties.txt tie structure and every tie-break rule, v0.6 mirror
\providecommand{\vsixFullFastSix}{}\renewcommand{\vsixFullFastSix}{51.26}
% E8-ties.txt tie structure and every tie-break rule, v0.6 mirror
\providecommand{\vsixFullExactSix}{}\renewcommand{\vsixFullExactSix}{48.10}
% E8-ties.txt tie structure and every tie-break rule, v0.6 mirror
\providecommand{\vsixFullBestSix}{}\renewcommand{\vsixFullBestSix}{98.94}
% E8-ties.txt tie structure and every tie-break rule, v0.6 mirror
\providecommand{\vsixDisagreeSix}{}\renewcommand{\vsixDisagreeSix}{23.19}
% E8-ties.txt tie structure, v0.5 mirror (reference)
\providecommand{\vsixTieNFive}{}\renewcommand{\vsixTieNFive}{3,773}
% E8-ties.txt tie structure, v0.5 mirror (reference)
\providecommand{\vsixTieShareFive}{}\renewcommand{\vsixTieShareFive}{54.74}
% E8-ties.txt tie structure, v0.5 mirror (reference)
\providecommand{\vsixPathFiveContested}{}\renewcommand{\vsixPathFiveContested}{98.77}
% E8-ties.txt tie structure, v0.5 mirror (reference)
\providecommand{\vsixTieGapFive}{}\renewcommand{\vsixTieGapFive}{1.537}
% E8-ties.txt tie structure, v0.5 mirror (reference)
\providecommand{\vsixTieSpreadFive}{}\renewcommand{\vsixTieSpreadFive}{7.536}
% E8-ties.txt tie structure, v0.5 mirror (reference)
\providecommand{\vsixTieSameCardFive}{}\renewcommand{\vsixTieSameCardFive}{4.59}
% E8-ties.txt tie structure, v0.5 mirror (reference)
\providecommand{\vsixExactBuildUsFive}{}\renewcommand{\vsixExactBuildUsFive}{428.6}
% E8-ties.txt tie structure, v0.5 mirror (reference)
\providecommand{\vsixTieSameSuitFive}{}\renewcommand{\vsixTieSameSuitFive}{93.77}
% E8-ties.txt tie structure, v0.5 mirror (reference)
\providecommand{\vsixTieExactSepFive}{}\renewcommand{\vsixTieExactSepFive}{0}
% E8-ties.txt tie structure, v0.5 mirror (reference)
\providecommand{\vsixTieShipMovedFive}{}\renewcommand{\vsixTieShipMovedFive}{66.23}
% E8-ties.txt tie structure, v0.5 mirror (reference)
\providecommand{\vsixTieExactSepPctFive}{}\renewcommand{\vsixTieExactSepPctFive}{0.00}
% E8-ties.txt tie structure, v0.5 mirror (reference)
\providecommand{\vsixTieHitArrayFive}{}\renewcommand{\vsixTieHitArrayFive}{43.81}
% E8-ties.txt tie structure, v0.5 mirror (reference)
\providecommand{\vsixTieHitFastFive}{}\renewcommand{\vsixTieHitFastFive}{43.81}
% E8-ties.txt tie structure, v0.5 mirror (reference)
\providecommand{\vsixTieHitShipFive}{}\renewcommand{\vsixTieHitShipFive}{43.76}
% E8-ties.txt tie structure, v0.5 mirror (reference)
\providecommand{\vsixTieHitExactFive}{}\renewcommand{\vsixTieHitExactFive}{43.81}
% E8-ties.txt tie structure, v0.5 mirror (reference)
\providecommand{\vsixTieHitBestFive}{}\renewcommand{\vsixTieHitBestFive}{70.24}
% E8-ties.txt tie structure, v0.5 mirror (reference)
\providecommand{\vsixFullFastFive}{}\renewcommand{\vsixFullFastFive}{51.96}
% E8-ties.txt tie structure, v0.5 mirror (reference)
\providecommand{\vsixFullExactFive}{}\renewcommand{\vsixFullExactFive}{48.36}
% E8-ties.txt tie structure, v0.5 mirror (reference)
\providecommand{\vsixFullBestFive}{}\renewcommand{\vsixFullBestFive}{98.98}
% E8-ties.txt tie structure, v0.5 mirror (reference)
\providecommand{\vsixDisagreeFive}{}\renewcommand{\vsixDisagreeFive}{22.91}
% E8-ties.txt tie structure against a NON-mirror opponent
\providecommand{\vsixTieNCross}{}\renewcommand{\vsixTieNCross}{4,081}
% E8-ties.txt tie structure against a NON-mirror opponent
\providecommand{\vsixTieShareCross}{}\renewcommand{\vsixTieShareCross}{55.16}
% E8-ties.txt tie structure against a NON-mirror opponent
\providecommand{\vsixPathCrossContested}{}\renewcommand{\vsixPathCrossContested}{99.08}
% E8-ties.txt tie structure against a NON-mirror opponent
\providecommand{\vsixTieGapCross}{}\renewcommand{\vsixTieGapCross}{1.544}
% E8-ties.txt tie structure against a NON-mirror opponent
\providecommand{\vsixTieSpreadCross}{}\renewcommand{\vsixTieSpreadCross}{7.767}
% E8-ties.txt tie structure against a NON-mirror opponent
\providecommand{\vsixTieSameCardCross}{}\renewcommand{\vsixTieSameCardCross}{4.02}
% E8-ties.txt tie structure against a NON-mirror opponent
\providecommand{\vsixExactBuildUsCross}{}\renewcommand{\vsixExactBuildUsCross}{396.7}
% E8-ties.txt tie structure against a NON-mirror opponent
\providecommand{\vsixTieSameSuitCross}{}\renewcommand{\vsixTieSameSuitCross}{93.87}
% E8-ties.txt tie structure against a NON-mirror opponent
\providecommand{\vsixTieExactSepCross}{}\renewcommand{\vsixTieExactSepCross}{0}
% E8-ties.txt tie structure against a NON-mirror opponent
\providecommand{\vsixTieShipMovedCross}{}\renewcommand{\vsixTieShipMovedCross}{66.70}
% E8-ties.txt tie structure against a NON-mirror opponent
\providecommand{\vsixTieExactSepPctCross}{}\renewcommand{\vsixTieExactSepPctCross}{0.00}
% E8-ties.txt tie structure against a NON-mirror opponent
\providecommand{\vsixTieHitArrayCross}{}\renewcommand{\vsixTieHitArrayCross}{44.84}
% E8-ties.txt tie structure against a NON-mirror opponent
\providecommand{\vsixTieHitFastCross}{}\renewcommand{\vsixTieHitFastCross}{44.84}
% E8-ties.txt tie structure against a NON-mirror opponent
\providecommand{\vsixTieHitShipCross}{}\renewcommand{\vsixTieHitShipCross}{43.91}
% E8-ties.txt tie structure against a NON-mirror opponent
\providecommand{\vsixTieHitExactCross}{}\renewcommand{\vsixTieHitExactCross}{44.84}
% E8-ties.txt tie structure against a NON-mirror opponent
\providecommand{\vsixTieHitBestCross}{}\renewcommand{\vsixTieHitBestCross}{71.11}
% E8-ties.txt tie structure against a NON-mirror opponent
\providecommand{\vsixFullFastCross}{}\renewcommand{\vsixFullFastCross}{55.23}
% E8-ties.txt tie structure against a NON-mirror opponent
\providecommand{\vsixFullExactCross}{}\renewcommand{\vsixFullExactCross}{50.76}
% E8-ties.txt tie structure against a NON-mirror opponent
\providecommand{\vsixFullBestCross}{}\renewcommand{\vsixFullBestCross}{98.78}
% E8-ties.txt tie structure against a NON-mirror opponent
\providecommand{\vsixDisagreeCross}{}\renewcommand{\vsixDisagreeCross}{25.49}
% E8-belief.txt exact row
\providecommand{\vsixBeliefExactNLL}{}\renewcommand{\vsixBeliefExactNLL}{1.42246}
% E8-belief.txt exact row
\providecommand{\vsixBeliefExactHit}{}\renewcommand{\vsixBeliefExactHit}{47.94}
% E8-belief.txt theta=0 row
\providecommand{\vsixBeliefNoPriorNLL}{}\renewcommand{\vsixBeliefNoPriorNLL}{1.39083}
% E8-belief.txt theta=0 row
\providecommand{\vsixBeliefNoPriorHit}{}\renewcommand{\vsixBeliefNoPriorHit}{49.99}
% E8-belief.txt shipped-theta row
\providecommand{\vsixBeliefShippedNLL}{}\renewcommand{\vsixBeliefShippedNLL}{1.38218}
% E8-belief.txt shipped-theta row
\providecommand{\vsixBeliefShippedHit}{}\renewcommand{\vsixBeliefShippedHit}{51.49}
% E9-throughput.txt
\providecommand{\vsixThroughputFour}{}\renewcommand{\vsixThroughputFour}{147.9}
% E9-throughput.txt
\providecommand{\vsixThroughputFive}{}\renewcommand{\vsixThroughputFive}{276.3}
% E9-throughput.txt
\providecommand{\vsixThroughputSix}{}\renewcommand{\vsixThroughputSix}{303.4}
% E10-gateaudit.txt: one audit run
\providecommand{\vsixGateN}{}\renewcommand{\vsixGateN}{1}
% E10-gateaudit.txt: (opportunity, half-suit) pairs
\providecommand{\vsixGateStarved}{}\renewcommand{\vsixGateStarved}{31,321,233}
% E10-gateaudit.txt
\providecommand{\vsixGateOpportunities}{}\renewcommand{\vsixGateOpportunities}{5,472,906}
% E10-gateaudit.txt
\providecommand{\vsixGateRejected}{}\renewcommand{\vsixGateRejected}{14,449,770}
% E10-gateaudit.txt
\providecommand{\vsixGateFalseNeg}{}\renewcommand{\vsixGateFalseNeg}{0}
% E10-gateaudit.txt
\providecommand{\vsixGateVerdict}{}\renewcommand{\vsixGateVerdict}{PASS}
% E12-search.jsonl: blueprint-forced control
\providecommand{\vsixSearchControl}{}\renewcommand{\vsixSearchControl}{49.31}
% E12-search.jsonl: unguarded argmax
\providecommand{\vsixSearchUnguarded}{}\renewcommand{\vsixSearchUnguarded}{13.61}
% E12-search.jsonl: kappa=2.5
\providecommand{\vsixSearchGuarded}{}\renewcommand{\vsixSearchGuarded}{53.75}
% E12-search.jsonl: endgame search with the deliberate miss
\providecommand{\vsixSearchDead}{}\renewcommand{\vsixSearchDead}{50.28}
% E12-search.jsonl: range of the deviation-threshold ladder
\providecommand{\vsixSearchSpread}{}\renewcommand{\vsixSearchSpread}{40.28}
% E12-search.jsonl: blueprint-forced control minus the unguarded argmax
\providecommand{\vsixCurseGap}{}\renewcommand{\vsixCurseGap}{35.69}
% E14-searchdev.txt kappa=0
\providecommand{\vsixSearchDevZero}{}\renewcommand{\vsixSearchDevZero}{65.67}
% E14-searchdev.txt kappa=1
\providecommand{\vsixSearchDevOne}{}\renewcommand{\vsixSearchDevOne}{41.51}
% E14-searchdev.txt kappa=2.5
\providecommand{\vsixSearchDevTwoFive}{}\renewcommand{\vsixSearchDevTwoFive}{33.49}
% E14-searchdev.txt kappa=2.5 mean rollout depth
\providecommand{\vsixSearchDepth}{}\renewcommand{\vsixSearchDepth}{59.8}
% E14-searchdev.txt kappa=2.5 single-thread games/s
\providecommand{\vsixSearchGps}{}\renewcommand{\vsixSearchGps}{0.144}
% E14-searchdev.txt kappa=4
\providecommand{\vsixSearchDevFour}{}\renewcommand{\vsixSearchDevFour}{31.40}
% E14-searchdev.txt kappa=6
\providecommand{\vsixSearchDevSix}{}\renewcommand{\vsixSearchDevSix}{31.59}
% E15-deliberate-miss.txt unbanned-dead pathology
\providecommand{\vsixDeadPathEvents}{}\renewcommand{\vsixDeadPathEvents}{134.12}
% E15-deliberate-miss.txt unbanned-dead pathology
\providecommand{\vsixDeadPathPNinety}{}\renewcommand{\vsixDeadPathPNinety}{401}
% E15-deliberate-miss.txt unbanned-dead pathology
\providecommand{\vsixDeadPathPNineNine}{}\renewcommand{\vsixDeadPathPNineNine}{406}
% E15-deliberate-miss.txt unbanned-dead pathology
\providecommand{\vsixDeadPathHit}{}\renewcommand{\vsixDeadPathHit}{36.85}
% E15-deliberate-miss.txt unbanned-dead pathology
\providecommand{\vsixDeadPathDead}{}\renewcommand{\vsixDeadPathDead}{35.85}
% E15-deliberate-miss.txt unbanned-dead pathology
\providecommand{\vsixDeadPathLongest}{}\renewcommand{\vsixDeadPathLongest}{365}
% E15-deliberate-miss.txt unbanned-dead pathology
\providecommand{\vsixDeadPathRunSix}{}\renewcommand{\vsixDeadPathRunSix}{19.00}
% E15-deliberate-miss.txt unbanned-dead pathology
\providecommand{\vsixDeadPathRepeatACT}{}\renewcommand{\vsixDeadPathRepeatACT}{36.67}
% E15-deliberate-miss.txt unbanned-dead pathology
\providecommand{\vsixDeadPathOwnlocked}{}\renewcommand{\vsixDeadPathOwnlocked}{14.54}
% E15-deliberate-miss.txt unbanned-dead pathology
\providecommand{\vsixDeadPathDeclwrong}{}\renewcommand{\vsixDeadPathDeclwrong}{3.31}
% E15-deliberate-miss.txt unbanned-dead pathology
\providecommand{\vsixDeadPathLimit}{}\renewcommand{\vsixDeadPathLimit}{10.00}
% E15-deliberate-miss.txt unbanned-dead pathology
\providecommand{\vsixDeadPathForced}{}\renewcommand{\vsixDeadPathForced}{14}
% E15-deliberate-miss.txt unbanned-dead pathology
\providecommand{\vsixDeadPathForcedwrong}{}\renewcommand{\vsixDeadPathForcedwrong}{10}
% E15-deliberate-miss.txt win rate
\providecommand{\vsixDeadWin}{}\renewcommand{\vsixDeadWin}{51.1667}
% E15-deliberate-miss.txt CI low
\providecommand{\vsixDeadWinLo}{}\renewcommand{\vsixDeadWinLo}{48.3393}
% E15-deliberate-miss.txt CI high
\providecommand{\vsixDeadWinHi}{}\renewcommand{\vsixDeadWinHi}{53.9866}
% fitA.jsonl header
\providecommand{\vsixFitAPanel}{}\renewcommand{\vsixFitAPanel}{v05+v03+lockout+withholder}
% fitA.jsonl header
\providecommand{\vsixFitASeed}{}\renewcommand{\vsixFitASeed}{20260823}
% fitA.jsonl header
\providecommand{\vsixFitAPop}{}\renewcommand{\vsixFitAPop}{20}
% fitA.jsonl header
\providecommand{\vsixFitADeals}{}\renewcommand{\vsixFitADeals}{150x2}
% fitA.jsonl header
\providecommand{\vsixFitASigmaRel}{}\renewcommand{\vsixFitASigmaRel}{0.050}
% fitA.jsonl: generation records
\providecommand{\vsixFitAGens}{}\renewcommand{\vsixFitAGens}{30}
% fitA.jsonl: worst generation-0 clip fraction
\providecommand{\vsixFitAClipGenZero}{}\renewcommand{\vsixFitAClipGenZero}{0.00}
% fitC.jsonl header
\providecommand{\vsixFitCPanel}{}\renewcommand{\vsixFitCPanel}{v05+v03+withholder+feint}
% fitC.jsonl header
\providecommand{\vsixFitCSeed}{}\renewcommand{\vsixFitCSeed}{20260824}
% fitC.jsonl header
\providecommand{\vsixFitCPop}{}\renewcommand{\vsixFitCPop}{14}
% fitC.jsonl header
\providecommand{\vsixFitCDeals}{}\renewcommand{\vsixFitCDeals}{200x2}
% fitC.jsonl header
\providecommand{\vsixFitCSigmaRel}{}\renewcommand{\vsixFitCSigmaRel}{0.040}
% fitC.jsonl: generation records
\providecommand{\vsixFitCGens}{}\renewcommand{\vsixFitCGens}{14}
% fitC.jsonl: worst generation-0 clip fraction
\providecommand{\vsixFitCClipGenZero}{}\renewcommand{\vsixFitCClipGenZero}{50.00}
% engine/experiments_v06.sh: --rotations on every held-out cell
\providecommand{\vsixRotations}{}\renewcommand{\vsixRotations}{6}
% engine/src/arena.hpp pairedBootstrap/clusterBootstrap default B
\providecommand{\vsixBootReps}{}\renewcommand{\vsixBootReps}{20,000}
% engine/src/arena.hpp: 2.5/97.5 percentile bootstrap interval
\providecommand{\vsixBootAlpha}{}\renewcommand{\vsixBootAlpha}{95}
% engine/experiments_v06.sh E3: seed list
\providecommand{\vsixHeldOutBanks}{}\renewcommand{\vsixHeldOutBanks}{5}
% measurement host: 15 cores
\providecommand{\vsixCores}{}\renewcommand{\vsixCores}{15}
% a mirror cell at --rotations=2 duplicates exactly; see E2
\providecommand{\vsixMirrorDuplication}{}\renewcommand{\vsixMirrorDuplication}{2}
% engine/experiments_v06.sh E1
\providecommand{\vsixVerifyGames}{}\renewcommand{\vsixVerifyGames}{600}
% fitting seeds, disjoint from every evaluation bank
\providecommand{\vsixReservedSeeds}{}\renewcommand{\vsixReservedSeeds}{20260823, 20260824}
% engine/src/tuner.hpp self-test: a policy paired against itself
\providecommand{\vsixPairedSelfDelta}{}\renewcommand{\vsixPairedSelfDelta}{0.0000}
% engine/src/v06.hpp frozen vector, coordinate NFEAT+9
\providecommand{\vsixPriorTheta}{}\renewcommand{\vsixPriorTheta}{0.37062}
% engine/src/v06.hpp frozen vector, coordinate NFEAT+10
\providecommand{\vsixPriorPhi}{}\renewcommand{\vsixPriorPhi}{0.14525}
% engine/src/v06.hpp provenance stamp
\providecommand{\vsixFitProvenance}{}\renewcommand{\vsixFitProvenance}{fitC.jsonl [final weights record] sha256:be4a4c31d9153208 gens=14 obj=minimaxregret paired=1 panel=v05+v03+withholder+feint seed=20260824 deals=200x2 sigmaRel=0.04}
% research/v05/runs/fit-round1.jsonl: OLS slope of the v0.5 incumbent objective
\providecommand{\vsixFitSlope}{}\renewcommand{\vsixFitSlope}{+0.00049}
% research/v05/runs/fit-round1.jsonl: t statistic
\providecommand{\vsixFitSlopeT}{}\renewcommand{\vsixFitSlopeT}{1.60}
% research/v05/runs/fit-round1.jsonl: generations
\providecommand{\vsixFitBanks}{}\renewcommand{\vsixFitBanks}{40}
% research/v05/runs/fit-round1.jsonl: mean best-minus-incumbent
\providecommand{\vsixFitBestGap}{}\renewcommand{\vsixFitBestGap}{0.0404}
% expected maximum of 24 iid draws at the observed per-generation sd
\providecommand{\vsixFitOrderStat}{}\renewcommand{\vsixFitOrderStat}{0.0447}
% research/v05/runs/fit-round1.jsonl: population
\providecommand{\vsixFitPop}{}\renewcommand{\vsixFitPop}{24}
% E4 + F4 pooled: worst style cell
\providecommand{\vsixPoolWorstSix}{}\renewcommand{\vsixPoolWorstSix}{50.00}
% E4 + F4 pooled: argmin style
\providecommand{\vsixPoolWorstOppSix}{}\renewcommand{\vsixPoolWorstOppSix}{v04}
% E4 + F4 pooled: mean over the style set
\providecommand{\vsixPoolMeanSix}{}\renewcommand{\vsixPoolMeanSix}{78.63}
% E4 + F4 pooled: minimax regret over the style set
\providecommand{\vsixPoolRegretSix}{}\renewcommand{\vsixPoolRegretSix}{3.06}
% E4 + F4 pooled: worst style cell
\providecommand{\vsixPoolWorstFive}{}\renewcommand{\vsixPoolWorstFive}{50.00}
% E4 + F4 pooled: argmin style
\providecommand{\vsixPoolWorstOppFive}{}\renewcommand{\vsixPoolWorstOppFive}{v05}
% E4 + F4 pooled: mean over the style set
\providecommand{\vsixPoolMeanFive}{}\renewcommand{\vsixPoolMeanFive}{77.64}
% E4 + F4 pooled: minimax regret over the style set
\providecommand{\vsixPoolRegretFive}{}\renewcommand{\vsixPoolRegretFive}{8.60}
% E4 + F4 pooled: worst style cell
\providecommand{\vsixPoolWorstFour}{}\renewcommand{\vsixPoolWorstFour}{48.39}
% E4 + F4 pooled: argmin style
\providecommand{\vsixPoolWorstOppFour}{}\renewcommand{\vsixPoolWorstOppFour}{v05}
% E4 + F4 pooled: mean over the style set
\providecommand{\vsixPoolMeanFour}{}\renewcommand{\vsixPoolMeanFour}{76.77}
% E4 + F4 pooled: minimax regret over the style set
\providecommand{\vsixPoolRegretFour}{}\renewcommand{\vsixPoolRegretFour}{11.65}
% E4 + F4 pooled: withholder cell, v0.6
\providecommand{\vsixPoolWithholderSix}{}\renewcommand{\vsixPoolWithholderSix}{79.15}
% E4 + F4 pooled: withholder cell, v0.5
\providecommand{\vsixPoolWithholderFive}{}\renewcommand{\vsixPoolWithholderFive}{70.54}
% E4 + F4 pooled: withholder gain
\providecommand{\vsixPoolWithholderGain}{}\renewcommand{\vsixPoolWithholderGain}{8.60}
% E4 + F4 pooled: withholder games
\providecommand{\vsixPoolWithholderN}{}\renewcommand{\vsixPoolWithholderN}{4,800}
% E4 (515253) + F4 (90210, 424242)
\providecommand{\vsixPoolBanks}{}\renewcommand{\vsixPoolBanks}{3}
% E4 + F4 pooled: styles where v0.6 beats v0.5
\providecommand{\vsixPoolAhead}{}\renewcommand{\vsixPoolAhead}{7}
% E4 + F4 pooled: styles where it does not
\providecommand{\vsixPoolBehind}{}\renewcommand{\vsixPoolBehind}{5}
% E4 + F4 pooled: styles where the two are equal
\providecommand{\vsixPoolTied}{}\renewcommand{\vsixPoolTied}{1}
% E4 + F4 pooled: largest per-style loss against v0.5
\providecommand{\vsixPoolWorstLoss}{}\renewcommand{\vsixPoolWorstLoss}{1.72}
% E3 + E7-default + E5: every v0.6-vs-v0.5 cell in the battery
\providecommand{\vsixHeadPooled}{}\renewcommand{\vsixHeadPooled}{50.53}
% E3 + E7-default + E5: total games
\providecommand{\vsixHeadPooledN}{}\renewcommand{\vsixHeadPooledN}{11,300}
% number of head-to-head cells pooled
\providecommand{\vsixHeadCells}{}\renewcommand{\vsixHeadCells}{7}
% cells above parity
\providecommand{\vsixHeadCellsAbove}{}\renewcommand{\vsixHeadCellsAbove}{5}
% naive z on the pooled cells; ignores rotation clustering and therefore OVERSTATES significance
\providecommand{\vsixHeadZ}{}\renewcommand{\vsixHeadZ}{1.13}
% E7-rules.jsonl: dialects measured
\providecommand{\vsixDialectN}{}\renewcommand{\vsixDialectN}{4}
% E7-rules.jsonl: range of the win rate across dialects
\providecommand{\vsixDialectSpan}{}\renewcommand{\vsixDialectSpan}{0.47}
% E7-rules.jsonl: games per dialect
\providecommand{\vsixDialectDeals}{}\renewcommand{\vsixDialectDeals}{1,500}
% engine/experiments_v06.sh E7
\providecommand{\vsixDialectSeed}{}\renewcommand{\vsixDialectSeed}{828282}
% E8-ties.txt: mean exact count-law build, milliseconds (see vsixExactBuildUs*)
\providecommand{\vsixCountCost}{}\renewcommand{\vsixCountCost}{0.4}
% engine/src/belief.hpp DP_STATE_CAP and the capacity-vector bound
\providecommand{\vsixCountStates}{}\renewcommand{\vsixCountStates}{$10^5$}
% research/v06/notes/R10 section 1: us/event for the reduced blueprint, all six seats
\providecommand{\vsixSinkhornCost}{}\renewcommand{\vsixSinkhornCost}{9.07}
% research/v06/notes/R11 section 4: sd of one determinization s terminal set differential
\providecommand{\vsixLeafSd}{}\renewcommand{\vsixLeafSd}{2.5}
% research/v06/notes/R10 section 2: median event at which the consistent-deal count falls below 10^5
\providecommand{\vsixCollapseEvents}{}\renewcommand{\vsixCollapseEvents}{76}
% research/v05/results/P3-deception.md section 4 (transcribed)
\providecommand{\vsixPriorDeleteCost}{}\renewcommand{\vsixPriorDeleteCost}{4.60}
% research/v05/results/P3-deception.md section 4 (transcribed)
\providecommand{\vsixPriorDeleteCI}{}\renewcommand{\vsixPriorDeleteCI}{[2.63, 6.58]}
% F0-search-confirm.jsonl: replication banks at kappa=2.5
\providecommand{\vsixSearchReplBanks}{}\renewcommand{\vsixSearchReplBanks}{2}
% F0-search-confirm.jsonl: mean over the replication banks
\providecommand{\vsixSearchReplMean}{}\renewcommand{\vsixSearchReplMean}{52.08}
% F0-search-confirm.jsonl: games
\providecommand{\vsixSearchReplN}{}\renewcommand{\vsixSearchReplN}{1,440}
% F0-search-confirm.jsonl: worst replication bank
\providecommand{\vsixSearchReplMin}{}\renewcommand{\vsixSearchReplMin}{52.08}
% F0-search-confirm.jsonl: v0.6-Search against v0.6, pooled over every cell
\providecommand{\vsixSearchLift}{}\renewcommand{\vsixSearchLift}{52.08}
% F0-search-confirm.jsonl: games
\providecommand{\vsixSearchLiftN}{}\renewcommand{\vsixSearchLiftN}{2,880}
% F0-search-confirm.jsonl: cells
\providecommand{\vsixSearchLiftCells}{}\renewcommand{\vsixSearchLiftCells}{4}
% F0-search-confirm.jsonl: worst cell
\providecommand{\vsixSearchLiftMin}{}\renewcommand{\vsixSearchLiftMin}{50.83}
% F0-search-confirm.jsonl: cells above parity
\providecommand{\vsixSearchLiftAbove}{}\renewcommand{\vsixSearchLiftAbove}{4}
% E12 + F0: every bank of the kappa=2.5 search against v0.5
\providecommand{\vsixSearchPooled}{}\renewcommand{\vsixSearchPooled}{52.64}
% E12 + F0: games
\providecommand{\vsixSearchPooledN}{}\renewcommand{\vsixSearchPooledN}{2,160}
% E12 + F0: banks
\providecommand{\vsixSearchPooledBanks}{}\renewcommand{\vsixSearchPooledBanks}{3}
% E12 + F0: worst bank
\providecommand{\vsixSearchPooledMin}{}\renewcommand{\vsixSearchPooledMin}{52.08}
% F9-tieonly.txt: mean over the replication banks, vs v0.5
\providecommand{\vsixTieSearchMean}{}\renewcommand{\vsixTieSearchMean}{50.88}
% F9-tieonly.txt: per-bank
\providecommand{\vsixTieSearchBankList}{}\renewcommand{\vsixTieSearchBankList}{51.94, 50.69, 50.00}
% F9-tieonly.txt: banks
\providecommand{\vsixTieSearchBanks}{}\renewcommand{\vsixTieSearchBanks}{3}
% F9-tieonly.txt: worst bank
\providecommand{\vsixTieSearchMin}{}\renewcommand{\vsixTieSearchMin}{50.00}
% F9-tieonly.txt: tie-only search against v0.6 itself
\providecommand{\vsixTieSearchLift}{}\renewcommand{\vsixTieSearchLift}{50.00}
% F9-tieonly.txt
\providecommand{\vsixTieSearchLiftBanks}{}\renewcommand{\vsixTieSearchLiftBanks}{2}
% F8-tiesearch.txt A: search ONLY inside the tie group, 12 determinizations
\providecommand{\vsixSearchTieOnlyTwelve}{}\renewcommand{\vsixSearchTieOnlyTwelve}{54.1667}
% F8-tiesearch.txt A: search ONLY inside the tie group, 12 determinizations
\providecommand{\vsixSearchTieOnlyTwelveLo}{}\renewcommand{\vsixSearchTieOnlyTwelveLo}{50.5147}
% F8-tiesearch.txt A: search ONLY inside the tie group, 12 determinizations
\providecommand{\vsixSearchTieOnlyTwelveHi}{}\renewcommand{\vsixSearchTieOnlyTwelveHi}{57.7744}
% F8-tiesearch.txt B: the same, ONE determinization (near-random tie choice)
\providecommand{\vsixSearchTieOnlyOne}{}\renewcommand{\vsixSearchTieOnlyOne}{49.5833}
% F8-tiesearch.txt B: the same, ONE determinization (near-random tie choice)
\providecommand{\vsixSearchTieOnlyOneLo}{}\renewcommand{\vsixSearchTieOnlyOneLo}{45.9432}
% F8-tiesearch.txt B: the same, ONE determinization (near-random tie choice)
\providecommand{\vsixSearchTieOnlyOneHi}{}\renewcommand{\vsixSearchTieOnlyOneHi}{53.2279}
% F8-tiesearch.txt D: the FREE version -- uniform random choice inside the tie group, seeded from the public event stream, no rollouts
\providecommand{\vsixSearchTieRandom}{}\renewcommand{\vsixSearchTieRandom}{50}
% F8-tiesearch.txt D: the FREE version -- uniform random choice inside the tie group, seeded from the public event stream, no rollouts
\providecommand{\vsixSearchTieRandomLo}{}\renewcommand{\vsixSearchTieRandomLo}{47.1756}
% F8-tiesearch.txt D: the FREE version -- uniform random choice inside the tie group, seeded from the public event stream, no rollouts
\providecommand{\vsixSearchTieRandomHi}{}\renewcommand{\vsixSearchTieRandomHi}{52.8244}
% F7-winrate.txt: kappa applied inside the tie group too
\providecommand{\vsixSearchTieGuardedWin}{}\renewcommand{\vsixSearchTieGuardedWin}{49.8611}
% F7-winrate.txt
\providecommand{\vsixSearchTieGuardedLo}{}\renewcommand{\vsixSearchTieGuardedLo}{46.2194}
% F7-winrate.txt
\providecommand{\vsixSearchTieGuardedHi}{}\renewcommand{\vsixSearchTieGuardedHi}{53.5043}
% F7-tieguard.txt kappa=2.5, tie group ALSO guarded (kappatie=kappa)
\providecommand{\vsixSearchDevGuardedTwoFive}{}\renewcommand{\vsixSearchDevGuardedTwoFive}{2.70}
% F7-tieguard.txt kappa=4, tie group ALSO guarded
\providecommand{\vsixSearchDevGuardedFour}{}\renewcommand{\vsixSearchDevGuardedFour}{0.11}
% F6-reconstruction.txt seed 31
\providecommand{\vsixReconStates}{}\renewcommand{\vsixReconStates}{3,216}
% F6-reconstruction.txt seed 31
\providecommand{\vsixReconChecks}{}\renewcommand{\vsixReconChecks}{103,644}
% F6-reconstruction.txt seed 31: mask mismatch rate
\providecommand{\vsixReconMaskPct}{}\renewcommand{\vsixReconMaskPct}{0.3338}
% F6-reconstruction.txt seed 31: mismatches where the reconstruction is WIDER
\providecommand{\vsixReconWider}{}\renewcommand{\vsixReconWider}{346}
% F6-reconstruction.txt seed 31: mismatches where it is NARROWER
\providecommand{\vsixReconNarrower}{}\renewcommand{\vsixReconNarrower}{0}
% F6-reconstruction.txt: seed banks
\providecommand{\vsixReconBanks}{}\renewcommand{\vsixReconBanks}{2}
% E5 + F10: paired panel comparisons
\providecommand{\vsixPairedBanks}{}\renewcommand{\vsixPairedBanks}{3}
% E5 + F10: mean paired margin
\providecommand{\vsixPairedMean}{}\renewcommand{\vsixPairedMean}{2.16}
% E5 + F10: smallest paired margin
\providecommand{\vsixPairedMin}{}\renewcommand{\vsixPairedMin}{1.13}
% E5 + F10: banks whose CI excludes zero
\providecommand{\vsixPairedExcl}{}\renewcommand{\vsixPairedExcl}{2}
% F1-chain2x2.json: reference pooled win rate
\providecommand{\vsixFourwayRef}{}\renewcommand{\vsixFourwayRef}{68.42}
% X1-lbr.jsonl: fitted best response against v04
\providecommand{\vsixLbrFour}{}\renewcommand{\vsixLbrFour}{50.69}
% X1-lbr.jsonl
\providecommand{\vsixLbrFourLo}{}\renewcommand{\vsixLbrFourLo}{49.06}
% X1-lbr.jsonl
\providecommand{\vsixLbrFourHi}{}\renewcommand{\vsixLbrFourHi}{52.31}
% X1-lbr.jsonl
\providecommand{\vsixLbrFourN}{}\renewcommand{\vsixLbrFourN}{3,600}
% X1-lbr.jsonl: fitted best response against v05
\providecommand{\vsixLbrFive}{}\renewcommand{\vsixLbrFive}{50.31}
% X1-lbr.jsonl
\providecommand{\vsixLbrFiveLo}{}\renewcommand{\vsixLbrFiveLo}{48.75}
% X1-lbr.jsonl
\providecommand{\vsixLbrFiveHi}{}\renewcommand{\vsixLbrFiveHi}{51.89}
% X1-lbr.jsonl
\providecommand{\vsixLbrFiveN}{}\renewcommand{\vsixLbrFiveN}{3,600}
% X1-lbr.jsonl: fitted best response against v06
\providecommand{\vsixLbrSix}{}\renewcommand{\vsixLbrSix}{48.36}
% X1-lbr.jsonl
\providecommand{\vsixLbrSixLo}{}\renewcommand{\vsixLbrSixLo}{46.75}
% X1-lbr.jsonl
\providecommand{\vsixLbrSixHi}{}\renewcommand{\vsixLbrSixHi}{49.97}
% X1-lbr.jsonl
\providecommand{\vsixLbrSixN}{}\renewcommand{\vsixLbrSixN}{3,600}
% research/v06/results/: artifact count
\providecommand{\vsixManifestFiles}{}\renewcommand{\vsixManifestFiles}{35}
% build_tables_v06.py run date
\providecommand{\vsixBuildDate}{}\renewcommand{\vsixBuildDate}{2026-08-23}
% macros used in sections_v06 that this script generates
\providecommand{\vsixProvGenerated}{}\renewcommand{\vsixProvGenerated}{201}
% macros used in sections_v06 transcribed from earlier studies
\providecommand{\vsixProvTranscribed}{}\renewcommand{\vsixProvTranscribed}{84}
% research/v06/results: span of the battery artifact mtimes, hours
\providecommand{\vsixWallClock}{}\renewcommand{\vsixWallClock}{1.5}
%% --- end numbers_v06_generated ---

\begin{document}
\maketitle
\thispagestyle{fancy}

\begin{abstract}
%% --- begin sections_v06/abstract ---
FishBot v0.5 removed a failure mode and gained nothing in win rate, and said so. This report asks
where the remaining value in six-player Canadian Fish actually is. The answer turns out to sit in
one place, and it is a place three studies had written off.

At \vsixTieShareFive\% of \vfive{}'s ask decisions two or more candidates score \emph{bit-for-bit
identically} and the winner is whichever the enumerator emitted first. \vsixTieSameSuitFive\% of
those ties are two cards of one half-suit at one target, and the \emph{exact} posterior --- not
merely the approximation the policy deploys --- assigns them identical probability in
\vsixTieExactSepPctFive\% of cases. Every rule that ranks them by that probability therefore
realises the same hit rate to three significant figures, and so does enumeration order. The obvious
conclusion, which an earlier draft of this work drew, is that the ties are irreducible.

They are not. Resolving the tie group uniformly at random, seeded from the public event stream, is
worth \emph{exactly} nothing: \vsixSearchTieRandom\% [\vsixSearchTieRandomLo,
\vsixSearchTieRandomHi], mean half-suits $4.500$ to $4.500$. Resolving it with one sampled deal is
worth nothing either: \vsixSearchTieOnlyOne\%. Resolving it with an ensemble of twelve sampled
deals --- each a full determinization of the hidden state, played forward by six players
reconstructed at their own information sets rather than made clairvoyant --- is worth
\textbf{\vsixSearchPooled\%} against \vfive{} over \vsixSearchPooledN{} games and
\vsixSearchPooledBanks{} banks, none below \vsixSearchPooledMin\%. The separating information is
real and it lives in the \emph{joint} posterior, which every marginal integrates away.

Two negative controls make that a claim about \emph{information} rather than about a deterministic
policy being readable, and they are the reason the result is worth stating. Where the gain sits
within the search --- inside the tie group, outside it, or both --- is \emph{not} resolved at this
budget: the restricted variants land between \vsixSearchTieGuardedWin\% and
\vsixTieSearchMean\% and their intervals all overlap the unrestricted search's.

Getting there required correcting the statistic first. An unguarded argmax over the same rollout
ensemble is \textbf{\vsixCurseGap{} points worse} than the blueprint it searches from
(\vsixSearchUnguarded\% against a \vsixSearchControl\% control at identical budget, seed and sample
size), because one determinization's return is the final half-suit differential and genuine
differences between candidates are an order of magnitude smaller than its spread. A paired
lower-confidence-bound deviation rule recovers all of it.

Two further results. The exact observer-conditioned posterior is a \emph{worse} predictor of card
location than the approximation the policy deploys --- \vsixBeliefExactHit\% argmax hit rate at
\vsixBeliefExactNLL{} nats against \vsixBeliefShippedHit\% at \vsixBeliefShippedNLL{} --- so the
matched-budget refit under exact inference that three studies listed as outstanding is not worth
running. And \vfive{}'s forty-generation refit is indistinguishable from sampling noise, so its
shipped vector is v0.4's; five mechanisms in this study cleared a 95\% paired interval at the
per-cell budget this project has always used and returned exactly zero at three times it.

What ships is two configurations. \vsix{} is the deployed policy: the same mechanisms as \vfive{}
with a vector found by a repaired optimiser, \vsixThroughputSix{} games per second, better than
\vfive{} on the paired panel at all three banks measured and \emph{not} separated from it head to
head (\vsixHeadPooled\% over \vsixHeadPooledN{} games, pooling every cell rather than the five that
agreed). \vsixsearch{} adds the tie-resolving search and is the strongest configuration measured,
at a cost of three orders of magnitude in throughput.
%% --- end sections_v06/abstract ---
\end{abstract}

%% --- begin sections_v06/01-introduction ---
\section{Introduction}
\label{sec:intro}

Six-player Canadian Fish is a team game of pure public information flow. Fifty-four cards are dealt
nine each into two alternating teams of three; a player asks a named opponent for a named card,
must already hold another card of that card's six-card half-suit, must not hold the card itself, and
keeps the turn on a hit and loses it on a miss; a team scores a half-suit by naming, correctly and
in full, which of its members holds each of the six cards. Every ask, every hit, every miss and
every declaration is seen by all six players. That single property makes the entire hidden state the
\emph{initial deal}, and it makes the posterior over that deal exactly computable --- which is why
this project has been able to treat inference as a solved problem and spend three studies on policy.

The three prior FishBots are, in order: v0.3, a fitted linear policy in a TypeScript simulator;
v0.4, a rules-exact C++ engine with an exact reference posterior, a fast deployed approximation, a
twenty-feature ask score and a value-based declaration rule; and v0.5, which diagnosed and removed a
failure mode in v0.4 --- a two-question deadlock in mirror play that ended in misdeclaration --- and
reported, in terms, that removing it was worth almost nothing in win rate.

That report is the starting point here. A project whose last two releases produced no measurable
strength gain is either at the ceiling of its policy class or is measuring badly, and the two
possibilities call for opposite responses. This study answers the question by building the
mechanisms the accumulated design register ranked highest, measuring each at a resolution the
project had not previously used, and then measuring the measurement.

\paragraph{Contributions.}
\begin{enumerate}
\item \textbf{Exchangeable ties are irreducible} (\S\ref{sec:diag-irreducible}). More than half of
  v0.5's ask decisions end in a bit-for-bit tie; the tied candidates are two cards of one half-suit
  at one target; the \emph{exact} posterior separates \vsixTieExactSepPctFive\% of them; and every
  tie-break rule realises the same hit rate. The largest item on the inherited register is a
  property of the information set, not a defect.
\item \textbf{Exactness is not the same as informativeness} (\S\ref{sec:belief-exactworse}). The
  exact posterior is a worse predictor of card location than the deployed approximation, because
  the approximation carries a soft policy prior and the exact object is exact under a uniform prior
  over deals. This closes an experiment three studies have listed as outstanding.
\item \textbf{The optimizer's curse in determinized search} (\S\ref{sec:search-curse}). A search
  that determinizes the deal but reconstructs all six continuation players at their own information
  sets --- not a double-dummy rollout --- is \vsixCurseGap{} points worse than its own blueprint
  when the action is chosen by an unguarded argmax, and a paired deviation rule recovers all of it.
  Guarded, it is the one mechanism in this study that beats a static rule: \vsixSearchLift\% against
  \vsix{} itself over \vsixSearchLiftN{} games with every cell above parity.
\item \textbf{A fitting harness that moves the policy} (\S\ref{sec:fitting}). Four measured defects
  in the previous optimiser, a repair, and a recovery test the previous one provably fails. The
  repaired optimiser finds a vector that trades nothing head to head for a large, replicated gain
  against the deception archetype.
\item \textbf{An evidence standard} (\S\ref{sec:protocol-paired}). Five mechanisms in this study
  cleared a $95\%$ paired interval at the per-cell budget this project has used since v0.3 and
  returned exactly zero at three times it, on two disjoint banks.
\item \textbf{Engine corrections} (\S\ref{sec:engine-repairs}, \S\ref{sec:corrections}), including
  an exact-inference object that returned another observer's tables whenever two observers held one,
  and a declaration pre-gate audit that had been dead code since v0.5.
\end{enumerate}

\paragraph{What this paper does not claim.} \vsix{} is \emph{not} separated from \vfive{} head to
head. Pooling every head-to-head cell in the battery --- not only the five that agree --- it is
\vsixHeadPooled\% over \vsixHeadPooledN{} games with \vsixHeadCellsAbove{} of \vsixHeadCells{}
cells above parity and a naive $z$ of \vsixHeadZ{}, and that $z$ ignores the correlation between
the six rotations of a deal, so it overstates rather than understates. What is separated is the
paired comparison over the opponent panel and the profile's worst case. The substantive claims are
robustness across playstyles, correctness, and a substantially narrowed design space.
%% --- end sections_v06/01-introduction ---
%% --- begin sections_v06/02-rules ---
\section{The game and the rules dialect studied}
\label{sec:rules}

Canadian Fish, also called Literature, is played by six players in two teams of
three, seated so that the teams alternate. The 54-card deck is partitioned into
nine \emph{half-suits} of six cards, each player is dealt nine, and the team
that takes at least five half-suits wins. On their turn a player asks one
opponent for one named card; the ask is legal only if the named card lies in a
live half-suit, the asker holds at least one other card of that half-suit and
not the named card, and the target is an opponent who still holds cards. A hit
transfers the card and retains the turn, a miss passes the turn to the target,
and every ask, answer, transfer and card count is public. A half-suit is claimed
by a \emph{declaration}: an exact allocation naming, for each of the six cards,
which member of the declarer's own team holds it, awarded to the declaring team
if every assignment is correct and to the opposing team otherwise.

That is the whole game, and it is the same game the two preceding reports
studied. The dialect is unchanged from \vfour{} and \vfive{}, so this section
states the selection and points at the treatment rather than repeating it:
\S\ref{app:dialect} gives the clause table and the clauses this study's
mechanisms depend on, and the v0.4 report's Appendix A gives the dialect at the
length required to re-implement it \cite{fishbot04}.

\subsection{The dialect studied}
\label{sec:rules-dialect}

Published descriptions of the game document mutually incompatible variants
\cite{pagat}, so the dialect below is one explicit selection among them and not
a canonical rule set. It matches the strategy corpus this line of work draws on
\cite{develin} and the rules of record supplied by the project owner, and every
contested choice is an engine switch, so its effect is measurable rather than
assumed. The consequential choices are: a 54-card deck in nine half-suits, nine
cards per player; declaration legal at any moment, in or out of turn, by a seat
holding cards or not; a misdeclaration awarding the half-suit to the opposing
team; lowest seat index first among simultaneous declarers; an eight-level
willingness ladder in the forced endgame that arises when one whole team is
cardless; and a 400-ask action cap whose residue is adjudicated neutrally by
physical majority. No prearranged signalling conventions are available to any
policy, and team communication is limited to that willingness bit and to the
successor chosen by a cardless turn-holder. Table~\ref{tab:dialect} in
\S\ref{app:dialect} lists each choice against the switch that changes it.
\S\ref{sec:results-dialects} reports the primary result under
\vsixDialectN{} dialect perturbations; unless a result says otherwise, every
number in this report was produced under the settings above, and the action cap
was reached in \vsixLimitRate{}\% of games.

Two of those choices are places where the engine is \emph{more restrictive} than
the rules of record allow, and both are coordination channels a policy could in
principle use. The v0.5 study measured both, and this report inherits its
conclusions rather than re-opening them.

\paragraph{Turn transfer.} When the turn reaches a cardless seat whose teammates
both still hold cards, the rules of record permit the team to negotiate openly
over who receives it, sharing nothing but willingness. The engine instead has
the cardless seat choose unilaterally from its own belief, and solicits no
willingness bit. The channel is empty: a genuine multi-candidate transfer
decision arises \numTransferEvents{} times per game, and replacing the
unilateral rule with a \emph{ground-truth} chooser on one team, paired on common
deals, is worth $\numTransferDelta \pm \numTransferCI$ sets per game over
\numTransferGames{} paired games, with the two arms diverging at all in
\numTransferDiverged{}\% of games
(\filepath{research/v05/results/P8-coordination.md}, \S1.4). A legal mechanism
cannot beat the clairvoyant one, so that interval bounds the whole channel, and
no version of this policy has ever been built to exploit it.

\paragraph{Declaration arbitration.} The rules do not say who acts first when
two seats offer a declaration at the same instant, so the order is a modelling
choice. The engine scans by lowest seat index, deliberately declining to rank
candidates by stated confidence, because a confidence comparison moves a
statistic of one seat's hand to five others. The cost of that refusal is
\numArbCost{}~pp of win rate pooled over \numArbGames{} games, range
\numArbRange{}~pp across five seeds --- and a clairvoyant arbitrator, which
picks a proposer whose named allocation is actually correct, is worth
\numArbOracle{}~pp on the same design
(\filepath{research/v05/results/P6-verify-arbitration-cost.md}, \S\S3--4). The
information-safe rule therefore forgoes a third of a percentage point, and no
rule of any kind, safe or unsafe, informed or omniscient, recovers appreciably
more.

\subsection{What is unchanged from the v0.4 and v0.5 studies}
\label{sec:rules-unchanged}

Nothing in the dialect, the engine's rule implementation or the deal generator
changed between \vfive{} and \vsix{}. A reader who follows the citation to the
v0.4 report's rules treatment will find it accurate on every clause except two,
which no report has yet retracted and which \S\ref{sec:corr-dialect} corrects
here:

\begin{enumerate}
  \item \textbf{An undeclared half-suit does not score nothing.} The v0.4 report
        justifies the unconditional second stage of its forcing rule on the
        premise that an unclaimed half-suit is worthless. Under the dialect
        actually evaluated --- by both studies --- residue at the action cap is
        awarded by physical majority, and on a wrong post-horizon declaration
        the declaring team holds a mean \numResidueHeld{} of the six cards. The
        stage converts a probable award into a certain concession. The v0.4
        paper states the dialect correctly in its rules section and contradicts
        it two sections later.
  \item \textbf{Restarting the forced-endgame ladder sharpens nothing.} The v0.4
        report argues that restarting the willingness ladder after each
        declaration lets early, safer declarations resolve cards and so sharpen
        the allocations available for later ones. The argument is sound and
        inoperative: on every occasion a team went cardless with live half-suits
        --- \numLadderOccasions{} of them at seed \numLadderSeed{} --- exactly
        \numLadderLiveSets{} half-suit was live, so there is no ordering to
        exploit.
\end{enumerate}

Both are restated correctly in \S\ref{app:dialect-corrections} for a reader who
reaches the appendix first, and both bear on this study: the first is a premise
the declaration mechanism of \S\ref{sec:mech-declare} must not inherit, and the
second describes a structure that a more patient declaration policy would begin
to produce.
%% --- end sections_v06/02-rules ---
%% --- begin sections_v06/03-related ---
\section{Related work}
\label{sec:related}

This section covers the work the rest of the report builds on, and one class of
work it deliberately declines. The background the main argument does not depend
on is collected in \S\ref{app:related}, the evaluation-methodology literature is
taken up in \S\ref{sec:protocol}, and \S\ref{app:rel-declined} states, for each
declined method, why it is declined. Absence is a design decision in this study,
not an omission, and it is recorded as one.

Three of this report's results sit against specific literatures and are placed
here rather than left implicit. The optimiser's-curse result of
\S\ref{sec:search-curse} belongs to the determinization line, where an unguarded
argmax over sampled worlds is a known-in-principle failure that is rarely
quantified. The belief result of \S\ref{sec:belief-exactworse} --- a
policy-aware approximation beating the exact posterior --- belongs to the
trick-taking inference line, which is the only literature that has measured the
same trade-off. And the shared common-knowledge object of
\S\ref{sec:mech-knowledge} is the common-information reduction, restricted to
the one case in which it is cheap.

\subsection{Fish and Literature}
\label{sec:related-fish}

The published material on this game is human and informal, and neither of its
two substantial sources is quantitative or describes a program. McLeod's Pagat
entry \cite{pagat} is the canonical rules source and records several points on
which dialects differ; ours is specified in \S\ref{sec:rules} and listed in full
in \S\ref{app:dialect}. Develin's chapter \cite{develin} is the most detailed
strategy writing we located, describes the same dialect, and supplies two
observations this design uses: an ask teaches the opponents as much as it
teaches a teammate, and a player may not declare a half-suit they hold entirely.
Secondary strategy pages \cite{dorsa,strategy-site,wikiliterature} restate
blackballing, asking the asker, and an endgame delegation heuristic; they are
evidence about how the game is played and talked about, not sources of
quantitative claims. \S\ref{app:rel-informal} records where the two substantial
sources disagree.

Computational prior art remains close to absent. The v0.4 study's searches found
no academic work on this game, no public repository containing search, planning
or a working learner for it, and no coverage in RLCard \cite{rlcard}; the v0.6
literature refresh re-ran that question against 2025--2026 work and found
nothing that changes it. The one agent located, Somani's \texttt{literature}
\cite{somani,somani-server}, plays the four-player, 48-card variant, which does
not raise the six-player allocation problem, and publishes no strength numbers;
a closed-source Android application advertises bots with no stated methodology
\cite{playstore}. Every novelty statement in this report therefore takes the
form ``we found no prior published work that \dots'' and should be read against
that scope. One 2026 item is worth naming because it may become the external
calibration this line of work lacks: Valet packages twenty-one traditional
imperfect-information card games in a common description language with measured
branching factors and durations \cite{goadrich2026valet}, though whether any of
the twenty-one is a team game is not stated on its abstract page and we did not
fetch the full text, so nothing here is compared against it. The three prior
FishBot reports \cite{fishbot03,fishbot04,fishbot05} are the record this study
extends and, in \S\ref{sec:corrections}, corrects.

\subsection{Determinization, and why double-dummy analysis fails here}
\label{sec:related-determinization}

Determinization, or perfect-information Monte Carlo --- sample hidden cards
consistently with the public history, solve the resulting perfect-information
game, and play the move with the best average value --- is the standard recipe
for hidden-hand card games \cite{levy,ginsberg-gib}, and is unsound however many
worlds are drawn. Frank and Basin named the two mechanisms: \emph{strategy
fusion}, in which the searcher plays a different strategy in each sampled world
although a real player must commit to one strategy per information set, and
\emph{non-locality}, in which a node's value depends on parts of the tree
outside its own subtree \cite{frank-basin-ai,frank-basin-aaai}. Long, Sturtevant,
Buro and Furtak characterise when the method nevertheless performs well, through
leaf correlation, bias and a disambiguation factor \cite{long-pimc}; Information
Set Monte Carlo Tree Search replaces independent trees with a single tree over
information sets \cite{cowling-ismcts,goodman-ris}; the $\alpha\mu$ family backs
up Pareto fronts of world-vectors instead of scalar averages, repairing both
mechanisms, and reports Bridge 3NT declarer play at 60.2\% to 63.0\% with twenty
worlds \cite{cazenave2019alphamu}; and knowledge-based paranoia search reasons
about the worst case over deduced-consistent worlds rather than the average,
reporting over 1,000 points on the extended Seeger scale in Skat
\cite{edelkamp2021paranoia}. Its runtime is not stated on the abstract page, so
no compute cost is attributed to it anywhere in this report.

Fish breaks the method for a reason none of that literature addresses. Its
perfect-information game is not merely fused, it is degenerate: a clairvoyant
player names only cards the target actually holds, so never fails an ask, never
loses the turn, and declares each half-suit correctly as soon as their team
holds all of it. Almost every reachable position carries the same
perfect-information value, and averaging that value over deals sampled from the
public history collapses to choosing the ask most likely to hit --- a one-ply
heuristic blind to the information an ask leaks, the information a refusal
supplies, and the option value of postponing a declaration.
$\alpha\mu$ repairs unsoundness, not degeneracy, so it does not rescue the
construction here. \S\ref{app:rel-determinization} gives the collapse in full.

What \S\ref{sec:search} adds is that the binding problem in this game is not the
abstract failure but the statistic. One determinization's value is the final
half-suit differential, with standard deviation \vsixLeafSd{} --- an order of
magnitude larger than the differences between the leading candidates --- so an
argmax over $D$ sampled means selects on residual noise, and does so more
confidently the more candidates it is offered. That is the optimiser's curse
rather than strategy fusion, and it is measured here at \vsixCurseGap{} points
by a control that forces the same code to return the blueprint's own pick
(\S\ref{sec:search-curse}). This report therefore uses determinization only as a
research probe, behind the paired deviation guard of \S\ref{sec:search-falsifier},
and reports the result including the configurations that land at zero
(\S\ref{sec:search-deadask}). The transferable claim is the guard, not the
search.

\subsection{Search in imperfect-information and cooperative games}
\label{sec:related-search}

States in an imperfect-information game do not have well-defined values, so
depth-limited search needs something at the leaf that represents the opponent's
ability to choose. Brown, Sandholm and Amos supply it: let the opponent select
among $k$ continuation strategies at the depth limit, each yielding its own
leaf-value vector, and require the searcher to be good against that choice ---
demonstrated at master level in heads-up no-limit hold'em on a four-core CPU and
16\,GB \cite{brown2018depthlimited}. That construction is the one piece of this
literature this study uses as method rather than context: it is the multi-valued
leaf of \S\ref{sec:mech-leaves} and the prerequisite for
\S\ref{sec:search}. The soundness argument is stated for two-player zero-sum
games, and Fish's team-level game presents three uncorrelated opponent seats at
the leaf, so nothing is inherited. It is adopted here as a robustification
heuristic with a measurable consequence, and audited rather than assumed
(\S\ref{sec:mech-exploit}).

Two lines address the other half of the problem, which is what a search may
prune. Zhang and Sandholm solve subgames over low-order knowledge rather than
the full common-knowledge closure, which is the published answer to a closure
too large to enumerate --- and they prove that the naive form can \emph{increase}
exploitability \cite{zhang2021subgamesolvingwithoutck}. That warning applies
directly to the live-ask gate this policy already ships, which is a
knowledge-limited pruning rule that has never been audited, and it is why
\S\ref{sec:mech-exploit} was built before any search work rather than after it.
Subgame solving in adversarial team games builds a gadget over the team belief
DAG and runs in polynomial time given a best-response oracle \emph{when the
blueprint is sparse} \cite{zhang2022teamsubgame}; a Fish blueprint is not
sparse, so that escape hatch does not open here, and \S\ref{app:rel-declined}
records this as a verified non-transfer rather than an unexamined gap.

In cooperative games the constraint is different and it transfers verbatim.
SPARTA improves a blueprint purely at test time, 24.08 to 24.61 in Hanabi, and
distinguishes single-agent search --- which assumes every other agent strictly
follows the blueprint --- from multi-agent search, in which all agents run the
\emph{same} common-knowledge procedure \cite{lerer2020sparta}. The moment a
second teammate also searches, the first teammate's belief over its partner is
computed under a policy that partner is no longer playing, and the improvement
guarantee is void. This is the Dec-POMDP common-knowledge caveat, and it is the
constraint most easily violated in silence. Fish satisfies it cheaply: all card
movement is public, so any \emph{deterministic} procedure that reads only public
state and its own hand is automatically common knowledge among the three
teammates, with no shared seed and no communication. The qualifier is
load-bearing --- a randomised tie-break loses the property unless its
randomisation is seeded from the public history, which is a one-integer decision
with a real correctness consequence (\S\ref{sec:mech-tiebreak}).

The remainder of this literature is cited as context. The full-solver line
\cite{moravcik2017deepstack,brown2020rebel,schmid2023sog} established that
test-time re-solving over public belief states is principled rather than a hack,
and ReBeL's safety result is the licence under which any test-time search here
operates; the algorithms are declined, because each requires iterating or
learning over a closure whose cardinality at deal time is
$45!/(9!)^5 \approx 1.9 \times 10^{28}$ deals from one seat. Update equivalence
recasts search as replicating a last-iterate-convergent update locally, needing
no public-information enumeration at all, and reports matching or exceeding
public-belief-state search in Hanabi at two orders of magnitude less search time
\cite{sokota2024updateequivalence}; its improvement guarantee is for
common-payoff games, which Fish is not, so it informs the shape of the mechanism
set without licensing it. Learned Belief Search buys 55--91\% of exact search's
benefit at 4.6--42$\times$ less compute \cite{hu2021lbs} and is simply redundant
here, because this game's exact belief is already computed and oracle-validated
(\S\ref{sec:belief-count}). Online planning by policy fine-tuning
\cite{fickinger2021rlfinetuning} requires gradient updates at decision time and
does not transfer to a CPU-only engine with no network. The depth-limited
counter-strategy line \cite{milec2021cdr,milec2025abd} carries one warning this
study honours --- robust-response methods do not compose with limited-look-ahead
solving off the shelf --- and the learned-model and gadget-refinement work
\cite{kubicek2025lamir,kubicek2026refinements} is recorded for its meta-lesson
rather than its machinery: when a search subproblem has many near-equivalent
optima, which one is selected is not a tie-break, and refinement is reported at
more than 50\% of exploitability. That is the same observation
\S\ref{sec:diag-ties} makes about an argmax resolved by array order.

Finally, every method above is published with a caveat that it can increase
exploitability, so none of them can be evaluated by win rate alone. Local best
response gives a cheap lower bound: a small greedy search, at most two actions
deep, showed every entrant in a 2016 poker competition to be more than
3,180\,mBB/h exploitable \cite{lisy2017lbr}. Its two documented pitfalls are
carried into \S\ref{sec:results-lbr} --- restricting where the response may
begin deviating understates exploitability, and a negative score proves nothing
--- and the figure this report publishes is a lower bound within its response
class, not an exploitability number.

\subsection{Policy-aware inference}
\label{sec:related-inference}

Skat engines enumerate the deals consistent with the public record and reweight
them by a model of what each player would have done --- a supervised per-card
location model \cite{solinas-supervised}, or reach probabilities taken from a
policy \cite{rebstock-inference} --- and a dedicated inference module was
reported worth $+217$ tournament points per 36 games to the engine Kermit on its
own \cite{buro-kermit}. The policy prior of \S\ref{sec:belief-prior} is the same
idea in a much simpler form: a fitted soft reweighting of the rules posterior,
with two parameters and no learned opponent model.

This literature is where the central result of \S\ref{sec:belief} belongs, and
it is the only literature that has measured the same trade-off. What is new here
is the margin and its direction. In this game the exact conditional posterior is
computable in closed form and cheap (\S\ref{sec:belief-count}), so the usual
excuse for approximating --- that exactness is unaffordable --- does not apply;
and yet substituting the exact posterior into an otherwise identical policy
\emph{loses}, at two independent seed banks. The exact path resolves more
distinctions and predicts hits worse, and the reason is structural rather than
numerical: the block construction consumes the deduction state alone and has no
parameter through which behaviour could enter, while the approximate path is
policy-aware by construction. Exactness and behaviour-awareness are not
available in the same object here, and \S\ref{sec:mech-tiebreak} is built on
that split. The report is careful about what this does and does not show: the
weights were fitted for the approximate belief, so it is a substitution result
rather than a verdict on exact inference (\S\ref{sec:limitations}).

Two adjacent lines are named and not used. Belief modelling in cooperative games
--- the Bayesian Action Decoder's public-belief formulation
\cite{foerster2019bad}, its factorised-prescription successor
\cite{sokota2021capi}, and the learned belief models above --- is aimed at games
whose exact belief is intractable, which is the opposite of the situation here.
And the exact-counting machinery this study does not need to approximate ---
Ryser's inclusion--exclusion formula \cite{ryser}, Glynn's sign-vector formula
\cite{glynn}, the Jerrum--Sinclair--Vigoda chain \cite{jsv}, matrix scaling
\cite{lsw} and sequential importance sampling \cite{chen-sis} --- is recorded in
\S\ref{app:rel-counting}, along with why the block structure of this game
sidesteps all of it. History filtering is FNP-complete in general
\cite{solinas-history}; Fish escapes because its public record determines every
card movement, so the only object to reconstruct is the initial deal. No learned
belief model appears anywhere in this study, and that is a stated non-goal
rather than an oversight.

\subsection{Team games and common information}
\label{sec:related-team}

Three teammates share a payoff, hold different information and cannot
communicate. The formal statement of that situation is the common-information
approach: reformulate the decentralised problem from the point of view of a
coordinator who knows the common information and chooses \emph{prescriptions},
maps from each controller's private information to its action
\cite{nayyar2013common}. In Fish the common information is the entire public
history, so the coordinator's information state is exactly the belief object the
engine already maintains --- which is the formal content of the mechanism in
\S\ref{sec:mech-knowledge} that replaces six per-seat deduction states with one
shared object plus per-seat refinement. It is also the reason the SPARTA
constraint above exists: if one seat searches and its teammates do not, they are
no longer executing the prescription the first seat's belief update assumes, and
the reduction silently fails. Solving the coordinator's problem is out of reach
--- a prescription maps every possible nine-card hand to an ask --- so this is a
correctness lens, not an algorithm.

The team-game equilibrium literature is cited to delimit scope. The target
concept for a Fish team is a team-maxmin equilibrium with correlation
\cite{celli-gatti}, connected to extensive-form correlation by Farina and
coauthors \cite{farina-collusion} and made computable through the team belief
DAG and the two-player conversion \cite{zhang-tbdag,carminati-tpi}, at a
representation cost exponential in how much private information distinguishes
teammates inside a public state. Three unrestricted nine-card hands is exactly
the case that cost forbids. Hidden-role machinery \cite{carminati2024hiddenrole}
does not apply at all, because Fish teams are public from the deal. Nothing in
this report is an equilibrium result of any kind.

Conventions are the other half of team play, and the coordination literature
supplies both the vocabulary and the warning. Hanabi is the reference benchmark
for coordination through observable actions \cite{bard-hanabi}; Other-Play
demonstrated how much of a self-play score is convention rather than skill
\cite{hu-otherplay}, the Simplified Action Decoder showed what conventions need
in order to survive exploratory training \cite{hu-sad}, and Off-Belief Learning
supplies an explicit dial on convention depth together with grounded play at
level one \cite{hu2021obl}. Grounded play is unusually strong in this game
because ask legality is itself a hard certificate: an ask proves its asker holds
another card of the named half-suit, so reading a teammate's ask grounded is
exact rather than approximate. Self-explaining deviations formalise the opposite
move --- deviate so that a teammate's theory of mind concludes circumstances are
abnormal --- and produced the first algorithmic finesse plays in Hanabi
\cite{hu2022sed}; in Fish the same channel is supplied by the rules rather than
learned, which is a structural advantage and, in cross-play, a structural risk.
Joint policy search is the bridge-bidding analogue, scoring simultaneous changes
at several information sets, at $+0.63$ IMPs per board against a strong
commercial program \cite{tian2020jps}; the Hanabi construction that most clearly
fails to port is the hat-guessing code, which works because each receiver sees
every other receiver's hand \cite{cox-hanabi}, and in Fish no player sees any
hand.

Human-regularised search is where that risk is priced. piKL anchors a search
result to a reference policy with a bounded divergence \cite{jacob2022pikl}, and
the same construction was extended to coordination with real human partners
\cite{hu2022humanregularized} and used as the planning core of a Diplomacy agent
whose game family --- adversarial teams acting through public actions --- is the
one Fish belongs to \cite{fair2022cicero}. The evidence that a small human
corpus is not the way in is also published: in the ad-hoc human-AI coordination
challenge, an agent trained with no human data at all outperformed a
best-response-to-behavioural-cloning baseline \cite{dizdarevic2025ah2ac2}. This
study uses none of these as method. It takes from them the requirement that the
partner regime be a stated parameter of the policy rather than an assumption,
which is \S\ref{sec:mech-partner}, and that a convention-bearing mechanism be
reported with its cross-play cell and not only its self-play cell.

The last entry is a counterweight to the rest of this section. A 2026
gold-standard study of lightweight game-playing agents concludes that the
binding constraint on strength is information rather than model capacity, and
that disciplined evaluation against a held-out expert beats deeper search
\cite{kelidari2026goldstandard}. That is this study's own thesis stated by
someone else, and it is the reason two of this report's four principal results
are negative and are reported as such.
%% --- end sections_v06/03-related ---
%% --- begin sections_v06/04-engine ---
\section{The simulator: implementation, information safety, verification, and the v0.6 repairs}
\label{sec:engine}

The simulator is the instrument every number in this report is read from, so its
properties are stated before its results. Three of them matter: it cannot leak a
hidden field to a policy, it checks its own deductions against the truth while it
runs, and a configuration in which every v0.6 mechanism is switched off is
\vfive{} bit for bit. The last of those is not a courtesy. It is what makes the
ablation table of \S\ref{sec:results-ablations} exact rather than approximate,
and it is the first thing \filepath{engine/experiments_v06.sh} runs.

This section also records four repairs to the engine itself. They are reported
here rather than folded silently into a diff because three of them are the
reason a measurement in this report was possible at all, and because one of them
--- a dead command-line option --- means a published pass in the v0.4 record was
vacuous (\S\ref{sec:corr-vfour}).

\subsection{Implementation}
\label{sec:engine-impl}

The engine is unchanged in structure from the v0.4 report, which describes it in
full \cite{fishbot04}: a single-header C++20 core in \filepath{engine/src/}, a
hand as one \verb|uint64_t| over the 54 card indices $c = 6h + i$, and
counter-based random streams keyed by \emph{(seed, seat, purpose)} so that a deal
and its entire play reproduce from the seed alone. The duplicate-block protocol
of \S\ref{sec:protocol} depends on that last property and on nothing else.

What moved in v0.6 is the policy layer. \texttt{V06Agent} derives from
\texttt{V05Agent} and defers to it whenever no v0.6 mechanism is enabled, so the
v0.5 decision path is not reimplemented but inherited; the determinized search of
\S\ref{sec:search} and the rollout blueprint it calls live in two new headers,
\filepath{engine/src/v06.hpp} and \filepath{engine/src/v06_rollout.hpp}. The
diagnostics this report leans on are a new subcommand rather than scratch
probes: \texttt{fish v6probe --mode=ties} censuses ask decisions and every
tie-break rule including hindsight, \texttt{--mode=belief} scores several
posteriors as predictors on identical states, and \texttt{--mode=search} reports
how often a guarded search actually deviates from its blueprint. All three are in
\filepath{engine/src/probe_v06.hpp} and all three are run by the battery, so the
figures they produce are regenerated rather than transcribed.

Throughput on the shipped configuration is \vsixThroughputSix{} games/s against
\vsixThroughputFive{} for \vfive{} and \vsixThroughputFour{} for \vfour{}, all on
the same machine at \vsixCores{} threads. Every compute figure in this report is
single-machine and was measured under load, so all of them are lower bounds.

\subsection{Information safety}
\label{sec:engine-safety}

The acting policy never receives the game state. It receives its own
\texttt{Knowledge} object --- the constraint bookkeeping of \S\ref{sec:belief}
--- and the public record, in which each event carries only the actor, the
target, the named card, the outcome, any declared allocation, and the resulting
public hand counts. There is no interface through which a policy can read
another player's cards, so the class of leak the \vthree{} audit reported in its
predecessor \cite{fishbot03} is excluded structurally rather than by inspection.
The property and the argument for it are unchanged from the v0.4 report
\cite{fishbot04}.

Soundness is the stronger claim, and it is checked rather than argued. After
every public event the engine's audit mode verifies, for every player and against
the true deal, that every card the constraint state records as resolved is held
by the seat it is resolved to; that the true holder of every unresolved card lies
in that card's surviving support; that the derived capacities equal the true
counts of unresolved cards per player; and that every live ask-legality
certificate is satisfied by the true deal. A single violation would mean the
belief engine had excluded the truth.

\subsection{Verification suite}
\label{sec:engine-verify}

\texttt{fish verify} plays the ordered cross-product of a \numVerifyPool{}-policy
pool --- \numVerifyPairs{} ordered pairs, so both orderings of every pair and
every self-play pairing are exercised --- with the audit above enabled
throughout, and additionally checks deck integrity and card conservation, that
exactly nine half-suits are resolved in every completed game, that a declaration
only ever assigns cards to the declarer's own teammates, that ask legality is
enforced independently on the move generator and on the applier, and that
replaying a seed at a fixed thread count reproduces the game bit-identically. It
returns a non-zero exit status on failure, which is what lets the battery use it
as a gate rather than as a report.

Over \vsixVerifyGames{} deals the suite recorded \vsixAuditViolations{}
violations in \vsixAuditChecks{} checks, with an overall verdict of
\vsixVerify{}. The random streams are keyed so that thread count cannot affect a
game's outcome, but the suite tests replay identity only at a fixed thread count
and does not test that claim directly; this is carried forward from the v0.4
study as an unclosed gap and is listed as one in \S\ref{app:repro-gaps}.

\subsection{The v0.6 engine repairs}
\label{sec:engine-repairs}

Each repair below is stated as a defect, the measurement that shows the defect is
real, what the defect blocked, and the measurement that shows it is fixed. Three
of the four are correctness-neutral for \vfive{}: they change nothing a \vfive{}
game does, which is why they survived a release. The fourth changes the Fast
posterior and is therefore a flag that defaults off.

\paragraph{BlockDP tables aliased across instances.}
\filepath{engine/src/blockdp.hpp} parked every instance's Unit, Group, forward
and backward tables in one \verb|thread_local| buffer pool, and
\texttt{BlockDP::build} took its storage from that pool without owning it. Two
\texttt{BlockDP} objects alive on one thread therefore shared storage, and a
second instance's \texttt{build} silently repointed the first instance's tables
at its own data. The driver makes this routine rather than exotic: agents own one
\texttt{block} member each and rebuild only when dirty, while
\texttt{Game::forcedEndgame} and \texttt{Game::declarationRound} poll several
agents back to back between public events.

The defect is real and was measured directly. Reading one instance's own group
table twice with an unrelated \texttt{build} in between gave
\numAliasMismatch{} mismatches in \numAliasChecks{} checks: one observer's
posterior tables had been replaced by another observer's
(\filepath{research/v05/results/P2-forced-endgame.md}, \S6). An earlier form of
the same check that issued a real query after the clobber died with
\texttt{SIGBUS}, because the stale forward and backward pointers index a table
sized for a different capacity vector.

What it blocked is the whole of \S\ref{sec:belief}. The deployed approximate
belief never builds a \texttt{BlockDP}, so nothing in the v0.4 or v0.5 release
path could observe the defect; anything exact hits it on the first
multi-agent poll. That asymmetry is exactly how a fatal defect survives into a
release, and it is why the exact posterior could not be evaluated as a predictor
before this study.

The repair is a per-thread generation stamp plus a stored copy of the source
\texttt{Knowledge}. Every \texttt{build} increments the counter and records its
value; every public query calls \texttt{ensureCurrent}, which re-derives the
instance lazily from its stored \texttt{Knowledge} when it finds its own stamp
stale. The check cannot be forgotten at a call site because it lives inside the
queries, and it costs $O(1)$ memory. The probe in
\filepath{engine/src/probe_forcedendgame.hpp} was extended to separate the two
things the original check conflated: \texttt{QUERY} mismatches, where the same
query issued on the same instance before and after an unrelated build returns
different answers, and raw shared-pool field reads, where the instance's own
pointers still address the pool and the fields behind them have changed. The
first must be zero and is \vsixAliasQueryMismatch{}. The second is
\vsixAliasRawMismatch{} and is \emph{not} fixed: the pool is still shared, the
raw fields still change, and the guard's contract is that no public query ever
reads them stale. Reporting only the first would overstate the repair.

\paragraph{\texttt{gateaudit} was dead code for v0.5.}
The declaration pre-gate audit is instrumentation that counts how many
declaration opportunities a cheap pre-gate rejects and how many of those
rejections were false negatives --- candidates the expensive test would have
accepted. The switch that enables it was parsed only in the \vfour{} branch of
\filepath{engine/src/factory.hpp}. A v0.5 specification therefore parsed
\texttt{gateaudit=1}, ignored it, and reported a pass over zero opportunities:
\texttt{fish gateaudit --a=v05:gateaudit=1} could not fail, because it never
looked at anything. \S\ref{sec:corr-vfour} records what that means for the v0.4
record's corresponding claim.

The option is now parsed in the v0.5 and v0.6 branches as well. The audit over
\vsixGateOpportunities{} declaration opportunities rejected \vsixGateRejected{}
at the cheap gate with \vsixGateFalseNeg{} false negatives, verdict
\vsixGateVerdict{}. The instrument is only as good as its opportunity count, so
that count is reported first and the verdict second.

\paragraph{Two defects in the tuner.}
\filepath{engine/src/tuner.hpp} chose between the two parameter-passing spellings
with \verb|w.size() > 18| where the ask-feature count \verb|NFEAT| ---
\numAskFeats{} --- belongs. The expression is correct today by coincidence, and
it is the same class of defect, in the same place, that cost the v0.4 study a
fitting round when the ask-feature count grew. It now reads \verb|NFEAT|. The
second defect is worse because it is silent and quantitative: the win rate was
computed as \verb|st.winsA / (st.games * 2)|, hard-coding the tuner's two
rotations into the denominator. Fitting at six rotations --- which the
duplicate-block design of \S\ref{sec:protocol} makes the natural choice --- would
have reported a third of the true win rate for every candidate, with no error and
no warning. It now reads \verb|st.games * mc.rotations|.

\paragraph{Certificate de-duplication, opt-in.}
\texttt{Knowledge::onEvent} in \filepath{engine/src/belief.hpp} appended a fresh
ask-legality certificate on every repeated ask without checking whether an
identical one, or one strictly stronger, was already stored. The store therefore
grows carrying no information, and every consumer that iterates it pays for the
duplicates. The obvious repair --- skip a certificate that duplicates a stored
one, or whose card set is a superset of a stored one for the same player --- is
implemented, and it is a flag that defaults off. The reason is that the
duplicates are not inert. The approximate conditioning step of
\S\ref{sec:belief} applies a certificate once per stored copy, so removing
copies \emph{changes} the Fast posterior and therefore changes play. Making
de-duplication the default would have silently altered every \vfour{} and
\vfive{} number this report compares against. With the flag off, both stay bit
for bit what they were.

\paragraph{The identity control, and why it is a result.}
\texttt{v06:legacy=1} restores \vfive{}'s parameter vector and zeroes the three
v0.6 ask terms, and with every mechanism switch off \texttt{V06Agent} defers to
its base class. The battery checks that this is an identity and not a close
agreement, by comparing the md5 digest of the entire \texttt{fish pathology}
output --- every pathology counter, not a win rate --- between
\texttt{v06:legacy=1} and \texttt{v05} on the same deals and seed. The verdict is
\vsixIdentity{}.

This matters more than a build-hygiene check. Every ablation in
\S\ref{sec:results-ablations} and \S\ref{app:ablations} is a paired comparison
between the v0.6 binary carrying mechanism $X$ and the v0.6 binary carrying
nothing. If the all-switches-off configuration merely resembled \vfive{}, each
row would confound the mechanism with whatever else differed, and the table would
be an approximation whose error had no bound. Because the configuration is
\vfive{} bit for bit, the difference each row reports is the mechanism and
nothing else. The control is cheap, it is run before any strength number is
collected, and it is the reason the ablation table can be read as exact.
%% --- end sections_v06/04-engine ---
%% --- begin sections_v06/05-belief ---
\section{The observer-conditioned deal posterior}
\label{sec:belief}

This section builds the object the rest of the report measures, and then reports
the first of the study's two principal negative results: the exact posterior is a
worse predictor than the approximation the engine actually deploys.

The order matters, because the result is easy to misread in either direction.
Two facts about the exact posterior are both true. It resolves distinctions the
approximation cannot --- it is the exact conditional count under the rules, with
no independence assumption anywhere in it --- and it predicts where cards are
worse than the approximation does. Resolution and calibration are different
properties, and this section keeps them apart throughout. The reason the exact
object is the worse predictor is not that the approximation is secretly better at
counting. It is that the exact object is exact \emph{with respect to a uniform
prior over deals}, and that prior is wrong.

\subsection{The hidden state is the initial deal}
\label{sec:belief-state}

In the dialect of \S\ref{sec:rules} a card changes hands in exactly one way, a
successful ask, which every player observes; declarations remove cards from play,
and that too is public. The v0.4 report proves the consequence
\cite{fishbot04}: a card whose location has never been publicly revealed is still
held by the player to whom it was dealt, the joint hidden state at any time is
exactly the initial deal, and the current holder of every card is a
deterministic function of the initial deal and the public history. The argument
is one line --- every change of holder is an announced event, so a card with no
such event has not moved --- and it is cited here rather than reproved.

What that buys is the elimination of the filtering recursion that dominates
inference in most imperfect-information card games. There is no belief that
decays and no dynamics to approximate, only a fixed hidden object and an
accumulating set of logical constraints on it. Every player computes the
posterior from the same public history but conditions additionally on their own
private hand, so the six posteriors differ and no single distribution is shared.

The constraints are accumulating rather than merely present, and the rate at
which they accumulate governs whether an exact endgame treatment is worth
building. A game reaches a decided state after a mean of \vsixCollapseEvents{}
public events, which is the figure \S\ref{sec:mechanisms} uses to size the
endgame regime.

\paragraph{Notation.} $\mathcal{H}$ is the set of live half-suits and $H \in
\mathcal{H}$ one of them. Fix an observer $i$. $U$ is the set of cards in live
half-suits whose holder is not publicly known to $i$, and for $c \in U$ the map
$\sigma : U \to \{0,\dots,5\}$ gives that card's holder, constant in time.
$M_c \subseteq \{0,\dots,5\}$ is the surviving support of $\sigma(c)$, and $q_p$
is the capacity of seat $p$, the number of cards of $U$ that seat must hold. $Z$
is the number of assignments $\sigma$ consistent with the constraints below, the
partition function of the posterior. The three quantities the policy consumes are
the per-card ownership marginal $\mu_{c,p} = \Pr[\sigma(c) = p]$; the probability
$\tau(H,T)$ that team $T$ owns every card of half-suit $H$; and the probability
$\alpha(A)$ that a named allocation $A$ --- an assignment of all six cards of one
half-suit to named seats --- is correct. Unhatted symbols are values from the
exact engine, hatted ones ($\hat\mu_{c,p}$, $\hat\tau$, $\hat\alpha$) the same
values from the approximate path. Probabilities are conditional on the observer's
information throughout and the conditioning bar is suppressed.

\paragraph{The constraint system.} The observer's information is exactly five
constraints on $\sigma$.

\begin{description}[style=unboxed,leftmargin=0pt,itemsep=4pt,parsep=0pt]
  \item[(C1) Own hand.] $\sigma(c) \ne i$ for every $c \in U$, and every card $i$
        holds is resolved.
  \item[(C2) Public transfers and correct declarations.] Any card whose holder
        has been revealed leaves $U$ with its holder known exactly. A successful
        ask reveals the holder \emph{before} the transfer, and that is the value
        relevant to every earlier constraint.
  \item[(C3) Exclusions.] An ask for $c$ proves the asker did not hold $c$; a
        miss proves the target did not hold $c$. Both are permanent, because
        unresolved cards never move. The surviving support $M_c$ records them.
  \item[(C4) Capacities.] Hand counts are public, so for every player $p$
        \begin{equation}
          \bigl|\{c \in U : \sigma(c) = p\}\bigr| \;=\; q_p \;:=\; h_p - k_p ,
          \label{eq:capacity}
        \end{equation}
        where $h_p$ is $p$'s public card count and $k_p$ the number of live cards
        already known to be $p$'s. Note $\sum_p q_p = |U|$ automatically, and
        that a declaration --- correct or not --- is absorbed by
        \eqref{eq:capacity} without special handling, because the count drop it
        causes is public.
  \item[(C5) Ask legality.] An ask for $c$ in half-suit $H$ by player $a$ proves
        that $a$ held some \emph{other} card of $H$ at that moment. Because
        unresolved cards never move, this is the time-independent disjunction
        \begin{equation}
          \bigvee_{d \in D} \bigl[\sigma(d) = a\bigr],
          \qquad D = \{d \in H \setminus \{c\} : d \text{ unresolved},\; a \in M_d\},
          \label{eq:disj}
        \end{equation}
        vacuous if some card of $H$ is already known to be $a$'s. Every such
        certificate is confined to a single half-suit, which is the structural
        fact \S\ref{sec:belief-count} exploits.
\end{description}

The prior is uniform over deals consistent with the rules, so the posterior is
uniform over the assignments $\sigma$ satisfying (C1)--(C5) and $\Pr[\sigma] =
1/Z$ for each of them. That sentence contains the whole of
\S\ref{sec:belief-exactworse} in advance: the object is exact with respect to
\emph{that} prior, and the prior is a modelling choice, not a rule of the game.

\subsection{The exact count law}
\label{sec:belief-count}

Set (C5) aside. Counting assignments subject to (C3) and (C4) is a
degree-constrained bipartite counting problem, solved exactly by a
forward--backward recursion over the vector $n = (n_0,\dots,n_5)$ of
\emph{remaining} capacities. Ordering $U$ as $c_1,\dots,c_{|U|}$, $F_j(n)$ counts
placements of the first $j$ cards leaving $n$ unused and $B_j(n)$ placements of
the rest out of $n$:
\begin{align}
  F_j(n) &= \sum_{p \,\in\, M_{c_j} :\, n_p < q_p} F_{j-1}(n + e_p),
  &F_0(q) &= 1, \label{eq:fwd}\\
  B_j(n) &= \sum_{p \,\in\, M_{c_{j+1}} :\, n_p \ge 1} B_{j+1}(n - e_p),
  &B_{|U|}(\mathbf{0}) &= 1, \label{eq:bwd}\\
  Z &= F_{|U|}(\mathbf{0}) \;=\; B_0(q), \label{eq:Z}
\end{align}
with $e_p$ the $p$-th unit vector. Every card consumes exactly one unit of
capacity, so the layer index is recoverable from $n$ and the $|U|$-layer table
collapses to two arrays of at most $10^5$ entries (\S\ref{app:inf-count}). These
recursions count rather than sample, so they and the per-card marginal they
contract to are exact for (C1)--(C4).

\paragraph{Ask legality by half-suit block enumeration.} Constraint (C5) is
disjunctive and breaks the product form of \eqref{eq:fwd}--\eqref{eq:bwd}, and
inclusion--exclusion over the tens of certificates a game accumulates would cost
$2^k$ passes. The escape is structural: every certificate lives inside one
half-suit, so $U$ partitions into half-suit blocks admitting one of two exact
treatments. A block with no live certificate is expanded card by card as in
\eqref{eq:fwd}--\eqref{eq:bwd}. A block with a live certificate is enumerated
exhaustively --- at most $\prod_{c \in H}|M_c| \le 5^6$ assignments --- each
tested against \eqref{eq:disj} directly, and the survivors aggregated into a
count-vector table
\begin{equation}
  g_H(t) \;=\; \bigl|\{\text{assignments } \beta \text{ of } H :
    \mathrm{count}(\beta) = t,\; \beta \models \text{certificates of } H\}\bigr| ,
  \label{eq:blocktable}
\end{equation}
where $\mathrm{count}(\beta)_p$ is the number of $H$'s unresolved cards $\beta$
gives to seat $p$. Either way a block consumes a known number of cards, so the
layer identity survives and the dynamic program runs over blocks rather than
cards. Writing $S_t$ for the path mass of count vector $t$ at its block's entry
layer,
\begin{equation}
  S_t \;=\; \sum_{|n| \,=\, \text{entry layer}} F(n)\, B(n - t),
  \label{eq:pathmass}
\end{equation}
the three quantities of \S\ref{sec:belief-state} are read straight off the
tables:
\begin{align}
  \mu_{c,p} &= \frac{1}{Z}\sum_t g_H^{(c \to p)}(t)\, S_t
  &&\text{(per-card ownership)} \label{eq:q1}\\
  \alpha(A) &= \frac{S_{t(A)}}{Z}
  &&\text{(the declaration query)} \label{eq:q2}\\
  \tau(H,T) &= \frac{1}{Z}\sum_{t \,:\, \mathrm{supp}(t) \subseteq T} g_H(t)\, S_t
  &&\text{(does one team hold the half-suit?)} \label{eq:q3}
\end{align}
All three are exact under the uniform prior; no independence assumption and no
sampling enters. \S\ref{app:inf-blocks} gives the mechanics of both block kinds.

\paragraph{Exactness here is cheap.} An exact rebuild costs \vsixCountCost{}~ms
mean over \vsixCountStates{} states, against \vsixSinkhornCost{}~microseconds for
the approximate path. That ratio is large but the absolute cost is small, and it
is small enough that the result of \S\ref{sec:belief-exactworse} is not a budget
story. The exact posterior is not too expensive to deploy. It is worse.

\paragraph{What is verified and what is merely unfalsified.} Two exactness claims
in this machinery are of different standing, and the distinction is stated here
rather than buried in the appendix. The count law and the block enumeration are
checked against a brute-force enumerator that factorises nothing
(\S\ref{app:inf-audit}), and the agreement is at integer rather than rounding
scale. The team-allocation search, by contrast, does not re-check the survival
constraint of \eqref{eq:disj} on the allocation it returns; it has failed zero
times, which is not the same as a proof. \S\ref{app:inf-latent} lists that case
and two others --- a count routine that falls back to a greedy approximation for
heterogeneous non-singleton masks, and an enumeration truncation that leaves a
partial table which is then normalised as if complete --- each with its
reachability argument.

\paragraph{The approximate path.} Six observers invoke the posterior after every
public event, which puts it in the inner loop of every large experiment, so the
engine also implements an approximation and that is what ships. The Fast path
keeps the exact bookkeeping of (C1)--(C3) and the exact capacities (C4), fits
them by Sinkhorn scaling of the support matrix, and interleaves a conditioning
step for (C5) that reweights a certificate's cards by treating their entries as
independent Bernoulli indicators before capacity fitting is re-run
(\S\ref{app:inf-sinkhorn}). The independence assumption is false, because the
capacity constraints that make \eqref{eq:q2} necessary also couple the cards
inside a certificate. Every performance number in this report was produced by
this path.

\subsection{The policy prior}
\label{sec:belief-prior}

The posterior of (C1)--(C5) is exact with respect to the rules and is not the
full Bayesian posterior over deals, which would additionally weight each deal $x$
by the likelihood that the observed actions would have been chosen given $x$.
That likelihood is playstyle-dependent: a player who has taken many turns and
never asked in half-suit $H$ is, under any plausible policy, less likely to hold
cards of $H$, and the rules alone say nothing about this.

The engine implements a two-parameter family of such corrections as prior weights
inside the same machinery,
\begin{equation}
  w(c,p) \;=\; \exp\!\bigl(\theta\, a_{p,H(c)} - \phi\,(a_p - a_{p,H(c)})\bigr),
  \label{eq:prior}
\end{equation}
where $a_{p,H}$ counts $p$'s public asks in half-suit $H$, $a_p$ their total
public asks, and $H(c)$ is the half-suit of card $c$. The weights replace the
uniform initialisation of the Sinkhorn fit, so $\theta = \phi = 0$ recovers the
policy-agnostic posterior of \S\ref{sec:belief-state}. Both parameters are
fitted jointly with the rest of the policy (\S\ref{sec:fitting}); the shipped
values are $\theta = \vsixPriorTheta{}$ and $\phi = \vsixPriorPhi{}$. Deleting
the prior from the shipped configuration, leaving every other coordinate at its
fitted value, costs \vsixPriorDeleteCost{} points, interval
[\vsixPriorDeleteCI{}].

Two architectural facts follow, and they are what makes
\S\ref{sec:belief-exactworse} possible. \emph{First}, \eqref{eq:prior} is the
only channel by which behaviour enters the belief; everything else in
(C1)--(C5) is a rule. \emph{Second}, the approximate path has that channel and
the exact path structurally cannot: the block construction of
\S\ref{sec:belief-count} takes the deduction state alone and has no parameter
through which an opponent model could reach it. The two paths are therefore not
a fast approximation and its exact reference. They are two different posteriors
over the same hidden state, one policy-agnostic and exact under a uniform deal
prior, the other policy-aware and approximate.

\subsection{The exact posterior is a worse predictor than the approximation}
\label{sec:belief-exactworse}

Both posteriors were scored as predictors, on the same states, by two measures
that do not involve playing a game: predictive log loss over the unresolved cards
of each state, and the realised hit rate of each posterior's own argmax ask ---
the ask the posterior itself rates most likely to succeed, scored against the
truth. The probe is \texttt{fish v6probe --mode=belief}
(\S\ref{sec:engine-impl}), and the source of record is
\filepath{research/v06/notes/R12-mechanism-trials.md}, \S3.

The comparison carries no fitted weights on either side. It scores posteriors,
not policies, so nothing in it depends on which posterior the ask weights were
tuned against. That is why it can settle a question that a substitution
experiment in play cannot.

\begin{table}[ht]
\centering
\caption{Each posterior scored as a predictor on identical states: mean
predictive negative log likelihood over the unresolved cards, the probability
each posterior assigns to its own argmax ask, and the realised hit rate of that
ask. Rows are the exact count law under the uniform deal prior, the approximate
path with the policy prior deleted, and the approximate path at the shipped
prior, followed by the $\theta$ sweep. Lower NLL and higher hit rate are better.
The unit is one card for the NLL column and one decision for the hit column.}
\label{tab:belief}
\small
%
%% --- begin tables_v06/belief ---
\begin{tabular}{lrrrr}
\toprule
Posterior & Cards & Mean NLL & Argmax $p$ & Argmax hit \\
\midrule
exact (uniform prior) & 140,661 & 1.42246 & 0.4660 & 47.94 \\
sinkhorn th=0.000 ph=0.122 & 140,661 & 1.39083 & 0.4896 & 49.99 \\
sinkhorn th=0.200 ph=0.122 & 140,661 & 1.38663 & 0.4939 & 50.75 \\
sinkhorn th=0.300 ph=0.122 & 140,661 & 1.38465 & 0.4972 & 51.06 \\
sinkhorn th=0.445 ph=0.122 & 140,661 & 1.38218 & 0.5035 & 51.49 \\
sinkhorn th=0.600 ph=0.122 & 140,661 & 1.38055 & 0.5113 & 51.59 \\
sinkhorn th=0.800 ph=0.122 & 140,661 & 1.38055 & 0.5211 & 51.64 \\
sinkhorn th=1.100 ph=0.122 & 140,661 & 1.38311 & 0.5336 & 51.64 \\
sinkhorn th=1.500 ph=0.122 & 140,661 & 1.38758 & 0.5459 & 51.51 \\
sinkhorn th=2.000 ph=0.122 & 140,661 & 1.39221 & 0.5555 & 51.40 \\
\bottomrule
\end{tabular}
%% --- end tables_v06/belief ---
\end{table}

The ordering is unambiguous and it is the wrong way round. The exact count law
under a uniform deal prior scores \vsixBeliefExactNLL{} nats per card and
\vsixBeliefExactHit{}\% argmax hit rate. The approximate path with only the
certificate half of the prior deleted ($\theta = 0$) scores \vsixBeliefNoPriorNLL{} and
\vsixBeliefNoPriorHit{}\%; with the prior deleted outright ($\theta = \phi = 0$,
\filepath{research/v06/results/F2-belief-noprior.txt}) it scores $1.39339$ and $49.14\%$ --- still
ahead of the exact law on both measures. The shipped prior scores \vsixBeliefShippedNLL{} and
\vsixBeliefShippedHit{}\%. The exact object is the worst predictor of the three,
and the gap between the second and third rows is the policy prior alone:
\numThetaPriorHitGain{} points of argmax hit rate and \numThetaPriorNatGain{}
nats.

The same ordering appears on the full live candidate set rather than on the card
marginals. Over the ask decisions of a \vfive{} mirror, the approximate argmax
converts \vsixFullFastFive{}\% of the time and the exact argmax
\vsixFullExactFive{}\%, on identical states, with the two argmaxes disagreeing at
\vsixDisagreeFive{}\% of decisions; hindsight converts
\vsixFullBestFive{}\%. Substituting the exact posterior into the shipped policy
therefore does not merely fail to help. It replaces the better predictor with the
worse one at every decision on which they differ.

Three consequences follow, and the note draws all three.

\begin{enumerate}
  \item \textbf{``Exact'' is exact under a uniform-deal prior, and that prior is
        wrong.} The exact count law is the worst predictor in
        Table~\ref{tab:belief}. Part of the gap is \eqref{eq:prior}: deleting
        the prior costs the approximation about a point and a half of argmax hit
        rate. But it is not the whole gap --- with both prior parameters deleted
        the approximation still predicts better than the exact law, so the
        remainder is the Sinkhorn fit's own maximum-entropy smoothing, which
        hedges better than an exact conditional distribution taken under a prior
        that is wrong. Exactness with respect to the
        rules is not exactness with respect to the game, because the observed
        actions carry information the rules do not encode. Calling the block
        posterior ``the exact posterior'' without that qualification --- as the
        earlier reports in this corpus do --- names it after a property it has
        only relative to an assumption nobody tested.
  \item \textbf{The residual belief error is not an approximation error to be
        engineered away.} This settles the corpus's standing open question
        --- whether a matched-budget refit of the ask weights under
        \texttt{belief=block} would recover the deficit the v0.4 substitution
        ablation measured --- in the negative. The refit is not worth running,
        because the object it would fit is a strictly less informative posterior.
        The question stayed open for two studies because it was posed as a
        budget question, and it is not one: \S\ref{sec:belief-count} shows the
        exact path is affordable, and this subsection shows it is worse anyway.
  \item \textbf{The predictive optimum is not the playing optimum.} Raising
        $\theta$ improves prediction well past its fitted value: the sweep in
        Table~\ref{tab:belief} minimises NLL at $\theta \approx
        \numThetaPredictiveOpt{}$, against the shipped $\numThetaShipped{}$. In
        play the same change loses --- \numThetaSixtyCost{} points of win rate at
        $\theta = \numThetaTrialLow{}$ and \numThetaElevenCost{} at $\theta =
        \numThetaTrialHigh{}$, at \numThetaTrialN{} games per cell. The ask weights were fitted at the
        shipped $\theta$ and do not transfer to a different one. Belief and
        policy are therefore not separable objects to be optimised in turn; any
        change to $\theta$ has to be refitted jointly with the weights that
        consume it, which is what the fit of \S\ref{sec:fitting} does.
\end{enumerate}

The third consequence is the one that generalises beyond this game. A better
posterior is not a better agent, and an inference component evaluated on
predictive loss can be improved into a policy that plays worse. The two
objectives coincide only when the downstream policy is optimal for whatever
posterior it is handed, and no fitted linear policy is.
%% --- end sections_v06/05-belief ---
%% --- begin sections_v06/06-diagnosis ---
\section{Diagnosis: what \vfive{} does wrong, and what it does not}
\label{sec:diagnosis}

The v0.5 study closed with an unusually candid non-claim: removing a failure mode was worth almost
nothing in win rate, and that was the result. This section asks the obvious follow-up question ---
where, then, is the remaining value --- and answers it with four measurements, two of which are
negative and are reported as findings rather than as caveats. The negative pair is load-bearing:
between them they close off the two largest items on the design register this study inherited, and
they are the reason \vsix{} spends its effort on the measurement apparatus rather than on the
policy's feature set.

\subsection{The instrument}
\label{sec:diag-instrument}

Every claim below is produced by \texttt{fish v6probe}, a diagnostics subcommand added for this
study, and is reproducible from a clean build by the commands in Appendix~\ref{app:repro}. Three
definitions fix what is being counted.

\begin{definition}[Live candidate]
An ask $(c,t)$ is \emph{live} for an actor if it is legal and the actor cannot prove from its own
hand and the public record that $t$ does not hold $c$. This is the M1 predicate \vfive{} already
gates on; it involves no posterior and no policy prior, so a candidate set restricted to it cannot
have been narrowed by a modelling error.
\end{definition}

\begin{definition}[Exact tie at the top]
A decision has an \emph{exact tie at the top} when two or more live candidates receive scores that
are equal to within $10^{-9}$ under the policy's own pre-search score, that is
$\lambda\, w^{\!\top}\! f + \nu\, Q$ where $Q$ is the one-ply expectimax over the value function.
The gap and the spread are reported alongside so that a tie can be distinguished from a flat
objective.
\end{definition}

\begin{definition}[Belief-symmetric tie]
A tie is \emph{belief-symmetric} when the tied candidates additionally receive equal probability
under the \emph{exact} posterior of \S\ref{sec:belief}, not merely under the Sinkhorn fit the
deployed policy runs. A belief-symmetric tie is a statement about the information set, not about
the approximation.
\end{definition}

\subsection{More than half of all ask decisions end in a tie}
\label{sec:diag-ties}

Over \vsixTieNFive{} tied decisions in \vfive{} mirror play, \vsixTieShareFive\% of contested
decisions --- decisions with at least two live candidates, which is \vsixPathFiveContested\% of all
of them --- carry an exact tie at the top. The objective is not flat: the mean spread across the
candidate set is \vsixTieSpreadFive{} while the mean gap between the leaders is
\vsixTieGapFive{}. The dimension is blind, not the surface.

The tie is structural, and it is structural in a specific place. \vsixTieSameSuitFive\% of ties are
two \emph{cards of one half-suit at one target}; only \vsixTieSameCardFive\% are one card offered to
two different targets. Three causes stack. Eighteen of the twenty ask features are computed per
(half-suit, target) and the three per-card features are all functions of the same marginal; the
Sinkhorn filter makes exchangeable cards numerically identical by construction; and
\texttt{askExpectedValue} discards its \texttt{target} argument outright, so the expectimax half of
the score is constant across targets. The last of these was the defect the v0.5 design register
headlined. It is aimed at \vsixTieSameCardFive\% of the tie mass.

\subsection{Negative result: the ties are exchangeable, and therefore irreducible}
\label{sec:diag-irreducible}

The natural reading of \S\ref{sec:diag-ties} is that a subroutine \vfive{} already runs is throwing
away a large amount of value, and the register this study inherited priced that value at over three
half-suits per game using a hindsight oracle. That reading is wrong, and the measurement that
refutes it is simple: re-score the tied candidates under the exact count law of
\S\ref{sec:belief-exactworse} and ask how many ties it separates.

\begin{center}
\begin{tabular}{lr}
\toprule
tie-break rule & realised hit rate \\
\midrule
enumeration order (whichever candidate was emitted first) & \vsixTieHitArrayFive\% \\
argmax of the deployed Sinkhorn marginal, within the tie & \vsixTieHitFastFive\% \\
the shipped policy's actual pick (its chain/threat pass) & \vsixTieHitShipFive\% \\
argmax of the \emph{exact} posterior, within the tie & \vsixTieHitExactFive\% \\
\midrule
hindsight best within the tie & \vsixTieHitBestFive\% \\
\bottomrule
\end{tabular}
\end{center}

The exact posterior separates \vsixTieExactSepPctFive\% of the ties, and every rule realises the
same hit rate to three significant figures. Two cards of one half-suit at one target, tied under the
policy's score, are tied under the exact posterior as well. The \vsixTieHitBestFive\% column prices
clairvoyance, not headroom.

\paragraph{The scope of that result, and the experiment that fixes it.} What the table measures is
equality of the \emph{marginal} probability that the target holds the card. It follows that no rule
ranking the tied candidates by that marginal --- the deployed one, the exact one, or any other ---
can beat enumeration order, and none does. An earlier draft of this section concluded from that
that the ties are irreducible. They are not, and the four rows below are the correction. Each plays
the same policy against \vfive{}, changing only how the tie group is resolved.

\begin{center}
\begin{tabular}{lr}
\toprule
how the tie group is resolved & win rate against \vfive{} \\
\midrule
enumeration order (the shipped rule) & $50.00$ by construction \\
uniformly at random, seeded from the public event stream & \vsixSearchTieRandom\% [\vsixSearchTieRandomLo, \vsixSearchTieRandomHi] \\
by a one-deal rollout & \vsixSearchTieOnlyOne\% [\vsixSearchTieOnlyOneLo, \vsixSearchTieOnlyOneHi] \\
\textbf{by a twelve-deal rollout ensemble} & \textbf{\vsixSearchTieOnlyTwelve\%} [\vsixSearchTieOnlyTwelveLo, \vsixSearchTieOnlyTwelveHi] \\
\bottomrule
\end{tabular}
\end{center}

Randomising the choice is worth \emph{exactly} nothing --- $50.00\%$, mean half-suits $4.500$ to
$4.500$ --- so this is not a story about a deterministic policy being predictable. One sampled deal
is worth nothing either. Twelve are worth several points at the bank in the table, and pooled over
every bank the configuration was measured on, \vsixSearchPooled\% over \vsixSearchPooledN{} games
with no bank below \vsixSearchPooledMin\%.

The separating information is therefore real, and it is not the mere fact of varying a
deterministic choice --- the random row rules that out exactly. It is in the \emph{joint}
posterior. Two cards can be equally likely to sit with the target while being differently correlated
with the rest of the deal --- differently placed in the allocations that survive the certificates,
differently useful to the half-suits the opponents are collecting --- and every marginal, exact or
not, integrates that structure away. An ensemble of whole sampled deals does not.

\paragraph{What is not resolved.} Whether the search's advantage is \emph{confined} to the tie
group is a separate question and this study does not answer it. Restricting the search to the tie
group gives \vsixTieSearchBankList\% over \vsixTieSearchBanks{} banks (mean
\vsixTieSearchMean\%); guarding the tie group
instead gives \vsixSearchTieGuardedWin\%; the unrestricted search gives \vsixSearchPooled\%. At
720 games a cell those intervals all overlap. What is established is the pair of negative controls
and the pooled effect, not the attribution. \S\ref{sec:search} says so again where it matters.

Two corrections to the record follow, and are carried into \S\ref{sec:corrections}.

\begin{enumerate}
\item \textbf{Enumeration order does not decide half of all asks.} \vfive{}'s top-$K$ chain and
  threat pass re-scores the leaders and moves the pick at \vsixTieShipMovedFive\% of ties, so array
  order settles roughly a fifth of contested decisions rather than more than half --- and, per the
  table, settles them exactly as well as any alternative.
\item \textbf{The channel is the card dimension, not the target dimension.} The emphasis the v0.5
  register placed on \texttt{(void)target;} is aimed at a twentieth of the tie mass.
\end{enumerate}

The design consequence is stated here and used twice later. Because the tie is a property of the
information set, a rollout ensemble drawn from that information set is symmetric over it too: the
determinizations carry no signal about which of two exchangeable cards to name, only variance. The
tie-break question is therefore not a search question, and \S\ref{sec:search} does not treat it as
one.

\subsection{Negative result: the \vfive{} fit was sampling noise}
\label{sec:diag-fitnoise}

The second negative result concerns the interpretation of \vfive{}'s refit rather than its shipped
configuration. Read from the fit trace, the incumbent objective moves over forty generations by
less than one generation's standard deviation, with an ordinary-least-squares slope whose $t$
statistic is below two, and the gap between the best candidate and the incumbent tracks the expected
maximum of a population of independent draws of exactly that noise.

The cause is visible in the sampler rather than in the objective. One scalar step size was applied
to every coordinate of a vector whose coordinates have ranges spanning two and a half orders of
magnitude, so the search was bang-bang on the bounded probability knobs --- the worst of them
clipped at \vsixFitAClipGenZero\% of generation-zero draws under the old setting --- and effectively
frozen on the unbounded weights. The per-coordinate step-size flag the interface advertised was
parsed and discarded. \S\ref{sec:fitting} is the repair, and \S\ref{sec:results} is what a working
optimiser then finds.

Taken with the independent observation that running the v0.5 \emph{mechanisms} on v0.4's frozen
parameter vector reproduces essentially the whole v0.5 result, the conclusion is that the
v0.4$\rightarrow$v0.5 transition was mechanical and that its parametric machinery never moved. That
is not a small correction: it means the shipped vector this study starts from is v0.4's, and v0.4's
was produced by a search carrying a separate aliasing defect of its own.

\subsection{What is actually still leaking}
\label{sec:diag-channels}

With the tie channel removed from the register and the parametric channel re-opened, the residual
loss channels are smaller than the inherited register suggested, and they are ranked here with the
paired figure each was measured at.

\begin{itemize}
\item \textbf{Asks into a half-suit the actor's own team already owns outright.}
  \vsixPathFiveOwnlocked\% of all asks --- guaranteed misses that donate the turn. M1 cannot see
  them: the enumerator only offers opponents as targets, and no single seat can \emph{prove} a
  teammate holds a card. Pricing the posterior mass the team already carries is the graded
  substitute, and \S\ref{sec:mechanisms} reports that it is worth nothing at the resolution
  \S\ref{sec:protocol} establishes as necessary.
\item \textbf{Declaration errors}, \vsixPathFiveDeclwrong\% of declarations in mirror play. The
  pre-gates that produce them are, as \S\ref{sec:engine} reports, exact: \vsixGateFalseNeg{} false
  negatives in \vsixGateRejected{} rejections. The error is in the allocation, not in the gate.
\item \textbf{The move class M1 deletes.} A provably-dead ask at a chosen opponent is how a human
  blackballs, donates the turn deliberately, or pays for a signal. \S\ref{sec:search} shows that
  \vfive{}'s score cannot price it even when it is offered, and that unbanning it wholesale trades
  a higher win rate for a policy that kills \vsixDeadPathLimit\% of its games at the action limit.
\end{itemize}

None of these is large. That is the diagnosis: \vfive{} sits on a plateau of its own policy class,
and the leverage available to \vsix{} is in the apparatus that measures and fits it, not in another
feature.
%% --- end sections_v06/06-diagnosis ---
%% --- begin sections_v06/07-mechanisms ---
\section{\vsix{}: the mechanism set}
\label{sec:mechanisms}

\vsix{} is implemented as a policy that \emph{derives} from \vfive{} rather than forking it, and
this is not a packaging decision. With every \vsix{} switch off and the v0.5 parameter vector
restored, each override defers to the base class and the two policies are identical --- checked, not
asserted, by comparing the md5 of the complete diagnostic output of a $60$-deal mirror run
(\vsixIdentity{}). Every ablation in \S\ref{sec:results} is therefore exact: there is no
``\vsix{} without mechanism $X$ differs from \vfive{} for unrelated reasons'' confound of the kind
that has cost this project two studies.

The mechanisms are grouped by what they change. Each carries the defect it was keyed to, its
measured verdict, and --- where the verdict is negative --- what that verdict rules out for the next
study. Four of the seven are reported as null, and they are reported here rather than buried in an
appendix, because the register they close is the register a reader of the v0.5 study would otherwise
work from.

\subsection{Shipped: the apparatus}
\label{sec:mech-cheap}

Three changes ship unconditionally because they are corrections, not strategies, and they make no
strength claim of their own.

\begin{itemize}
\item \textbf{The exact posterior's table aliasing} (\S\ref{sec:engine-repairs}), which made every
  exact-belief mechanism unsound in multi-agent play and which no strength measurement could ever
  have caught, because the deployed policy never queried the affected object.
\item \textbf{The repaired optimiser} (\S\ref{sec:fitting}): per-coordinate step sizes, a paired
  per-deal objective, an explicit minimax-regret dispatch, and a recovery test the previous harness
  provably fails.
\item \textbf{The commit gate} (\S\ref{sec:protocol}): the pathology KPIs are evaluated before any
  strength number is collected, because two of the highest-scoring rows of the v0.5 ablation table
  are configurations that kill $14\%$ of their games at the action limit.
\end{itemize}

\subsection{Measured and null: the tie-break}
\label{sec:mech-tiebreak}

The mechanism the inherited register ranked first was to resolve the exact ties of
\S\ref{sec:diag-ties} with a signal the Sinkhorn filter cannot represent --- specifically, the exact
count law, which costs about \vsixExactBuildUsFive{}~\textmu s to build and is affordable at a
decision. It was built, and \S\ref{sec:diag-irreducible} is its verdict: the exact posterior
separates \vsixTieExactSepPctFive\% of the ties, and every tie-break rule realises the same hit
rate. The mechanism is retained as an \emph{instrument} and not as a strategy, and the distinction is
literal: the exact-posterior switch is not read by the ask rule and the shipped policy never builds
the exact posterior. What survives is the diagnostic that measured the null.

What this rules out for the next study is broader than the mechanism. Any tie-break candidate must
first be shown to separate the tied set \emph{under the exact posterior}. Signals that cannot ---
per-card ask history, rank priors, teammate-derived terms --- are re-descriptions of the same
symmetric information and cannot beat enumeration order.

\subsection{Measured and null: three ask terms}
\label{sec:mech-knowledge}

Three terms were added to the ask score, each fitted rather than hand-set, and each keyed to a
measured channel.

\begin{itemize}
\item \textbf{Void creation}: the probability that a hit takes the target's \emph{last} card of the
  half-suit. Ask legality requires the actor to hold a card of the half-suit, so voiding an opponent
  permanently removes their right to ask in it --- a one-way asset that $41.65\%$ of \vfive{}'s
  successful asks create by accident while nothing in its score rewards creating one on purpose.
\item \textbf{Team-ownership discount}: the posterior mass the actor's own team already carries on
  the asked card, which is the graded substitute for the proof M1 cannot have and is aimed at the
  \vsixPathFiveOwnlocked\% of asks that are guaranteed misses into a half-suit the team already owns.
\item \textbf{The last-live-opponent split}, separating ``empty one of three opponents'' from
  ``empty the last one and trigger the forced endgame''.
\end{itemize}

All three are null as \emph{mechanisms}. One consequence has to be stated because it makes the
obvious ablation confounded: a non-zero extra weight is what selects the v0.6 scoring path, and the
v0.6 scoring path is the one that does not run v0.5's chain and threat re-scoring. Switching the
extra terms off therefore switches that pass back on. \S\ref{sec:results-ablations} reports the
four-way ablation that separates them. The void term cleared a $95\%$ paired interval at four
hundred games per cell on one bank and returned zero at a thousand on two; the last-live term fires
too rarely to change any output. The fit gives all three small coefficients and the ablation that
removes them outright is in Table~\ref{tab:ablations}. \S\ref{sec:protocol} is where that pattern
--- significant at the corpus's budget, zero at this one --- is turned into a protocol rather than
an anecdote.

\subsection{Measured and rejected: the deliberate miss}
\label{sec:mech-declare}

M1 deleted the move class every multi-step human tactic is built from. Re-admitting it raises the
win rate and destroys the policy (\S\ref{sec:search-deadask}); rationing it against the linear score
is inert, because the linear score provably cannot price it. It is offered to the search instead,
which is the only evaluator in the system that can, and it is the best-performing search
configuration measured. \S\ref{sec:search} reports what the search as a whole is worth.

\subsection{Shipped off: test-time search}
\label{sec:mech-leaves}

\S\ref{sec:search} in full. Off, instrumented, and reported with the conditions under which its
verdict would change.

\subsection{Shipped: the fitted vector}
\label{sec:mech-partner}

What is left, once four mechanisms have been measured null and one rejected, is the parameter vector
--- and \S\ref{sec:diag-fitnoise} is the reason that is not an anticlimax. \vfive{}'s vector is
v0.4's, because v0.5's forty-generation refit moved nothing, and v0.4's was produced by a search
carrying an aliasing defect of its own. Nobody has ever fitted this policy class with an optimiser
that works. \S\ref{sec:fitting} builds one and \S\ref{sec:results} reports what it finds.

\subsection{Reported separately: the partner regimes}
\label{sec:mech-exploit}

The evaluation harness built one policy object for all three seats of a team, so ``\vsix{} with two
bot partners'' and ``\vsix{} with two human-model partners'' --- the two regimes the project owner
asked to see reported separately --- could not be expressed at all. Per-seat specification is added
and both regimes appear as separate columns in \S\ref{sec:results-panel}. The self-play row is never
the headline.
%% --- end sections_v06/07-mechanisms ---
%% --- begin sections_v06/08-search ---
\section{Test-time search, and the statistic that decides whether it helps}
\label{sec:search}

Every prior FishBot chooses its ask by maximising a fitted score over the legal set. The obvious
next step, and the one the project owner asked for, is to reason forward: to evaluate a candidate by
what happens after it rather than by a static function of the position it is played from. This
section builds that and measures it. It is the one mechanism in this study that beats a static
rule --- and getting there required correcting a statistic whose cost, \vsixCurseGap{} points, is
the most transferable result in the paper.

\subsection{Why determinization here is not double-dummy search}
\label{sec:search-notpimc}

The v0.4 study eliminated perfect-information Monte Carlo on structural grounds and this study
re-confirms them: under perfect information the team on turn takes every half-suit not dealt outright
to its opponents, so a double-dummy evaluator scores every legal action alike and the average over
determinizations cannot discriminate. Fish is a game about who knows what; a clairvoyant evaluator
prices none of it.

The search built here determinizes the \emph{deal} --- which is the entire hidden state because every card movement in this game is public --- but does not make anybody clairvoyant. Each of the six players in the
continuation is reconstructed \emph{at its own information set}: the public deduction state, which
is common knowledge and is therefore one object rather than six, refined by the hand that
determinization gives that seat. They miss asks, emit certificates and misdeclare in the proportions
the real game does. This is the cooperative single-agent search of Lerer et al.\
\cite{lerer2020sparta}, with the belief taken from an exactly enumerable posterior rather than a learnt
one.

The reconstruction is checked against ground truth rather than argued. Over \vsixReconStates{} states and \vsixReconChecks{} card checks, the reconstructed
information set differs from the seat's real one on \vsixReconMaskPct\% of cards and is \emph{wider}
in \vsixReconWider{} of \vsixReconWider{} cases (narrower in \vsixReconNarrower{}), replicated at
\vsixReconBanks{} seed banks: it is a strict under-approximation of
what the seat knows, so no simulated player is ever credited with knowledge it does not have.

\subsection{The search}
\label{sec:search-alg}

At an ask decision the actor
(i) builds the exact capacity dynamic program on its own deduction state and draws $D$ deals from
it, rejecting draws that violate an ask-legality certificate, so accepted draws are exact samples of
the policy-agnostic posterior, and re-weighting them by the policy prior as an importance weight so
that the sampler stays exact and the prior stays auditable;
(ii) takes the $K$ leading candidates under the blueprint score, extended through every candidate
that ties the $K$-th, because a plain top-$K$ would cut an exact tie arbitrarily and so reproduce the
very defect the search is being asked about;
(iii) seats six blueprint agents at their reconstructed information sets, applies the candidate, and
plays the continuation out;
(iv) and repeats every candidate on the \emph{same} determinizations, so the comparison is paired.

Two design choices are forced by measurement rather than taste. The blueprint used inside the
rollout is the deployed policy with its two-ply refinement disabled and its Sinkhorn iteration count
reduced; Table~\ref{tab:rollout} gives the cost/fidelity frontier, which the corpus had not swept.
The cheapest blueprint available is $62\times$ cheaper per event than the deployed policy and
$38$ points weaker, and a search whose continuation model is $38$ points weak is searching a
different game.

\begin{table}[t]\centering
\caption{Rollout-blueprint fidelity. Win rate against \vfive{} and single-thread throughput.
Reducing the Sinkhorn iteration count, which the corpus had not tried, is the cheap axis: it buys
most of the speed of the independent-marginal blueprint at a small fraction of its strength cost.}
\label{tab:rollout}
%% --- begin tables_v06/rollout ---
\begin{tabular}{lrrr}
\toprule
Rollout blueprint & Win rate vs \vfive{} & 95\% CI & Games/s \\
\midrule
\texttt{v05:belief=indep,topk=0} & 12.00 & [10.33, 13.75] & 996.2 \\
\texttt{v05:topk=0,souter=1,sinner=1} & 46.08 & [43.33, 48.83] & 146.8 \\
\texttt{v05:topk=0,souter=1,sinner=3} & 46.75 & [44.00, 49.50] & 114.9 \\
\texttt{v05:topk=0,souter=2,sinner=4} & 51.42 & [48.67, 54.17] & 78.7 \\
\texttt{v05:topk=0} & 47.75 & [44.92, 50.50] & 30.7 \\
\bottomrule
\end{tabular}
%% --- end tables_v06/rollout ---
\end{table}

\subsection{The result: the optimizer's curse is worth \vsixCurseGap{} points}
\label{sec:search-curse}

The first working build of the search chose the candidate with the highest mean rollout return and
scored \vsixSearchUnguarded\% against \vfive{}. The control that isolates the cause is the same
code, the same rollout blueprint, the same budget, the same seed bank and the same sample size, with
the blueprint's preference made to dominate: it scores \vsixSearchControl\%. The machinery is
correct and the \emph{statistic} was wrong, and the gap between those two rows ---
\vsixCurseGap{} points --- is the price of the argmax.

One caveat about that control, because it is the row the whole comparison rests on. The blueprint's
preference dominates only where the blueprint \emph{has} one: inside its own tie group the added
term is identically zero and the deviation penalty is waived, so the control still lets the search
choose among tied candidates. It scores at parity, which is a second, independent confirmation of
\S\ref{sec:diag-irreducible}: searching inside an exchangeable tie group changes nothing.

One determinization's return is the final half-suit differential, whose standard deviation is about
two and a half half-suits, while genuine differences between the leading candidates are an order of
magnitude smaller. An argmax over $D$ such means therefore selects on residual noise and lands on a
candidate the blueprint had ranked below the top. Replacing the argmax with a paired
lower-confidence-bound rule --- deviate from the blueprint's own choice only when the improvement
over it, on the same determinizations, clears $\kappa$ standard errors of the paired difference ---
recovers the whole deficit monotonically in $\kappa$ (Table~\ref{tab:search}).

\begin{table}[t]\centering
\caption{The deviation-threshold ladder, six rotations, held-out bank. The first row is the control
that forces the blueprint to decide. The range across the ladder is \vsixSearchSpread{} points.}
\label{tab:search}
%% --- begin tables_v06/search ---
\begin{tabular}{lrrr}
\toprule
Configuration & Win rate vs \vfive{} & 95\% CI & $n$ \\
\midrule
\texttt{s1=1,det=8,cand=6,blend=1000000} & 49.31 & [45.83, 52.78] & 720 \\
\texttt{s1=1,det=8,cand=6,kappa=0} & 13.61 & [10.97, 16.39] & 720 \\
\texttt{s1=1,det=12,cand=4,kappa=0} & 23.33 & [20.42, 26.39] & 720 \\
\texttt{s1=1,det=12,cand=4,kappa=1} & 42.78 & [39.31, 46.25] & 720 \\
\texttt{s1=1,det=12,cand=4,kappa=2.5} & 53.75 & [50.42, 57.22] & 720 \\
\texttt{s1=1,det=12,cand=4,kappa=4} & 53.89 & [50.42, 57.36] & 720 \\
\texttt{s1=1,det=16,cand=6,kappa=2,maxq=26} & 48.89 & [45.14, 52.64] & 720 \\
\texttt{s1=1,det=16,cand=6,kappa=2,maxq=26,deadsearch=2} & 50.28 & [46.67, 53.89] & 720 \\
\bottomrule
\end{tabular}
%% --- end tables_v06/search ---
\end{table}

This is a property of determinized evaluation in this game rather than of this implementation, and
it is stated as the section's main contribution. Any search over a high-variance terminal return
that selects by an unguarded argmax will pay it, and the price here --- \vsixSearchSpread{} points
across the ladder --- is larger than every mechanism the three prior studies fitted, combined.

\subsection{Where the advantage lives}
\label{sec:search-falsifier}

Guarded, the search stops losing and starts winning: at $\kappa = 2.5$ it takes
\vsixSearchPooled\% against \vfive{} pooled over \vsixSearchPooledBanks{} banks and
\vsixSearchPooledN{} games. The remaining
question is whether that win is the kind a working search produces, and the deviation rate answers
it --- but only once the rate is decomposed. The falsifier the literature supplies is explicit: a healthy test-time search differs from its
blueprint on a small percentage of decisions, and one that differs on tens of percent is mis-scaled
rather than smart. At $\kappa = 2.5$ the search moves the action on \vsixSearchDevTwoFive\% of the decisions it
searches (Table~\ref{tab:searchdev}), and that rate does not fall below about $31\%$ however far
$\kappa$ is raised. Read naively this is an order of magnitude above the band the literature
associates with a search that refines a good policy, and it was so read in an earlier draft. It is
the wrong reading, and the reason is \S\ref{sec:diag-irreducible}.

The deviation penalty is waived inside the blueprint's own tie group, because there the blueprint
expresses no preference. That group is more than half of all decisions, so the search always
re-chooses inside it, and those re-choices are neutral by construction --- the candidates are
exchangeable. The floor of $31\%$ is that churn, not signal. Applying the same penalty inside the
tie group collapses the rate to \vsixSearchDevGuardedTwoFive\% at $\kappa = 2.5$ and
\vsixSearchDevGuardedFour\% at $\kappa = 4$.

That decomposition also bears on where the win sits. Applying the penalty inside the tie group --- so that the search deviates only on the
\vsixSearchDevGuardedTwoFive\% of decisions where it has evidence against the blueprint's stated
preference --- takes the win rate to \vsixSearchTieGuardedWin\%
[\vsixSearchTieGuardedLo, \vsixSearchTieGuardedHi]: the advantage disappears. Forbidding deviation
\emph{outside} the tie group instead, so the search only ever re-chooses among candidates the
blueprint cannot separate, takes it to \vsixSearchTieOnlyTwelve\%
[\vsixSearchTieOnlyTwelveLo, \vsixSearchTieOnlyTwelveHi] --- better than the unrestricted search.

\textbf{What is established, and what is not.} The pooled effect is established twice over. On
\vfive{}'s parameter vector the guarded search takes \vsixSearchPooled\% against \vfive{} over
\vsixSearchPooledN{} games and \vsixSearchPooledBanks{} banks, none below
\vsixSearchPooledMin\%. On \vsix{}'s own vector --- so that the search is being asked to improve
the best static policy this project has, not the previous one --- it takes \vsixSearchLift\%
against \vsix{} over \vsixSearchLiftN{} games, with \vsixSearchLiftAbove{} of
\vsixSearchLiftCells{} cells above parity and the worst at \vsixSearchLiftMin\%. That second
result is what makes \vsixsearch{} a configuration worth naming. So are the two negative
controls of \S\ref{sec:diag-irreducible}: randomising the tie group is worth exactly zero and one
sampled deal is worth nothing, so what the ensemble is reading is joint information and not the
mere fact of varying a deterministic choice.

What is \emph{not} established is where inside the search the effect sits. Guarding the tie group
gives \vsixSearchTieGuardedWin\%; restricting the search to the tie group gives
\vsixTieSearchBankList\% across \vsixTieSearchBanks{} banks; the unrestricted search gives
\vsixSearchPooled\%. At 720 games a cell those intervals overlap each other and the unrestricted
figure, and an earlier draft of this section read a single favourable bank as an attribution. It is
not one. Resolving it needs roughly an order of magnitude more games at a configuration that costs
\vsixSearchGps{} games per second on one thread, against \vsixThroughputSix{} for the deployed
policy across all threads --- which is the reason it is not resolved here.

\begin{table}[t]\centering
\caption{Search deviation rate against the deviation threshold. A healthy search moves a small
percentage of decisions; this one either moves a third of them or none.}
\label{tab:searchdev}
%% --- begin tables_v06/searchdev ---
\begin{tabular}{lrrr}
\toprule
$\kappa$ & Decisions searched (\%) & Moved the action (\%) & Games/s \\
\midrule
0 & 98.58 & 65.67 & 0.173 \\
1 & 98.51 & 41.51 & 0.165 \\
2.5 & 98.51 & 33.49 & 0.144 \\
4 & 98.32 & 31.40 & 0.144 \\
6 & 98.32 & 31.59 & 0.144 \\
\bottomrule
\end{tabular}
%% --- end tables_v06/searchdev ---
\end{table}



\subsection{The one move class that exists only under search}
\label{sec:search-deadask}

M1 restricts the candidate set to asks the actor cannot prove will miss. That fixed v0.4's
two-question cycle, and it deleted the move class every multi-step human tactic is built out of:
blackballing, choosing which opponent receives the turn, and paying for a signal with a safe ask are
all a guaranteed miss at a \emph{chosen} seat. Such an ask exists at $79.2\%$ of \vfive{}'s
decisions and at two or more distinct chosen opponents at $50.1\%$.

Unbanning it wholesale is measured and rejected. With dead asks re-admitted and the ownership
features gated by hit probability so that the $p=0$ incentive behind v0.4's cycle cannot fire, the
configuration scores \vsixDeadWin\% [\vsixDeadWinLo, \vsixDeadWinHi] against \vfive{} --- higher ---
and fails the commit gate outright: \vsixDeadPathDead\% of its asks are provably dead, its longest
dead run is \vsixDeadPathLongest{}, and \vsixDeadPathLimit\% of its games are killed by the action
limit. This is the trap \S\ref{sec:protocol} exists to prevent, and it is the same trap the v0.5
ablation table contains.

Rationing it does not help either, and the reason is instructive: \emph{the linear score cannot
price a deliberate miss}. With the four ownership features zeroed on a dead candidate, the admission
margin was swept at four settings and the output was bit-identical at every one, because the hit
probability term is zero on such an ask by definition and every remaining term of the score is a
penalty. Its value is the value of the position the donated turn produces, which no static feature
of the ask computes.

The search can price it, and this is the clearest case in the study of a move class that exists only
under search: the deliberate miss is offered to the rollouts as a candidate, one per distinct target,
under an anti-repeat guard that is provably free. A dead $(c,t)$ pair stays dead unless the card
publicly moves to $t$, and that movement is observable, so forbidding the unmoved repeat forbids
nothing that could have become live. This is exactly the distinction the blanket repetition guard of
the v0.5 study missed, and why that guard cost six points: it also forbade \emph{live} repeats, which
are frequently the best ask on the board once a card has moved.

With the deliberate miss in the candidate set the endgame-restricted search scores
\vsixSearchDead\%. It is the best search configuration measured and it is still not separated from
the blueprint.

\subsection{What ships}
\label{sec:search-ships}

The search ships \textbf{off}, and the reason is cost and unresolved scope rather than a null
result. On the ladder's best setting it is ahead of \vfive{} with an interval that excludes parity,
but it is three orders of magnitude slower than the deployed policy and it moves a third of all
decisions --- so what it is is not a refinement of the
shipped policy but a distinct, far more expensive one whose evaluation this study could not complete
at the resolution \S\ref{sec:protocol} demands of everything else. It is retained, switchable and
fully instrumented. \S\ref{sec:limitations} states what would have to be measured before it could
ship on.
%% --- end sections_v06/08-search ---
%% --- begin sections_v06/09-fitting ---
\section{Fitting}
\label{sec:fitting}

The v0.5 study ran \numFitTraceGensVfive{} generations of a cross-entropy search over
\numVfiveParams{} coordinates and reported that the refit was worth nothing
measurable. \S\ref{sec:diag-fitnoise} shows that this was the correct reading of
the trace and the wrong conclusion to draw from it: the search made no progress
because the search was broken, in four separate and independently diagnosable
ways. This section presents the four defects, the repair for each, and --- for
each --- the measurement that shows the repair did what it was supposed to do
rather than the intention that it should.

The four defects are stated together first, because three of them compound.

\begin{enumerate}
  \item \textbf{One scalar step size for every coordinate.} A single
        $\sigma_0 = \numSigmaScalar{}$ was applied to all \numVfiveParams{}
        coordinates, whose admissible ranges span \numSigmaRangeMin{} to
        \numSigmaRangeMax{}. At generation zero this put
        \numClipDeclMargin{}\% of \texttt{declareMargin} proposals on a clamp
        bound while \texttt{valueWeight} moved \numValueWeightStep{}\% of its
        range: the optimiser was simultaneously bang-bang on the bounded
        probability knobs and frozen on the unbounded weight knobs. The
        command-line flag that would have fixed this, \texttt{--sigmaparams}, was
        parsed into a string that the rest of the program never referenced.
  \item \textbf{An unpaired objective.} Candidates were compared on absolute win
        rate against a common seed bank. Common random numbers pair the
        \emph{deal}; they do not pair the comparison, and the deal-level variance
        stayed in every candidate's score.
  \item \textbf{An objective documented as a soft minimum that is a weighted
        mean.} Differentiating the soft minimum shows its gradient is a convex
        combination of the panel's win rates with softmax weights, so at the
        compiled temperature it ranks candidates by mean win rate minus a small
        multiple of the spread. On \vfour{}'s shipped profile the
        maximum-to-minimum gradient weight ratio is \numSoftMinRatioFour{}
        against the $1$ a true mean would give and the $\infty$ a true minimum
        would give, and the ``minimum'' sits \numSoftMinBelowMin{} points away
        from the actual minimum. Two different temperatures were in use in one
        pipeline, and neither was recorded in any trace; the correction that owes
        to the earlier record is in \S\ref{sec:corr-self}.
  \item \textbf{The rotation defect.} The tuner computed each candidate's win
        rate as wins over \verb|games * 2|, hard-coding its own two rotations
        into the denominator (\S\ref{sec:engine-repairs}). Fitting at six
        rotations, which the duplicate-block design of \S\ref{sec:protocol}
        otherwise recommends, would have reported a third of the true win rate
        for every candidate, silently.
\end{enumerate}

Defects (i) and (ii) together explain the flat trace directly: an optimiser that
cannot move most of its coordinates, scoring what it can move through an
estimator that leaves the deal variance in, will produce a sequence of
generations whose best candidate is an order statistic of noise. Defect (iii)
means that what little signal there was pointed at the wrong objective. Defect
(iv) never fired, because the tuner never ran at six rotations, and is included
because it is a trap rather than a loss.

\subsection{Per-coordinate step sizes}
\label{sec:fit-sigma}

The repair is to derive each coordinate's step size from its own admissible
range. The optimiser now takes a relative step $\sigma_d = \sigma_{\mathrm{rel}}
\cdot (\mathrm{hi}_d - \mathrm{lo}_d)$, with an explicit per-coordinate vector
overriding it when one is supplied --- which is what \texttt{--sigmaparams} was
advertising and is now wired to. A coordinate's step is therefore a fixed
fraction of what that coordinate can be, rather than a fixed number of units of
whatever the coordinate happens to be measured in.

The check is not that the code changed. Every generation record in the run trace
now carries the per-coordinate clip fraction --- the fraction of that
generation's population whose proposal for that coordinate landed on a clamp
bound --- so the pathology the previous harness had can be read off the artifact
while a fit is running. At generation zero the worst coordinate of run~A clips at
\vsixFitAClipGenZero{}\%, against the \numClipDeclMargin{}\% the v0.5 trace
records for the same quantity. Run~C, which fits the wider v0.6 vector at a
smaller relative step, clips at \vsixFitCClipGenZero{}\%; that is reported as
measured rather than as passed, and \S\ref{app:params-bounds} prints the clip
column per coordinate so a reader can see which coordinates it is.

\subsection{Paired scoring on common random numbers}
\label{sec:fit-paired}

The incumbent is candidate zero of every generation and is played on exactly the
same deals as every other candidate. The paired per-deal margin over it is
therefore already available and costs nothing to compute, and the repair is
simply to score each candidate by that margin rather than by its absolute win
rate. The pairing now happens inside the objective rather than in
post-processing, so the selection, the elite set and the recorded score are all
the paired quantity.

The invariant that catches a mis-wired estimator is a policy scored against
itself. Under a correct paired estimator, the per-deal margin of a configuration
against its own incumbent is exactly \vsixPairedSelfDelta{} on every deal, so the
margin is exactly zero and the bootstrap interval has zero width --- not a small
number, and not an interval containing zero. This is asserted at startup. It is
worth stating why so weak-looking a test is the right one: every plausible wiring
error --- pairing against the wrong candidate, pairing across generations,
dividing by the wrong rotation count, resampling games instead of deals --- turns
that exact zero into something else, and none of them can be seen in a fit's
output otherwise.

\subsection{A minimax-regret objective}
\label{sec:fit-objective}

The soft minimum is retained as one option and is no longer the only one. The
optimiser now dispatches explicitly over $\{\texttt{softmin}, \texttt{min},
\texttt{mean}, \texttt{regret}, \texttt{minimaxregret}\}$, the chosen objective
is written into the run record rather than inferred from a compiled default, and
the fits of record use \texttt{minimaxregret}: for each opponent on the panel the
regret of a candidate is the best win rate any candidate in that generation
achieved against that opponent minus this candidate's, and the objective is the
negative of the worst such regret.

The reason for preferring regret to the minimum itself is that the panel cells
are not equally hard, so a minimum over raw win rates is dominated by whichever
opponent is hardest and carries almost no gradient from the rest. The reason for
preferring it to the soft minimum is that the soft minimum was never the object
it was named after. Two temperatures in one pipeline --- one compiled into the
tuner, one defaulted in the selection script, neither recorded --- is itself a
correction owed to the earlier record, and \S\ref{sec:corr-self} makes it.

\subsection{Commit-gate KPIs}
\label{sec:fit-gates}

No configuration in this study was committed on a win rate. The pathology
instrument runs as a gate with a non-zero exit status, as step one of
\filepath{engine/experiments_v06.sh} rather than as a report at the end, and the
identity controls of \S\ref{sec:engine-repairs} run in the same step. The
thresholds are \vsixGateN{} in number and the shipped configuration's measured
value on each is: dead asks \vsixGateDeadAsk{}\%, longest dead run
\vsixGateDeadRun{}, repeated asks \vsixGateRepeatAsk{}\%, wrong declarations
\vsixGateDeclWrong{}\%, asks inside a half-suit the actor's own team already owns
outright \vsixGateOwnLocked{}\%, and starved turns \vsixGateStarved{}\%.

A gate that has never failed has never been tested, so the battery runs a
negative control: a configuration that the gate must reject. Unbanning the
deliberate miss --- allowing asks the actor can prove will fail, with the
ownership features gated so that the zero-probability incentive cannot fire ---
scores \vsixDeadWin{}\% against \vfive{}, interval [\vsixDeadWinLo{},
\vsixDeadWinHi{}], which is higher than the shipped configuration and would be
selected by any procedure ordered on win rate. The same configuration plays
\vsixDeadPathDead{}\% of its asks dead, runs \vsixDeadPathLongest{} dead asks
without interruption at its worst, and reaches the action cap in
\vsixDeadPathLimit{}\% of games. The gate rejects it.
\S\ref{sec:search-deadask} treats what that configuration is actually doing;
the point here is narrower and is about ordering. A battery that collects
strength first and pathology afterwards selects this build, and the v0.5 study's
own ablation table contains two configurations of the same character.

\subsection{Schedule, selection, and the winner's curse}
\label{sec:fit-schedule}

Two fits of record were run. Run~A starts from \vfive{}'s vector and fits the
\numVfiveParams{} v0.5 coordinates over the panel \vsixFitAPanel{}, at population
\vsixFitAPop{}, \vsixFitADeals{} deals by rotations per cell, relative step
\vsixFitASigmaRel{}, for \vsixFitAGens{} generations from fitting seed
\vsixFitASeed{}. Run~C starts from the v0.6 base and fits the wider
\numVsixParams{}-coordinate vector over the panel \vsixFitCPanel{}, at population
\vsixFitCPop{}, \vsixFitCDeals{}, relative step \vsixFitCSigmaRel{}, for
\vsixFitCGens{} generations from fitting seed \vsixFitCSeed{}. Both are paired,
both use \texttt{minimaxregret}, and both write a header record as the first line
of the trace carrying the base specification, the objective, the temperature, the
pairing flag, the population and elite sizes, the generation count, the deals and
rotations per cell, the seed, the relative step and the full per-coordinate step
vector. That header is why \S\ref{app:params} can name the run behind the shipped
vector: the v0.5 study had to recover its own fitting temperature by regular
expression over a source file and a markdown document, because nothing in the
trace recorded it.

\paragraph{The winner's curse.} The per-generation best is a maximum over a
population evaluated on shared seeds and is biased upwards; on the v0.5 trace the
gap between the best and the incumbent was \vsixFitBestGap{}, against
\vsixFitOrderStat{} predicted by the order statistic of \vsixFitPop{} draws of
the observed noise alone. The optimiser therefore re-scores both the
distribution mean and the running best on a common guard bank and returns the
higher. The guard is retained and it is honestly weak: the guard bank is derived
from the same root seed as every generation bank, so it is disjoint from them but
not independent of the fitting process. Selection is guarded, not held out. The
held-out check is the evaluation protocol of \S\ref{sec:protocol}, and every
number in \S\ref{sec:results} comes from a bank the fit never touched.

\paragraph{Two acceptance tests.} The first is the self-comparison of
\S\ref{sec:fit-paired}: a policy scored against itself returns a paired margin of
exactly \vsixPairedSelfDelta{} with a zero-width interval. It is an invariant and
it passes by construction, which is the point --- it fails loudly the moment the
estimator is rewired.

The second is a recovery test on a known gap, and it is the one the previous
harness could not pass. A deliberately handicapped base was constructed by
zeroing the leading ask weight, the hit-probability coefficient, which is the
single most important coordinate in the vector; that configuration scores
\numFitSanityBase{}\% against \vfive{}. A \numFitSanityGens{}-generation fit from
that base with the repaired optimiser recovers to \numFitSanityAfter{}\% against
\vfive{} and to \numFitSanityVthree{}\% against \vthree{}. The v0.5 harness's own
\numFitTraceGensVfive{}-generation trace on the same objective moved its
incumbent at \vsixFitSlope{} per generation with $t = \vsixFitSlopeT{}$, which is
no movement at all; a recovery of this size in \numFitSanityGens{} generations is outside
anything that harness produced in \numFitTraceGensVfive{}. This test is not part of
\filepath{engine/experiments_v06.sh} and its command is given in
\S\ref{app:repro-commands} so that it can be re-run.

\paragraph{What the repaired optimiser was worth.} It is the only thing in this
study that produced a replicated effect. The mechanism trials of
\S\ref{sec:mechanisms} are null: single-coordinate and single-feature
perturbations of the v0.5 policy class are worth zero at the budget
\S\ref{sec:protocol} shows is required to see them. The joint refit is not, and
\S\ref{sec:results-h2h} reports it against the panel on two disjoint banks. The
plateau this line of work has been sitting on is in the mechanisms, not in the
parameters --- but the parameters were never actually searched, because the
searcher could not move.
%% --- end sections_v06/09-fitting ---
%% --- begin sections_v06/10-protocol ---
\section{Evaluation protocol}
\label{sec:protocol}

Fish has a single chance node --- the deal --- followed by a public transcript,
and that structure permits an unusually strong variance-reduction design. This
section fixes the experimental unit, the interval construction, the
fitting/evaluation separation, and the metrics of \S\ref{sec:results}.

It also fixes something the two preceding studies did not, and the omission cost
this one five mechanisms. A protocol section is normally read as bookkeeping. In
this study it is a result: \S\ref{sec:protocol-resolution} shows that the
budget most of the corpus's published ablation rows were run at cannot separate
an effect of the size those rows report from a seed draw, and that five separate
mechanisms cleared a nominal \vsixBootAlpha{}\% paired interval at that budget and returned
exactly zero when re-run larger on two disjoint banks. Every claim in this report
is paired and re-run on a second bank for that reason.

A reader should be able to tell, for every number in \S\ref{sec:results}, whether
it is paired or unpaired, which bank it came from, and whether that bank was used
in fitting. The earlier reports could not support that reading, which is why
several of their comparisons are corrected in \S\ref{sec:corrections}.

\subsection{Six-rotation duplicate blocks}
\label{sec:protocol-blocks}

Seats alternate between the two teams, so the cyclic group of six seat rotations
of a fixed deal forms a balanced block: three seat shifts by two team
orientations. The odd rotations exchange which team holds which hand-triple; the
even rotations permute hands within a team. Playing all \vsixRotations{}
rotations of every deal therefore has each policy hold every hand-triple of that
deal exactly once, so the intrinsic advantage of the deal cancels from the
comparison rather than merely being averaged down. A matchup reported over $D$
deals is $6D$ games organised into $D$ blocks of six. This is the duplicate-block
convention of competitive play \cite{acpc}, strengthened: the competition design
balances the team label, and the six-seat rotation additionally balances the
assignment of individual hands within a team.

One guarantee of that convention does not survive, and it is stated here rather
than in a footnote because it silently halves several denominators. In a
\emph{mirror} match --- both teams running the same policy specification --- the
two team orientations of a deal produce byte-identical games, so each game is
played twice and the effective sample is \vsixMirrorDuplication{}\% of the
nominal one. Rates are unaffected; denominators are not. Every table in
\S\ref{sec:results} states which convention its $n$ uses, and a mirror
\texttt{match} at two rotations returns exactly $50\%$ for that reason and not
because the policies are balanced.

\subsection{Deal-clustered bootstrap}
\label{sec:protocol-ci}

The six rotations of one deal are one correlated cluster: they share the same
card distribution and differ only in seat assignment. Resampling \emph{games}
would treat those six observations as independent and understate the variance.
Every interval in this report therefore resamples \emph{deals} --- a cluster
bootstrap of the per-deal win count for a single arm, and a paired bootstrap of
the per-deal difference for a comparison between two configurations played on the
same deals --- taking the $2.5\%$ and $97.5\%$ quantiles of \vsixBootReps{}
resamples, for \vsixBootAlpha{}\% intervals. Both are implemented in
\filepath{engine/src/arena.hpp}.

These are \emph{percentile} bootstrap intervals with the deal as the resampling
cluster, and they inherit the known coverage error of that construction. We keep
it for two reasons. Clustering by deal is the correction that matters most in
this design, because the within-deal correlation across the six rotations is
large and ignoring it understates the variance by a substantial factor, whereas
the percentile-versus-BCa distinction is second-order at these sample sizes. And
we have not implemented the bias-corrected and accelerated or bootstrap-$t$
corrections, so no claim of nominal coverage is made: intervals should be read as
approximate, and a difference whose interval sits close to zero should not be
treated as resolved. \S\ref{sec:protocol-resolution} is what that caution looks
like when it is applied rather than merely stated.

\subsection{Paired estimators}
\label{sec:protocol-paired}

The headline comparison is a paired difference on common random numbers, and
every ablation is paired against its own control. The pairing is what makes the
design competitive with the estimator-side variance reduction developed for
poker \cite{zinkevich-divat,burch-aivat} without needing an auxiliary value
function: in this game the deal is the whole chance node, so replaying it cancels
the dominant variance component directly. \S\ref{sec:fit-paired} gives the
measured justification, and the same estimator is used in fitting and in
evaluation so that a candidate is scored during the search by the quantity it
will be judged on afterwards.

One quantity in this report is not available paired. Throughput
(\S\ref{sec:results-throughput}) is a property of one binary on one machine
rather than of a matchup, so there is no common deal to pair on; it is reported
as a single-arm figure measured under load and is a lower bound for that reason.

\subsection{The resolution requirement}
\label{sec:protocol-resolution}

Five mechanisms in this study cleared a nominal \vsixBootAlpha{}\% paired
interval at \numResSmallBudgetLo{}--\numResSmallBudgetHi{} games per cell and
returned exactly zero at
\numResLargeBudget{} games per cell on two disjoint banks.
Table~\ref{tab:resolution} is the record.

\begin{table}[ht]
\centering
\caption{Five mechanisms that resolved at a small budget and did not replicate.
Sign convention is the corpus's: $\delta = \text{reference} - \text{variant}$, so
a \emph{negative} $\delta$ means the variant is better. Every entry is a paired
per-deal margin with a deal-clustered percentile interval. The large-budget
column gives two independent seed banks. Source:
\texttt{research/v06/notes/R12-mechanism-trials.md}, \S6.1.}
\label{tab:resolution}
\footnotesize
\begin{tabular}{@{}>{\raggedright\arraybackslash}p{0.235\linewidth}>{\raggedright\arraybackslash}p{0.265\linewidth}>{\raggedright\arraybackslash}p{0.42\linewidth}@{}}
\toprule
Mechanism & Small-budget paired $\delta$ & Large-budget paired $\delta$, two banks \\
\midrule
\texttt{wVoid} at $1.5$
  & \numResVoidSmall{} [\numResVoidSmallCI{}], $n{=}\numResVoidSmallN{}$
  & \numResVoidLargeA{} [\numResVoidLargeACI{}] and
    \numResVoidLargeB{} [\numResVoidLargeBCI{}], $n{=}\numResVoidLargeN{}$ \\
\addlinespace
\texttt{wVoid} at $3$
  & \numResVoidThreeSmall{} [\numResVoidThreeSmallCI{}], $n{=}\numResVoidThreeSmallN{}$
  & as above \\
\addlinespace
delete the perseveration bonus
  & \numResPerseverSmall{} [\numResPerseverSmallCI{}], $n{=}\numResPerseverSmallN{}$
  & \numResPerseverLargeA{} [\numResPerseverLargeACI{}] and
    \numResPerseverLargeB{} [\numResPerseverLargeBCI{}], $n{=}\numResPerseverLargeN{}$ \\
\addlinespace
\texttt{declareMargin} more negative
  & \numResDeclSmall{}~pp unpaired, $n{=}\numResDeclSmallN{}$
  & \numResDeclLarge{} [\numResDeclLargeCI{}] --- the reference is better \\
\addlinespace
delete the negative value-of-information term
  & \numResVoiSmall{} [\numResVoiSmallCI{}], $n{=}\numResVoiSmallN{}$
  & not separated \\
\bottomrule
\end{tabular}
\end{table}

The planning rule that follows from this is arithmetic, not a convention. For an
unpaired win-rate comparison near even the half-width of a nominal $95\%$
interval is
\begin{equation}
  1.96\sqrt{\frac{p(1-p)}{N}} \;\approx\; \frac{0.98}{\sqrt{N}}
  \quad\text{at } p \approx \tfrac12 ,
  \label{eq:halfwidth}
\end{equation}
that is, about $98/\sqrt{N}$ percentage points at $N$ games. At $N = \numResSmallBudgetLo{}$ that is a
half-width of nearly five points; at $N = \numResVoidLargeN{}$ it is about
three. An effect
of one to two points --- the size every mechanism in Table~\ref{tab:resolution}
appeared to have --- is therefore indistinguishable from a seed draw at the
budget most of this corpus's published ablation rows were run at. Pairing buys
back a large factor, but it buys it back from the deal variance and not from the
requirement to check that the result is not the bank.

Two rules follow and both are enforced in this report. Every claim is
\emph{paired}: the variant and its control are played on the same deals and the
interval is on the per-deal difference. And every claim is \emph{re-run on a
second, disjoint bank}, and reported as replicated only if it holds in sign and
size on both. Where a result resolves on one bank and not the other, that is
said, and the per-cell rows are given so a reader can see which half of the
comparison carries it.

The general lesson is uncomfortable and is stated plainly in
\S\ref{sec:limitations}: at the budget this project used before v0.6, a
mechanism that does nothing has a substantial chance of being published as a
mechanism that does something, and several were.

\subsection{Seed separation and what was held out}
\label{sec:protocol-seeds}

Fitting used base seeds \vsixFitASeed{} (run~A) and \vsixFitCSeed{} (run~C).
Every seed bank used from \S\ref{sec:results} onwards is disjoint from both by
construction, and the disjointness is enforced by a registry of the
\vsixReservedSeeds{} reserved seeds that the battery checks, rather than by
discipline. The evaluation banks are \texttt{90210}, \texttt{31337},
\texttt{515151}, \texttt{777001} and \texttt{424242} for the head-to-head;
\texttt{515253} for the per-style panel; \texttt{606060} for the frozen-policy
ablations; \texttt{717171} for calibration; \texttt{828282} for the rule
dialects; \texttt{313131} for the partner regimes; \texttt{90210} again for the
search ladder, the deviation-rate probe and the pre-gate audit; \texttt{31} for
the pathology KPIs and the tie and belief probes; and \texttt{515253} with
\texttt{6543210} for the exploitability probe's fitting and evaluation halves
respectively. \vsixHeldOutBanks{} of them are held out and \vsixFitBanks{} are
fitting banks. Banks carried forward from the v0.5 study are named in
\S\ref{app:repro-commands} so that the corresponding rows are directly
comparable.

The separation is narrower than ``held out'' usually implies, and the report
scopes its claims accordingly. Deals are held out: no deal used in any reported
experiment appears in any fitting bank. Opponent \emph{identities} largely are
not --- the fitting panels of \S\ref{sec:fit-schedule} include \vfive{},
\vthree{} and two scripted archetypes that also appear in the evaluation panel
--- so the head-to-head and per-style results measure generalisation to new deals
rather than to new strategic populations. The exceptions, which are held out in
both senses, are named in \S\ref{sec:results-panel}.

\subsection{Metrics}
\label{sec:protocol-metrics}

The primary metric is win rate, reported with its paired margin over the control.
Sets differential --- mean half-suits won out of nine --- accompanies it, because
a policy can win while converting a smaller fraction of its asks and this report
contains a matchup in which that happens. Ask hit rate and declaration accuracy
are diagnostic rather than primary: they describe the decision quality of one
component and a policy can trade either against the other.

The pathology KPIs of \S\ref{sec:fit-gates} are gates, not metrics. They do not
enter any comparison and no configuration is ranked on them; they decide only
whether a configuration is eligible to be measured at all.

One standing rule governs how the opponent panel is reported: \emph{never
headline an aggregate over it}. The headline is the worst case and the minimax
regret, with the full per-opponent profile printed alongside. A mean over a
scripted panel is a weighted average whose weights are the panel's composition,
and the panel's composition is a choice this study made. For the same reason
rating systems are deliberately not used as a headline here
\cite{coulom-whr,herbrich-trueskill,nelo,balduzzi-reeval}: a single scalar over a
fixed scripted panel is exactly the aggregate the rule forbids, and the
transitivity it assumes is not established for this population. That is stated
rather than the systems being omitted silently.
%% --- end sections_v06/10-protocol ---
%% --- begin sections_v06/11-results ---
\section{Results}
\label{sec:results}

Every number in this section comes from \texttt{engine/experiments\_v06.sh}, whose artifacts are
listed in Appendix~\ref{app:repro}. The battery is ordered so that the commit gate runs before any
strength measurement: a configuration that scores higher while failing the pathology KPIs is
rejected, and \S\ref{sec:protocol} explains why that ordering is not a formality.

\subsection{Controls}
\label{sec:results-verify}

The identity control passes: \vsix{} carrying v0.5's parameter vector with every mechanism switched
off produces a diagnostic transcript identical to \vfive{}'s under md5 (\vsixIdentity{}). The
verification suite passes with \vsixAuditViolations{} soundness violations in \vsixAuditChecks{}
checks (\vsixVerify{}); no agent's deduced knowledge ever excludes the truth, and the deal-sampling,
card-counting and block engines agree.

The block dynamic program's table-aliasing defect (\S\ref{sec:engine-repairs}) is repaired at the
query level: \vsixAliasQueryMismatch{} query mismatches, against \numAliasMismatch{} in
\numAliasChecks{} before the fix. Raw reads of the shared pool still differ
(\vsixAliasRawMismatch{}) and are documented rather than fixed, because the pool is what keeps the
memory footprint of six concurrent observers bounded.

The declaration pre-gate audit, which was dead code for \vfive{} because the option was parsed only
in the v0.4 branch, now runs: over \vsixGateOpportunities{} declaration opportunities and
\vsixGateRejected{} gate rejections it reports \vsixGateFalseNeg{} false negatives
(\vsixGateVerdict{}). \vfive{}'s pre-gates are exact, where v0.4's were not.

\subsection{The commit gate}
\label{sec:results-panel}

\begin{table}[t]\centering
\caption{Pathology KPIs in mirror play, 600 games per arm at seed 31. These gate every commit. The
v0.4 column is the failure mode v0.5 removed; the \vsix{} column must not reintroduce any of it.}
\label{tab:pathology}
\begin{tabular}{lrrr}
\toprule
KPI & \vfour{} & \vfive{} & \vsix{} \\
\midrule
provably dead asks (\% of asks)      & \vsixPathFourDead    & \vsixPathFiveDead    & \vsixPathSixDead \\
longest dead run                     & \vsixPathFourLongest & \vsixPathFiveLongest & \vsixPathSixLongest \\
games with a dead run $\ge 6$ (\%)   & \vsixPathFourRunSix    & \vsixPathFiveRunSix    & \vsixPathSixRunSix \\
exact repeat asks (\%)               & \vsixPathFourRepeatACT & \vsixPathFiveRepeatACT & \vsixPathSixRepeatACT \\
declarations wrong (\%)              & \vsixPathFourDeclwrong & \vsixPathFiveDeclwrong & \vsixPathSixDeclwrong \\
ask hit rate (\%)                    & \vsixPathFourHit     & \vsixPathFiveHit     & \vsixPathSixHit \\
events/game                          & \vsixPathFourEvents  & \vsixPathFiveEvents  & \vsixPathSixEvents \\
games killed by the action limit (\%)& \vsixPathFourLimit   & \vsixPathFiveLimit   & \vsixPathSixLimit \\
\bottomrule
\end{tabular}
\end{table}

\subsection{Head to head}
\label{sec:results-h2h}

Over the \vsixHeadFiveBanks{} dedicated held-out banks of E3, six rotations each and
\vsixHeadFiveN{} games in total, \vsix{} wins \vsixHeadFiveMean\% against \vfive{} with a
per-bank range of \vsixHeadFiveMin\% to \vsixHeadFiveMax\%, and \vsixHeadFourMean\% against
\vfour{} (\vsixHeadFourMin\% to \vsixHeadFourMax\%).

\textbf{Those five banks are not the whole evidence, and quoting them alone would be selection on
the experiment that agreed.} The battery contains two further \vsix{}-versus-\vfive{} cells at
default rules on further disjoint banks --- the default-dialect row of the rules experiment and the
\vfive{} column of the paired ablation panel --- and \vsix{} loses both. Pooling every such cell in
the battery gives \textbf{\vsixHeadPooled\% over \vsixHeadPooledN{} games,
\vsixHeadCellsAbove{} of \vsixHeadCells{} cells above parity, naive $z = \vsixHeadZ{}$} --- and
that $z$ treats the six rotations of a deal as independent, so it overstates the significance rather
than understating it.

\textbf{Head to head, \vsix{} and \vfive{} are not separated.} What is separated is the paired
comparison over the whole panel (\S\ref{sec:results-ablations}) and the worst case across
playstyles (\S\ref{sec:results-lbr}).

\begin{table}[t]\centering
\caption{Held-out head-to-head against \vfive{}, one row per seed bank, six rotations.}
\label{tab:h2h}
%% --- begin tables_v06/h2h_v05 ---
\begin{tabular}{rrrr}
\toprule
Win rate (\%) & 95\% cluster CI & Mean sets \vsix{} & Mean sets \vfive{} \\
\midrule
51.44 & [49.11, 53.78] & 4.560 & 4.440 \\
50.44 & [48.11, 52.78] & 4.526 & 4.474 \\
51.33 & [49.06, 53.67] & 4.541 & 4.459 \\
50.11 & [47.78, 52.44] & 4.523 & 4.477 \\
51.83 & [49.56, 54.11] & 4.546 & 4.454 \\
\bottomrule
\end{tabular}
%% --- end tables_v06/h2h_v05 ---
\end{table}

This is reported before anything else so that no reader has to hunt for it. What \vsix{} changes is
not here.

\subsection{Exploitability}
\label{sec:results-exploit}

The v0.5 study named the absence of any exploitability measurement as the single largest hole in its
own evaluation. It is measured here. A responder is fitted from the same policy family, with the
repaired optimiser and the frozen target as its entire opponent panel, and is then re-measured on a
fresh seed bank because the win rate it reached during fitting is a maximum over a population on
shared seeds and is upward biased.

\begin{table}[t]\centering
\caption{Local best-response probe. The response win rate is what a fitted exploiter achieved, so
\emph{lower is better} for the frozen policy. The \vfour{} row is the positive control: the v0.4
study published $51.19\%$ $[49.67, 52.72]$ for the same experiment.}
\label{tab:lbr}
%% --- begin tables_v06/lbr ---
\begin{tabular}{lrrr}
\toprule
Frozen target & Response win rate & 95\% CI & Games \\
\midrule
v04 & 50.69 & [49.06, 52.31] & 3,600 \\
v05 & 50.31 & [48.75, 51.89] & 3,600 \\
v06 & 48.36 & [46.75, 49.97] & 3,600 \\
\bottomrule
\end{tabular}
%% --- end tables_v06/lbr ---
\end{table}

The control reproduces: against \vfour{} the response reaches \vsixLbrFour\%
[\vsixLbrFourLo, \vsixLbrFourHi] over \vsixLbrFourN{} games, whose interval overlaps the published
figure. Against \vfive{} it reaches \vsixLbrFive\% [\vsixLbrFiveLo, \vsixLbrFiveHi]. Against
\vsix{} it reaches \textbf{\vsixLbrSix\%} [\vsixLbrSixLo, \vsixLbrSixHi] --- an interval that
excludes parity. The same optimiser, the same budget, the same policy family and the same fresh
evaluation bank fail to reach parity against \vsix{} and succeed against both of its predecessors.

The reading has to be careful, and the v0.4 findings document got it wrong in exactly this place
(\S\ref{sec:corr-vfour}). The number is a \emph{lower bound within the searched class}: a response
that fails to reach parity does not prove the target is hard to exploit, only that this search did
not find the exploit. What the three rows support is the \emph{comparison} --- identical class,
identical budget, identical protocol --- and on that comparison \vsix{} is the least exploitable of
the three. It is also the first exploitability number this project has for \vfive{}, whose own
study named its absence as the largest hole in its evaluation.

\subsection{The per-style profile and the worst case}
\label{sec:results-lbr}

\begin{table}[t]\centering
\caption{Per-opponent win rate over the full style set, pooled over \vsixPoolBanks{} disjoint seed
banks (515253, 90210, 424242), six rotations per cell. The headline is the \emph{minimum}, not the
mean. Minimax regret is measured against the best of the three arms on each opponent and is
therefore relative; \S\ref{sec:limitations} states what it is carried by.}
\label{tab:perstyle}
%% --- begin tables_v06/perstyle_pooled ---
\begin{tabular}{lrrrr}
\toprule
Opponent & \vsix{} & \vfive{} & \vfour{} & games \\
\midrule
v05 & 51.00 & 50.00 & 48.39 & 4,800 \\
v04 & 50.00 & 50.94 & 50.00 & 4,800 \\
v03 & 76.06 & 73.71 & 73.33 & 4,800 \\
v02 & 80.00 & 81.72 & 83.06 & 1,800 \\
lockout & 78.25 & 79.50 & 77.94 & 4,800 \\
detective & 76.71 & 75.88 & 77.28 & 4,800 \\
diversifier & 95.44 & 94.11 & 92.89 & 1,800 \\
hunter & 98.44 & 97.67 & 97.67 & 1,800 \\
bluffer & 99.83 & 99.89 & 99.94 & 1,800 \\
random & 100.00 & 100.00 & 100.00 & 1,800 \\
silent & 83.87 & 81.96 & 78.89 & 4,800 \\
feint & 53.37 & 53.44 & 51.11 & 4,800 \\
withholder & 79.15 & 70.54 & 67.50 & 4,800 \\
\midrule
\textbf{worst case} & \textbf{50.00} & \textbf{50.00} & \textbf{48.39} & \\
worst opponent & v04 & v05 & v05 & \\
mean (descriptive) & 78.63 & 77.64 & 76.77 & \\
\textbf{minimax regret} & \textbf{3.06} & \textbf{8.60} & \textbf{11.65} & \\
\bottomrule
\end{tabular}
%% --- end tables_v06/perstyle_pooled ---
\end{table}

Pooled over \vsixPoolBanks{} banks, \vsix{}'s worst cell is \vsixPoolWorstSix\% (against
\vsixPoolWorstOppSix{}) and \vfive{}'s is \vsixPoolWorstFive\% (its own mirror, which is
$50\%$ by construction), while minimax regret over the style set is \vsixPoolRegretSix{} against
\vsixPoolRegretFive{} and \vsixPoolRegretFour{}. \vsix{} is ahead of \vfive{} on \vsixPoolAhead{} of the thirteen styles, behind on
\vsixPoolBehind{} and equal on \vsixPoolTied{}, with a largest single-style loss of
\vsixPoolWorstLoss{} points.

The single large cell is the deception archetype the project owner brought here from live play: the
\textbf{withholder}, which holds cards of a half-suit it was asked for and then declines to ask
back in it. \vsix{} takes \vsixPoolWithholderSix\% against \vfive{}'s
\vsixPoolWithholderFive\%, a gain of \vsixPoolWithholderGain{} points over
\vsixPoolWithholderN{} games, and it replicates at all \vsixPoolBanks{} banks. It was in the
fitting panel, and \S\ref{sec:limitations} says what that means for how the number should be read.

\subsection{Belief as a predictor}
\label{sec:results-belief}

\begin{table}[t]\centering
\caption{Each posterior scored as a predictor on identical states: mean predictive log loss over the
unresolved cards, and the realised hit rate of that posterior's own argmax ask. The exact posterior
is the worst row.}
\label{tab:belief-results}
%% --- begin tables_v06/belief ---
\begin{tabular}{lrrrr}
\toprule
Posterior & Cards & Mean NLL & Argmax $p$ & Argmax hit \\
\midrule
exact (uniform prior) & 140,661 & 1.42246 & 0.4660 & 47.94 \\
sinkhorn th=0.000 ph=0.122 & 140,661 & 1.39083 & 0.4896 & 49.99 \\
sinkhorn th=0.200 ph=0.122 & 140,661 & 1.38663 & 0.4939 & 50.75 \\
sinkhorn th=0.300 ph=0.122 & 140,661 & 1.38465 & 0.4972 & 51.06 \\
sinkhorn th=0.445 ph=0.122 & 140,661 & 1.38218 & 0.5035 & 51.49 \\
sinkhorn th=0.600 ph=0.122 & 140,661 & 1.38055 & 0.5113 & 51.59 \\
sinkhorn th=0.800 ph=0.122 & 140,661 & 1.38055 & 0.5211 & 51.64 \\
sinkhorn th=1.100 ph=0.122 & 140,661 & 1.38311 & 0.5336 & 51.64 \\
sinkhorn th=1.500 ph=0.122 & 140,661 & 1.38758 & 0.5459 & 51.51 \\
sinkhorn th=2.000 ph=0.122 & 140,661 & 1.39221 & 0.5555 & 51.40 \\
\bottomrule
\end{tabular}
%% --- end tables_v06/belief ---
\end{table}

\subsection{The paired panel comparison}
\label{sec:results-paired}

The comparison the rest of this section rests on. Every variant plays the same deals against the
same panel; the statistic is the per-deal margin, bootstrapped resampling deals as clusters.

\begin{table}[t]\centering
\caption{\vsix{} against \vfive{} over a seven-opponent panel, 800--1{,}000 games per cell, on
\vsixPairedBanks{} disjoint seed banks. A positive margin means \vsix{} is ahead.}
\label{tab:pairedpanel}
%% --- begin tables_v06/pairedpanel ---
\begin{tabular}{lrr}
\toprule
Seed bank & \vsix{} $-$ \vfive{} & 95\% paired CI \\
\midrule
606060 & +2.69 & [+0.98, +4.35] \\
515253 & +2.67 & [+1.24, +4.11] \\
90210 & +1.13 & [-0.31, +2.56] \\
\bottomrule
\end{tabular}
%% --- end tables_v06/pairedpanel ---
\end{table}

All \vsixPairedBanks{} banks are positive, \vsixPairedExcl{} of them with an interval excluding
zero, mean \vsixPairedMean{} points and smallest \vsixPairedMin{}. This is the separated result;
the head-to-head of \S\ref{sec:results-h2h} is not.

\subsection{Paired ablations}
\label{sec:results-ablations}

\begin{table}[t]\centering
\caption{Paired mechanism ablations against \vsix{}: every variant plays the same deals against the
same panel, and the per-deal margin is bootstrapped resampling deals as clusters. A positive delta
means \vsix{} beats the variant.}
\label{tab:ablations}
%% --- begin tables_v06/ablations ---
\begin{tabular}{lrrr}
\toprule
Variant & Pooled & \vsix{} $-$ variant & 95\% paired CI \\
\midrule
\texttt{v05} & 65.73 & -2.69 & [-4.35, -0.98] \\
\texttt{v06:legacy=1} & 65.73 & -2.69 & [-4.35, -0.98] \\
\texttt{v06:xf=0} & 68.54 & +0.12 & [-1.67, +1.92] \\
\texttt{v06:s1=1,det=16,cand=6,kappa=2.0,maxq=26} & 69.77 & +1.35 & [-0.04, +2.79] \\
\bottomrule
\end{tabular}
%% --- end tables_v06/ablations ---
\end{table}

\subsection{The search ladder}
\label{sec:results-search}

Tables~\ref{tab:search} and \ref{tab:searchdev} in \S\ref{sec:search}. The result is that the
statistic, not the machinery, decides whether a determinized search helps, and that once the
statistic is corrected the search has nothing left to add.

\subsection{Partner regimes}
\label{sec:results-calib}

\begin{table}[t]\centering
\caption{Win rate against three \vfive{} opponents with the two remaining seats of the acting team
played by the named policy, 800 games per cell at seed 313131. The self-play row is not the
headline, and the three rows below it are why.}
\label{tab:partners}
%% --- begin tables_v06/partners ---
\begin{tabular}{lrr}
\toprule
Partners & \vsix{} & \vfive{} \\
\midrule
self & 52.25 & 50.00 \\
v03 & 34.50 & 33.12 \\
detective & 33.50 & 33.62 \\
withholder & 34.75 & 35.50 \\
\bottomrule
\end{tabular}
%% --- end tables_v06/partners ---
\end{table}

\vsix{}'s advantage over \vfive{} is $2.25$ points when its partners are copies of itself and
between $-0.8$ and $+1.4$ when they are not, with no mixed row separated from zero at this sample
size. A refit whose objective is scored in self-play buys a policy that coordinates with copies of
itself, and that coordination is exactly what a human partner does not supply. This is the project
owner's standing decision measured rather than asserted, and \S\ref{sec:limitations} carries it
forward as the sharpest limitation on everything in this section.

\subsection{Rule dialects}
\label{sec:results-dialects}

The result is reported in Appendix~\ref{app:dialect}: the ordering of \vsix{} and \vfive{} does not
change under the legacy dialect, with out-of-turn declaration disabled, or with cardless declaration
disabled.

\subsection{Throughput}
\label{sec:results-throughput}

\begin{table}[t]\centering
\caption{Games per second, self-play, all threads.}
\label{tab:throughput}
%% --- begin tables_v06/throughput ---
\begin{tabular}{lr}
\toprule
Policy & Games/s \\
\midrule
\texttt{v04} & 147.9 \\
\texttt{v05} & 276.3 \\
\texttt{v06} & 303.4 \\
\texttt{v05:topk=0} & 286.0 \\
\texttt{v05:belief=indep,topk=0} & 9289.0 \\
\bottomrule
\end{tabular}
%% --- end tables_v06/throughput ---
\end{table}
%% --- end sections_v06/11-results ---
%% --- begin sections_v06/12-corrections ---
\section{Corrections to the record}
\label{sec:corrections}

This project's convention is that corrections are stated as corrections, in their own section, and
not silently patched into the artifacts they correct. The v0.5 study opened a twelve-entry register
against the v0.4 record; four of its entries never reached print, and this study adds five of its
own. All are listed here with the measurement that forces them.

\subsection{Entries carried forward from the v0.5 register that were never published}
\label{sec:corr-vfour}
\label{sec:corr-dialect}

\begin{description}

\item[C7 --- ``an undeclared half-suit scores nothing''.] The v0.4 paper justifies its second
  forcing horizon on the ground that an unclaimed half-suit is worth zero. It is not, in this
  harness: residue at the action cap is adjudicated by physical majority, and on a wrong
  post-horizon declaration the declaring team holds a mean \numResidueHeld{} of the six cards. The
  v0.4 paper contradicts itself on this point --- its own rules section states the dialect
  correctly. Neither paper printed the correction. It is printed here, and it matters beyond
  bookkeeping: the argument it supports is the one v0.5 removed as mechanism M8, so the mechanism
  was right for a reason its own justification got wrong.

\item[C9 --- restarting the willingness ladder.] The v0.4 paper argues that restarting the
  forced-endgame ladder lets early, safer declarations resolve cards and sharpen the later ones.
  The premise is never satisfied: essentially every forced endgame in this dialect contains exactly
  \numLadderLiveSets{} live half-suit (\numLadderOccasions{} occasions at seed \numLadderSeed{}), so
  there are no later declarations to sharpen.

\item[The fitted adversary that ``fails to reach 50\%''.] The v0.4 findings document reports that a
  policy fitted specifically against v0.4 fails to reach parity with it. The same study's own local
  best-response table records the response achieving $51.19\%$ with a confidence interval of
  $[49.67, 52.72]$, and the paper states the null result correctly. The findings document turns a
  null result into a win. This is the single most consequential unretracted claim in the corpus,
  because it is the only place anything resembling an exploitability number has ever been reported.

\item[``Modelling declaration timing as optimal stopping is worth measurable win rate''.] The same
  document. The same study's own ablation puts the value-based stopping rule at $+0.12$ points with
  an interval of $[-1.23, +1.47]$ against a fixed threshold, and the v0.4 specification says in
  terms that the ablation does not resolve a benefit.

\end{description}

\subsection{Corrections this study adds}
\label{sec:corr-self}

\begin{description}

\item[The tie channel is not a defect, and it is not where it was said to be.] The v0.5 design
  register, and the recon that opened this study, priced the exact-tie channel as the largest
  available loss and located it in the target dimension. Both halves are wrong.
  \vsixTieSameSuitFive\% of ties are two \emph{cards} of one half-suit at one target, not one card
  at two targets; and the exact posterior separates \vsixTieExactSepPctFive\% of them, so no
  tie-break rule can beat enumeration order --- measured, all four rules realise the same hit rate
  to three significant figures. The associated ``+3.16 half-suits per game of headroom'' is a
  hindsight bound and prices clairvoyance. See \S\ref{sec:diag-irreducible}.

\item[Enumeration order decides a fifth of contested decisions, not more than half.] \vfive{}'s
  top-$K$ chain and threat pass re-scores the leaders and moves the pick at
  \vsixTieShipMovedFive\% of ties. The frequently-quoted figure conflates the tie rate with the
  rate at which array order is the deciding rule.

\item[``Exact inference costs six points'' is a statement about the prior, not about exactness.]
  Scored as predictors on identical states, the exact posterior under a uniform-deal prior is the
  \emph{worst} of the three inference paths measured, at \vsixBeliefExactNLL{} nats and
  \vsixBeliefExactHit\% argmax hit rate against the deployed approximation's
  \vsixBeliefShippedNLL{} and \vsixBeliefShippedHit\%. The policy prior is the entire difference.
  The matched-budget refit under the exact belief that three studies have listed as the outstanding
  experiment is therefore not worth running: the object it would fit is a less informative
  posterior. See \S\ref{sec:belief-exactworse}.

\item[The v0.5 fit resolved nothing, and the shipped vector is v0.4's.] \S\ref{sec:diag-fitnoise}.
  The consequence for the reader of the v0.5 paper is that its fitting section describes a procedure
  that did not, at the budget and step size used, move the policy; and the consequence for this
  study is that a working optimiser had to be built before any mechanism could be evaluated.

\item[The exact reference engine was unsound in multi-agent play.] The block dynamic program parked
  every instance's tables in a shared per-thread pool, so any second construction on the same thread
  silently repointed the first instance's tables: \numAliasMismatch{} mismatches in
  \numAliasChecks{} checks. Every previously published number computed with the exact belief in a
  configuration where more than one seat holds such an object is affected. The deployed
  approximate-belief numbers are not, because that path never queries the object --- which is
  exactly why the defect survived three studies. See \S\ref{sec:engine-repairs}.

\end{description}

\subsection{A note on comparability}
\label{sec:corr-comparability}

Two of the corrections above change how earlier numbers should be read rather than what they are,
and one changes the numbers themselves. Readers comparing across versions should note that this
study's evidence standard is not the corpus's: every claim here is paired, and every claim that
survived is re-run on a second, disjoint seed bank at a per-cell budget three to seven times larger
than the published ablation rows of the v0.3, v0.4 and v0.5 studies. \S\ref{sec:protocol-paired}
gives the measurements that forced that change, including five mechanisms in this study that cleared
a $95\%$ interval at the corpus's budget and returned zero at this one.
%% --- end sections_v06/12-corrections ---
%% --- begin sections_v06/13-limitations ---
\section{Threats to validity}
\label{sec:limitations}

\paragraph{\vsix{} is not separated from \vfive{} head to head, and this paper says so.}
Pooled over every head-to-head cell in the battery it is \vsixHeadPooled\% with a naive $z$ of
\vsixHeadZ{} (\S\ref{sec:results-h2h}). Its claim is on the axis this project has always said it
cares about --- the paired panel comparison and the worst case across playstyles --- together with a
set of corrections and negative results that close off most of the design register it inherited.

\paragraph{The worst-case improvement is carried by one column.} Minimax regret over the
thirteen-style set is \vsixPoolRegretSix{} for \vsix{} against \vsixPoolRegretFive{} for \vfive{}, but
delete the withholder column and the three arms converge. The withholder was in the fitting panel
of both fits. The correct reading is that the refit bought robustness against a deception archetype
it was shown, which is what a minimax-regret objective over a panel containing that archetype is
supposed to do, and not that \vsix{} is uniformly more robust.

\paragraph{The null results are null \emph{at this budget and against this panel}.}
Five mechanisms are reported as worth zero. Each was measured paired, at a thousand deals per cell,
on two disjoint banks, against a six- or seven-opponent panel. That resolves effects of roughly a
point; it does not resolve effects of a quarter of a point, and a policy improvement of a quarter of
a point per mechanism, compounded across a dozen mechanisms, would be a real policy improvement.
What this study establishes is that none of the mechanisms on the inherited register is large, not
that the sum of many small ones is zero.

\paragraph{The search verdict is conditional on two quantities.}
Test-time search is reported as not paying, and the conclusion rests on the variance of the terminal
return and the fidelity of the rollout blueprint. Both are movable. A shaped or truncated return
with a well-fitted leaf evaluator would cut the variance; a blueprint closer to the deployed policy
at the same cost would cut the bias. This study measured the frontier for the second
(Table~\ref{tab:rollout}) and did not build the first, because the leaf evaluator available is
algebraically close to a rescaling of the hit probability. A study that rebuilds the value function
first should re-open the search question; \S\ref{sec:search-ships} says so, and the machinery is
retained switched off for exactly that reason.

\paragraph{Exchangeability is established for the tied set, not for the whole candidate set.}
The finding of \S\ref{sec:diag-irreducible} is that candidates the policy's score cannot separate
are also candidates the \emph{exact posterior} cannot separate. It is not a claim that the exact
posterior is uninformative in general --- it disagrees with the deployed approximation about the
best ask at \vsixDisagreeFive\% of decisions --- only that where the deployed score ties, the
information ties too.

\paragraph{\vsix{}'s advantage is a self-play advantage and does not transfer to a change of
partner.} This is the sharpest limitation in the paper and it is measured, not conjectured. Against
three \vfive{} opponents, \vsix{} beats \vfive{} by $2.25$ points when its two partners are copies
of itself, by $1.4$ when they are \vthree{}, by $-0.1$ when they are the detective and by $-0.8$
when they are the withholder (Table~\ref{tab:partners}; 800 games a cell, none of the three mixed
rows separated from zero). A refit whose objective is scored in self-play buys a policy that
coordinates with copies of itself, and that coordination is exactly what a human partner does not
supply. The project owner's standing decision --- never headline the self-play configuration as
though it were the one a human will meet --- is why this table exists, and it is why the paper does
not claim \vsix{} is a better partner for a person.

\paragraph{The opponent panel is a modelling choice and the deception archetypes are hand-built.}
The withholder, the silent and the feint are stylised implementations of manoeuvres the project
owner reported from live play; they are not humans. A gain against the withholder is evidence about
robustness to a class of behaviour, not a measurement of play against a person. The partner-regime
columns have the same limitation from the other side: a ``human-model partner'' is a weaker or
deceptive bot, not a human.

\paragraph{Exploitability is a lower bound within a searched class.}
The best-response probe fits a responder from the same policy family with the same optimiser. A
stronger response outside that family would give a larger number. A response that fails to reach
parity therefore proves nothing, which is precisely the reading error \S\ref{sec:corr-vfour}
corrects in the v0.4 record.

\paragraph{The engine correction changes previously published numbers.}
Any earlier result computed with the exact block belief in a configuration where more than one seat
held such an object was affected by the table-aliasing defect. This study re-runs what it quotes;
it does not re-run the whole of the v0.4 and v0.5 studies, and it does not claim to know the size of
the effect on their published tables.

\paragraph{Fitting seeds are disjoint from evaluation banks, but the panel is not.}
The refit's panel contains opponents that also appear in the evaluation panel. The head-to-head and
per-style results are therefore in-panel for those cells and held-out only in the seed sense. The
cells that are held out in both senses --- the styles absent from the fitting panel --- are marked
as such in Table~\ref{tab:perstyle}, and they are the ones a sceptical reader should weigh.

\paragraph{Termination remains empirical.}
As in v0.5, no configuration shipped here hits the action cap in any measured run, but that is a
measurement rather than a theorem. The structural backstop is M1, whose gate is a hard deduction and
cannot be wrong; what is not proved is that M1 plus the stopping rule terminates on every reachable
position.
%% --- end sections_v06/13-limitations ---
%% --- begin sections_v06/14-conclusion ---
\section{Conclusion}
\label{sec:conclusion}

Three FishBot studies in a row have ended with a policy that is not measurably stronger than the one
before it. This one explains why.

The design register the project had accumulated --- a blind target dimension, a large tie channel, an
exact posterior held back only by its cost, a two-ply search that could be deepened --- described a
policy with several points of accessible headroom. Measured at a resolution the project had not
previously used, none of it is there. The ties are exchangeable and therefore irreducible. The exact
posterior is a worse predictor than the approximation, because exactness is bought by discarding the
policy prior. The search is dominated by the variance of its own return, to the tune of \vsixCurseGap{} points,
and once that is corrected it has nothing left to say. Four further mechanisms, each keyed to a
measured channel and each fitted rather than hand-set, are worth zero.

What was actually broken was the apparatus. The optimiser could not move a policy: one step size for
thirty-four coordinates spanning two and a half orders of magnitude, an unpaired objective, and an
objective documented as a worst case that was a weighted mean. The evaluation could not resolve a
point: five mechanisms in this study passed a 95\% interval at the budget the corpus uses and failed
at three times it. And the exact inference engine, the one component nobody doubted, was returning
another observer's tables whenever two observers held one.

\vsix{} fixes those, and the fixed optimiser then finds what forty generations of the old one could
not: a policy that is ahead of \vfive{} on the paired panel comparison at all three banks it was
measured on, and that gains nine points against the deception archetype the project owner brought
here from live play. Head to head it is a null --- \vsixHeadPooled\% over \vsixHeadPooledN{} games
--- and the paper says so in the abstract, because an earlier draft of this work quoted only the
five banks that agreed and the correction is exactly the kind this study exists to make. That is a
smaller headline than a new mechanism would have been. It is also the only one the evidence
supports, and stating it that way is the point.

Three things the next study should take from this. The tie-break question is closed --- any candidate
signal must first be shown to separate the tied set under the exact posterior, and none of the
obvious ones can. The search question is open but conditional: it should be re-opened only after the
leaf evaluator is rebuilt, because the present one is algebraically close to a rescaling of the hit
probability and cannot support a depth-limited search. And the evidence standard is now the binding
constraint on everything else: at a thousand deals per cell, paired, on two banks, this policy class
looks flat, and a study that wants to find a quarter-point mechanism will need a per-decision
objective rather than a per-game one.
%% --- end sections_v06/14-conclusion ---

%% --- begin sections_v06/bibliography ---
%% ---------------------------------------------------------------
%% Seeded from research/v06/notes/R7-literature-search.md section 11 (the
%% BibTeX block, 30 entries), converted to the manual \bibitem style the v0.4
%% and v0.5 papers use.  There is no .bib file anywhere in this repository and
%% no bibtex/biber step in any build command; that is deliberate.
%%
%% EVERY one of R7's 30 entries is present below.  None was dropped.  Eight of
%% them duplicate an entry that already existed in paper/sections/bibliography.tex
%% under a different key (nayyar-common, lerer-sparta, foerster-bad, hu-obl,
%% brown-rebel, schmid-sog, lisy-lbr, tian-jps); in each case R7's key is the one
%% kept, the older key is NOT defined here, and any \cite in sections_v06/ uses
%% the R7 spelling.  Do not reintroduce the old keys -- two \bibitem's for one
%% work is how a bibliography starts disagreeing with itself.
%%
%% R7's section 9 lists items it explicitly does NOT recommend for v0.6
%% (ReBeL/SoG/DeepStack implementations, the team belief DAG, learned belief
%% models, neural agents, LLM agents, hidden-role machinery, gadget-game
%% refinement).  "Not recommended as method" is not "not applicable as related
%% work": those entries are retained here and are cited in
%% \S\ref{sec:related}/\S\ref{app:related} as context, with
%% \S\ref{app:rel-declined} stating why each is declined.  Nothing is dropped
%% silently.
%%
%% R7's "Unfetchable" table lists four items that support no recommendation.
%% Three of them are gaps in an entry that IS present (Edelkamp's runtime,
%% Valet's team-game coverage, the PDF of the adversarial-team subgame paper),
%% and are recorded in the \bibitem notes below rather than dropped.  The fourth,
%% arXiv:2512.15435 (an outer-learning framework for Skat), was surfaced in
%% search only and never fetched, so there is no entry for it: DROPPED, reason
%% -- unverified metadata, cited by nothing.
%%
%% The remaining entries carry forward from paper/sections/bibliography.tex and
%% paper/sections_v05/bibliography.tex, unchanged, for the parts of
%% \S\ref{sec:related} that are unchanged since v0.4.
%% ---------------------------------------------------------------

\begin{thebibliography}{99}

%% ---------------------------------------------------------------
%% Fish and Literature
%% ---------------------------------------------------------------

\bibitem{fishbot03}
FishLab Research Project, ``FishBot v0.3: A Count-Conditioned, Empirically Tuned
Policy for Six-Player Canadian Fish,'' project technical report, 2026. Project
repository: \url{https://github.com/dylann29/fish-optimization}

\bibitem{fishbot04}
D. Nguyen, FishLab Research Project, ``FishBot v0.4: Observer-Conditioned Deal
Inference and Value-Based Declaration in Six-Player Canadian Fish,'' project
technical report, 2026. Project repository:
\url{https://github.com/dylann29/fish-optimization}

\bibitem{fishbot05}
D. Nguyen, FishLab Research Project, ``FishBot v0.5: A Self-Inflicted Deadlock,
Its Measurement, and the Repair of the Ask Policy That Caused It,'' project
technical report, 2026. Project repository:
\url{https://github.com/dylann29/fish-optimization}

\bibitem{pagat}
J. McLeod, ``Literature,'' \emph{Pagat}, 2006--2025, last updated 1 July 2026.\\
\url{https://www.pagat.com/quartet/literature.html}

\bibitem{develin}
M. Develin, ``Canadian Fish,'' chapter 9 of \emph{Card Games}, author's online
manual.\\
\url{https://web.archive.org/web/20170720075756/http://bantha.org/~develin/cardgames.html\#ch9}

\bibitem{dorsa}
Deposit Genius, ``Literature Game Strategy (Canadian Fish),'' online article.\\
\url{https://depositgenius.com/literature-strategy-canadian-fish/}

\bibitem{strategy-site}
Canadian Fish Project, ``Strategy notes,'' 2026.\\
\url{https://canadian-fish.vercel.app/strategy}

\bibitem{wikiliterature}
Wikipedia, ``Literature (card game),'' online encyclopedia entry.\\
\url{https://en.wikipedia.org/wiki/Literature_(card_game)}

\bibitem{somani}
N. Somani, ``literature: card game implementation and learning bot,'' GitHub
repository, MIT licence, Python, 2018--.\\
\url{https://github.com/neelsomani/literature}

\bibitem{somani-server}
N. Somani, ``literature-server,'' GitHub repository, with a live deployment at
\url{https://literature.neelsomani.com/}\\
\url{https://github.com/neelsomani/literature-server}

\bibitem{playstore}
``Literature: Fish Card Game'' (\texttt{com.cards.game.literature}), Google Play
store listing, closed source.\\
\url{https://play.google.com/store/apps/details?id=com.cards.game.literature}

\bibitem{rlcard}
D. Zha, K.-H. Lai, S. Huang, Y. Cao, K. Reddy, J. Vargas, R. Wei, J. Guo, and
X. Hu, ``RLCard: A Toolkit for Reinforcement Learning in Card Games,''
arXiv:1910.04376, 2019.

\bibitem{goadrich2026valet}
M. Goadrich, A. Morenville, and \'E. Piette, ``Valet: A Standardized Testbed of
Traditional Imperfect-Information Card Games,'' arXiv:2603.03252, 2026.
21 games in the RECYCLE description language; whether any of them is a team game
is not stated on the abstract page and the full text was not fetched.\\
\url{https://arxiv.org/abs/2603.03252}

%% ---------------------------------------------------------------
%% Determinization and search under imperfect information
%% ---------------------------------------------------------------

\bibitem{levy}
D. Levy, ``The Million Pound Bridge Program,'' \emph{Heuristic Programming in
Artificial Intelligence: The First Computer Olympiad}, Ellis Horwood, 1989.

\bibitem{ginsberg-gib}
M. L. Ginsberg, ``GIB: Imperfect Information in a Computationally Challenging
Game,'' \emph{Journal of Artificial Intelligence Research} 14:303--358, 2001.

\bibitem{frank-basin-ai}
I. Frank and D. Basin, ``Search in Games with Incomplete Information: A Case
Study Using Bridge Card Play,'' \emph{Artificial Intelligence}
100(1--2):87--123, 1998.

\bibitem{frank-basin-aaai}
I. Frank, D. Basin, and H. Matsubara, ``Finding Optimal Strategies for Imperfect
Information Games,'' \emph{AAAI-98}, pp.~500--507, 1998.

\bibitem{long-pimc}
J. Long, N. R. Sturtevant, M. Buro, and T. Furtak, ``Understanding the Success
of Perfect Information Monte Carlo Sampling in Game Tree Search,''
\emph{AAAI-10}, pp.~134--140, 2010.

\bibitem{cowling-ismcts}
P. I. Cowling, E. J. Powley, and D. Whitehouse, ``Information Set Monte Carlo
Tree Search,'' \emph{IEEE Transactions on Computational Intelligence and AI in
Games} 4(2):120--143, 2012.

\bibitem{goodman-ris}
J. Goodman, ``Re-determinizing Information Set Monte Carlo Tree Search in
Hanabi,'' arXiv:1902.06075, 2019.

\bibitem{cazenave2019alphamu}
T. Cazenave and V. Ventos, ``The $\alpha\mu$ Search Algorithm for the Game of
Bridge,'' arXiv:1911.07960, 2019.\\
\url{https://arxiv.org/abs/1911.07960}

\bibitem{edelkamp2021paranoia}
S. Edelkamp, ``Knowledge-Based Paranoia Search in Trick-Taking,''
arXiv:2104.05423, 2021. Skat; over 1,000 points on the extended Seeger scale.
The runtime is not stated on the abstract page, so no compute cost is attributed
to this method in this report.\\
\url{https://arxiv.org/abs/2104.05423}

\bibitem{buro-kermit}
M. Buro, J. R. Long, T. Furtak, and N. R. Sturtevant, ``Improving State
Evaluation, Inference, and Search in Trick-Based Card Games,'' \emph{IJCAI-09},
pp.~1407--1413, 2009.

\bibitem{solinas-supervised}
C. Solinas, D. Rebstock, and M. Buro, ``Improving Search with Supervised
Learning in Trick-Based Card Games,'' \emph{AAAI-19}, 2019; arXiv:1903.09604.

\bibitem{rebstock-inference}
D. Rebstock, C. Solinas, M. Buro, and N. R. Sturtevant, ``Policy Based Inference
in Trick-Taking Card Games,'' \emph{IEEE Conference on Games 2019}, 2019;
arXiv:1905.10911.

\bibitem{solinas-history}
C. Solinas, D. Rebstock, N. R. Sturtevant, and M. Buro, ``History Filtering in
Imperfect Information Games: Algorithms and Complexity,'' \emph{NeurIPS 2023},
2023; arXiv:2311.14651.

%% ---------------------------------------------------------------
%% Depth-limited solving, subgame solving, and decision-time planning
%% ---------------------------------------------------------------

\bibitem{brown2018depthlimited}
N. Brown, T. Sandholm, and B. Amos, ``Depth-Limited Solving for
Imperfect-Information Games,'' \emph{NeurIPS 2018}, 2018; arXiv:1805.08195.
Master-level heads-up no-limit hold'em on a 4-core CPU and 16 GB, via
multi-valued states. The soundness argument is for two-player zero-sum games and
is not inherited by the team game studied here.\\
\url{https://arxiv.org/abs/1805.08195}

\bibitem{zhang2021subgamesolvingwithoutck}
B. H. Zhang and T. Sandholm, ``Subgame Solving without Common Knowledge,''
\emph{NeurIPS 2021}, 2021; arXiv:2106.06068.\\
\url{https://proceedings.neurips.cc/paper/2021/hash/c96c08f8bb7960e11a1239352a479053-Abstract.html}

\bibitem{zhang2022teamsubgame}
B. H. Zhang, L. Carminati, F. Cacciamani, G. Farina, P. Olivieri, N. Gatti, and
T. Sandholm, ``Subgame Solving in Adversarial Team Games,'' \emph{NeurIPS 2022},
2022. Cited from the proceedings abstract page; the mirrored PDF did not
decode.\\
\url{https://proceedings.neurips.cc/paper_files/paper/2022/hash/aa5f5e6eb6f613ec412f1d948dfa21a5-Abstract-Conference.html}

\bibitem{sokota2024updateequivalence}
S. Sokota, G. Farina, D. J. Wu, H. Hu, K. A. Wang, J. Z. Kolter, and N. Brown,
``The Update-Equivalence Framework for Decision-Time Planning,'' \emph{ICLR
2024}, 2024; arXiv:2304.13138. Mirror-descent search at two orders of magnitude
less search time than public-belief-state search in Hanabi.\\
\url{https://openreview.net/forum?id=JXGph215fL}

\bibitem{lerer2020sparta}
A. Lerer, H. Hu, J. Foerster, and N. Brown, ``Improving Policies via Search in
Cooperative Partially Observable Games,'' \emph{AAAI 2020}, 2020;
arXiv:1912.02318. SPARTA; Hanabi 24.08 to 24.61 at test time.\\
\url{https://ojs.aaai.org/index.php/AAAI/article/view/6208}

\bibitem{hu2021lbs}
H. Hu, D. J. Wu, A. Lerer, J. Foerster, and N. Brown, ``Learned Belief Search:
Efficiently Improving Policies in Partially Observable Settings,''
arXiv:2106.09086, 2021.\\
\url{https://arxiv.org/abs/2106.09086}

\bibitem{fickinger2021rlfinetuning}
A. Fickinger, H. Hu, B. Amos, S. Russell, and N. Brown, ``Scalable Online
Planning via Reinforcement Learning Fine-Tuning,'' \emph{NeurIPS 2021}, 2021;
arXiv:2109.15316.\\
\url{https://arxiv.org/abs/2109.15316}

\bibitem{milec2021cdr}
D. Milec, O. Kub\'i\v{c}ek, and V. Lis\'y, ``Continual Depth-limited Responses
for Computing Counter-strategies in Sequential Games,'' arXiv:2112.12594, 2021.
Robust-response methods do not compose with limited-look-ahead solving off the
shelf.\\
\url{https://arxiv.org/abs/2112.12594}

\bibitem{milec2025abd}
D. Milec, V. Kova\v{r}\'ik, and V. Lis\'y, ``Adapting Beyond the Depth Limit:
Counter Strategies in Large Imperfect Information Games,'' arXiv:2501.10464,
2025.\\
\url{https://arxiv.org/abs/2501.10464}

\bibitem{kubicek2025lamir}
O. Kub\'i\v{c}ek and V. Lis\'y, ``Look-ahead Reasoning with a Learned Model in
Imperfect Information Games,'' arXiv:2510.05048, 2025; ICLR 2026 per the arXiv
PDF header.\\
\url{https://arxiv.org/abs/2510.05048}

\bibitem{kubicek2026refinements}
O. Kub\'i\v{c}ek, V. Lis\'y, and T. Sandholm, ``Equilibrium Refinements Improve
Subgame Solving in Imperfect-Information Games,'' arXiv:2601.17131, 2026.\\
\url{https://arxiv.org/abs/2601.17131}

%% ---------------------------------------------------------------
%% Full solvers, and the evaluation of their approximation quality
%% ---------------------------------------------------------------

\bibitem{brown2020rebel}
N. Brown, A. Bakhtin, A. Lerer, and Q. Gong, ``Combining Deep Reinforcement
Learning and Search for Imperfect-Information Games,'' \emph{NeurIPS 2020},
2020; arXiv:2007.13544.\\
\url{https://arxiv.org/abs/2007.13544}

\bibitem{schmid2023sog}
M. Schmid, M. Morav\v{c}\'ik, N. Burch, R. Kadlec, J. Davidson, K. Waugh,
N. Bard, F. Timbers, M. Lanctot, G. Z. Holland, E. Davoodi, A. Christianson, and
M. Bowling, ``Student of Games: A Unified Learning Algorithm for Both Perfect
and Imperfect Information Games,'' \emph{Science Advances} 9(46), 2023;
arXiv:2112.03178. Circulated 2021--2022 as ``Player of Games.''

\bibitem{moravcik2017deepstack}
M. Morav\v{c}\'ik, M. Schmid, N. Burch, V. Lis\'y, D. Morrill, N. Bard,
T. Davis, K. Waugh, M. Johanson, and M. Bowling, ``DeepStack: Expert-Level
Artificial Intelligence in Heads-Up No-Limit Poker,'' \emph{Science}
356(6337):508--513, 2017.

\bibitem{lisy2017lbr}
V. Lis\'y and M. Bowling, ``Equilibrium Approximation Quality of Current
No-Limit Poker Bots,'' arXiv:1612.07547, 2017. Local Best Response; every ACPC
2016 entrant is more than 3,180 mBB/h exploitable.\\
\url{https://arxiv.org/abs/1612.07547}

%% ---------------------------------------------------------------
%% Cooperative play, conventions, and belief modelling
%% ---------------------------------------------------------------

\bibitem{bard-hanabi}
N. Bard, J. N. Foerster, S. Chandar, N. Burch, M. Lanctot, H. F. Song,
E. Parisotto, V. Dumoulin, S. Moitra, E. Hughes, I. Dunning, S. Mourad,
H. Larochelle, M. G. Bellemare, and M. Bowling, ``The Hanabi Challenge: A New
Frontier for AI Research,'' \emph{Artificial Intelligence} 280, 2020;
arXiv:1902.00506.

\bibitem{foerster2019bad}
J. N. Foerster, H. F. Song, E. Hughes, N. Burch, I. Dunning, S. Whiteson,
M. Botvinick, and M. Bowling, ``Bayesian Action Decoder for Deep Multi-Agent
Reinforcement Learning,'' \emph{ICML 2019}, 2019; arXiv:1811.01458.\\
\url{https://arxiv.org/abs/1811.01458}

\bibitem{hu-sad}
H. Hu and J. N. Foerster, ``Simplified Action Decoder for Deep Multi-Agent
Reinforcement Learning,'' \emph{ICLR 2020}, 2020; arXiv:1912.02288.

\bibitem{hu-otherplay}
H. Hu, A. Lerer, A. Peysakhovich, and J. N. Foerster, ```Other-Play' for
Zero-Shot Coordination,'' \emph{ICML 2020}, 2020; arXiv:2003.02979.

\bibitem{hu2021obl}
H. Hu, A. Lerer, B. Cui, L. Pineda, N. Brown, and J. Foerster, ``Off-Belief
Learning,'' \emph{ICML 2021}, 2021; arXiv:2103.04000.\\
\url{https://arxiv.org/abs/2103.04000}

\bibitem{sokota2021capi}
S. Sokota, E. Lockhart, F. Timbers, E. Davoodi, R. D'Orazio, N. Burch,
M. Schmid, M. Bowling, and M. Lanctot, ``Solving Common-Payoff Games with
Approximate Policy Iteration,'' \emph{AAAI 2021}, 2021; arXiv:2101.04237.\\
\url{https://arxiv.org/abs/2101.04237}

\bibitem{hu2022sed}
H. Hu, S. Sokota, D. J. Wu, A. Bakhtin, A. Lupu, B. Cui, and J. N. Foerster,
``Self-Explaining Deviations for Coordination,'' \emph{NeurIPS 2022}, 2022;
arXiv:2207.12322. IMPROVISED; the first algorithmic Hanabi finesse plays.\\
\url{https://proceedings.neurips.cc/paper_files/paper/2022/hash/faa6276ea12d7afeb3e42b210c86f688-Abstract-Conference.html}

\bibitem{cox-hanabi}
C. Cox, J. De Silva, P. DeOrsey, F. H. J. Kenter, T. Retter, and J. Tobin, ``How
to Make the Perfect Fireworks Display: Two Strategies for Hanabi,''
\emph{Mathematics Magazine} 88(5):323--336, 2015.

\bibitem{tian2020jps}
Y. Tian, Q. Gong, and T. Jiang, ``Joint Policy Search for Multi-Agent
Collaboration with Imperfect Information,'' \emph{NeurIPS 2020}, 2020;
arXiv:2008.06495. Contract Bridge: $+0.63$ IMPs/board against WBridge5.\\
\url{https://arxiv.org/abs/2008.06495}

%% ---------------------------------------------------------------
%% Human-compatible play and regularised search
%% ---------------------------------------------------------------

\bibitem{jacob2022pikl}
A. P. Jacob, D. J. Wu, G. Farina, A. Lerer, H. Hu, A. Bakhtin, J. Andreas, and
N. Brown, ``Modeling Strong and Human-Like Gameplay with KL-Regularized
Search,'' \emph{ICML 2022}, 2022; arXiv:2112.07544.\\
\url{https://arxiv.org/abs/2112.07544}

\bibitem{hu2022humanregularized}
H. Hu, D. J. Wu, A. Lerer, J. Foerster, and N. Brown, ``Human-AI Coordination
via Human-Regularized Search and Learning,'' arXiv:2210.05125, 2022.\\
\url{https://arxiv.org/abs/2210.05125}

\bibitem{fair2022cicero}
Meta Fundamental AI Research Diplomacy Team (FAIR), A. Bakhtin, N. Brown,
E. Dinan, G. Farina, C. Flaherty, D. Fried, A. Goff, J. Gray, H. Hu, and others,
``Human-Level Play in the Game of Diplomacy by Combining Language Models with
Strategic Reasoning,'' \emph{Science} 378(6624):1067--1074, 2022.\\
\url{https://www.science.org/doi/10.1126/science.ade9097}

\bibitem{dizdarevic2025ah2ac2}
T. Dizdarevi\'c and others, ``Ad-Hoc Human-AI Coordination Challenge,''
\emph{ICML 2025}, 2025; arXiv:2506.21490.\\
\url{https://proceedings.mlr.press/v267/dizdarevic25a.html}

%% ---------------------------------------------------------------
%% Team games, hidden roles, and common information
%% ---------------------------------------------------------------

\bibitem{nayyar2013common}
A. Nayyar, A. Mahajan, and D. Teneketzis, ``Decentralized Stochastic Control
with Partial History Sharing: A Common Information Approach,'' \emph{IEEE
Transactions on Automatic Control} 58(7):1644--1658, 2013; arXiv:1209.1695.\\
\url{https://arxiv.org/abs/1209.1695}

\bibitem{celli-gatti}
A. Celli and N. Gatti, ``Computational Results for Extensive-Form Adversarial
Team Games,'' \emph{AAAI 2018}, 2018; arXiv:1711.06930.

\bibitem{farina-collusion}
G. Farina, A. Celli, N. Gatti, and T. Sandholm, ``Connecting Optimal Ex-Ante
Collusion in Teams to Extensive-Form Correlation: Faster Algorithms and Positive
Complexity Results,'' \emph{ICML 2021}, PMLR 139:3164--3173, 2021.

\bibitem{zhang-tbdag}
B. H. Zhang, G. Farina, A. Celli, and T. Sandholm, ``Team Belief DAG:
Generalizing the Sequence Form to Team Games for Fast Computation of Correlated
Team Max-Min Equilibria via Regret Minimization,'' \emph{ICML 2023}, 2023;
arXiv:2202.00789.

\bibitem{carminati-tpi}
L. Carminati, F. Cacciamani, M. Ciccone, and N. Gatti, ``A Marriage between
Adversarial Team Games and 2-Player Games: Enabling Abstractions, No-Regret
Learning and Subgame Solving,'' \emph{ICML 2022}, 2022; arXiv:2206.09161.

\bibitem{carminati2024hiddenrole}
L. Carminati, B. H. Zhang, G. Farina, N. Gatti, and T. Sandholm, ``Hidden-Role
Games: Equilibrium Concepts and Computation,'' \emph{ACM Conference on Economics
and Computation (EC) 2024}, 2024; arXiv:2308.16017.\\
\url{https://arxiv.org/abs/2308.16017}

%% ---------------------------------------------------------------
%% Counting, permanents and constrained sampling
%% ---------------------------------------------------------------

\bibitem{ryser}
H. J. Ryser, \emph{Combinatorial Mathematics}, Carus Mathematical Monographs 14,
Mathematical Association of America, 1963.

\bibitem{glynn}
D. G. Glynn, ``The Permanent of a Square Matrix,'' \emph{European Journal of
Combinatorics} 31(7):1887--1891, 2010.

\bibitem{jsv}
M. Jerrum, A. Sinclair, and E. Vigoda, ``A Polynomial-Time Approximation
Algorithm for the Permanent of a Matrix with Non-Negative Entries,''
\emph{Journal of the ACM} 51(4):671--697, 2004.

\bibitem{lsw}
N. Linial, A. Samorodnitsky, and A. Wigderson, ``A Deterministic Strongly
Polynomial Algorithm for Matrix Scaling and Approximate Permanents,''
\emph{Combinatorica} 20(4):545--568, 2000.

\bibitem{chen-sis}
Y. Chen, P. Diaconis, S. P. Holmes, and J. S. Liu, ``Sequential Monte Carlo
Methods for Statistical Analysis of Tables,'' \emph{Journal of the American
Statistical Association} 100(469):109--120, 2005.

%% ---------------------------------------------------------------
%% Evaluation methodology
%% ---------------------------------------------------------------

\bibitem{acpc}
N. Bard, J. Hawkin, J. Rubin, and M. Zinkevich, ``The Annual Computer Poker
Competition,'' \emph{AI Magazine} 34(2):112--118, 2013; and the competition
rules describing duplicate matches and common seeds.

\bibitem{zinkevich-divat}
M. Zinkevich, M. Bowling, N. Bard, M. Kan, and D. Billings, ``Optimal Unbiased
Estimators for Evaluating Agent Performance,'' \emph{AAAI-06}, pp.~573--578,
2006.

\bibitem{burch-aivat}
N. Burch, M. Schmid, M. Morav\v{c}\'ik, D. Morrill, and M. Bowling, ``AIVAT: A
New Variance Reduction Technique for Agent Evaluation in Imperfect Information
Games,'' \emph{AAAI-18}, 2018; arXiv:1612.06915.

\bibitem{waugh-abstraction}
K. Waugh, D. Schnizlein, M. Bowling, and D. Szafron, ``Abstraction Pathologies
in Extensive Games,'' \emph{AAMAS 2009}, pp.~781--788, 2009.

\bibitem{balduzzi-reeval}
D. Balduzzi, K. Tuyls, J. P\'erolat, and T. Graepel, ``Re-evaluating
Evaluation,'' \emph{NeurIPS 2018}, 2018; arXiv:1806.02643.

\bibitem{coulom-whr}
R. Coulom, ``Whole-History Rating: A Bayesian Rating System for Players of
Time-Varying Strength,'' \emph{Computers and Games 2008}, 2008.

\bibitem{herbrich-trueskill}
R. Herbrich, T. Minka, and T. Graepel, ``TrueSkill\texttrademark: A Bayesian
Skill Rating System,'' \emph{Advances in Neural Information Processing Systems}
19, 2006.

\bibitem{nelo}
M. Van den Bergh, ``Comments on Normalized Elo,'' Fishtest technical note; and
Stockfish contributors, ``Statistical Methods and Algorithms in Fishtest.''\\
\url{https://cantate.be/Fishtest/normalized_elo_practical.pdf}

\bibitem{kelidari2026goldstandard}
N. Kelidari, M. Haghi, and M. Salmani, ``A Gold-Standard Study of What Makes a
Lightweight Game-Playing Agent Strong,'' arXiv:2607.06854, 2026. Gin Rummy and
Leduc; concludes that the binding constraint is information, not capacity.\\
\url{https://arxiv.org/abs/2607.06854}

\end{thebibliography}
%% --- end sections_v06/bibliography ---

\appendix
%% --- begin sections_v06/A-dialect ---
\section{The rules dialect in full}
\label{app:dialect}

The dialect is unchanged from the v0.4 and v0.5 studies, so this appendix is a
summary and a delta rather than a re-derivation. Table~\ref{tab:dialect} is the
index: it lists every consequential choice against the engine switch that
changes it. The v0.4 report's Appendix A states each choice at the length
required to re-implement it, including the ask-legality predicate, the deal and
shuffle, declaration resolution, the arbitration orders and the scoring rule
\cite{fishbot04}; that treatment is accurate except on the two clauses of
\S\ref{app:dialect-corrections}. The two files that carry the definitions are
\filepath{engine/src/fish.hpp} (deck, indexing, the \texttt{Rules} record, the
legality predicate) and \filepath{engine/src/game.hpp} (the driver: declaration
rounds, turn handling, the forced endgame, the action cap).

\begin{table}[ht]
\centering
\caption{The rules dialect used for every experiment in this report unless
stated otherwise, and the engine switch that changes each choice. The switches
are command-line flags of the \texttt{fish} driver; the underlying fields are
members of \texttt{Rules} in \texttt{engine/src/fish.hpp}. Rows marked
``modelling choice'' are not fixed by any published statement of the rules.
Unchanged from the v0.4 and v0.5 studies.}
\label{tab:dialect}
\small
\begin{tabular}{@{}p{0.29\linewidth}p{0.46\linewidth}p{0.19\linewidth}@{}}
\toprule
Choice & Value used & Switch \\
\midrule
Deck composition & 54 cards in nine half-suits (low and high band of each suit,
plus the four eights and two jokers); nine cards per player &
\texttt{-{}-sets=8} gives the 48-card, eight-half-suit form \\
Ask legality & Live half-suit; asker holds another card of it and not the named
card; target is an opponent and still holds cards & --- \\
Turn transfer & A successful ask retains the turn; a miss passes it to the
target & --- \\
Out-of-turn declarations & Permitted at any moment, including during an
opponent's turn & \texttt{outOfTurnDeclare} (\texttt{-{}-no-out-of-turn}) \\
Cardless players may declare & Permitted &
\texttt{cardlessMayDeclare} (\texttt{-{}-no-cardless-declare}) \\
Incorrect declaration & Awards the half-suit to the opposing team & --- \\
Declaration aftermath & The six cards leave play and the turn returns to the
player who held it & --- \\
Cardless turn-holder & Chooses which live teammate receives the turn, from its
own belief (a modelling choice, and more restrictive than the rules of record;
\S\ref{sec:rules-dialect}) & --- \\
Forced endgame (whole team cardless) & The live team declares every remaining
half-suit, sharing only willingness, over an eight-level ladder
$\{0.995, 0.98, 0.95, 0.90, 0.80, 0.65, 0.50, \text{best guess}\}$ &
\texttt{-{}-legacy} gives the two-level ladder $\{0.38, \text{best guess}\}$ \\
Simultaneous declarations & Lowest seat index declares first (a modelling
choice; cost measured in \S\ref{sec:rules-dialect}) &
\texttt{declArbitration} (\texttt{-{}-arb=low\textbar high\textbar turn}) \\
Action cap & 400 asks (360 under the v0.3 dialect), any residue adjudicated
neutrally by majority physical holding, ties to the holder of the lowest card &
\texttt{-{}-maxasks} \\
Conventions & No prearranged signalling conventions; team communication is
limited to the willingness bit described above & --- \\
\bottomrule
\end{tabular}
\end{table}

\subsection{Clause table and what changed}
\label{app:dialect-table}

Nothing in Table~\ref{tab:dialect} changed between \vfour{}, \vfive{} and
\vsix{}. The engine repairs of \S\ref{sec:engine-repairs} are inference and
bookkeeping defects, not rule changes, and none of them alters a game's outcome
given the same actions. Two entries in the table are worth reading against the
v0.4 appendix rather than in place of it. The cardless-turn-holder row is
narrower than the rules of record permit, and the row for simultaneous
declarations is a choice the rules do not make at all; both are quantified in
\S\ref{sec:rules-dialect}, and both measure at approximately zero, which is why
neither is a \vsix{} mechanism.

Dialect robustness is reported rather than assumed: \S\ref{sec:results-dialects}
replays the primary head-to-head result under \vsixDialectN{} perturbations of
this table over \vsixDialectDeals{} deals at seed \vsixDialectSeed{}, spanning
\vsixDialectSpan{}\%.

\subsection{Clauses the \vsix{} mechanisms depend on}
\label{app:dialect-depends}

Four clauses carry weight in this report, in the sense that a mechanism of
\S\ref{sec:mechanisms} or a measurement of \S\ref{sec:belief} would be different
under a different choice. They are stated here in full so that the dependence is
auditable.

\paragraph{Residue is adjudicated by physical majority.} The driver counts asks
and stops the game at \texttt{maxAsks}, 400 under the dialect studied. Awarding
the unresolved residue to either side would bias the outcome towards whichever
policy stalls more readily, so each remaining half-suit is instead awarded to
the team physically holding the majority of its six cards, ties going to the
team holding the lowest-indexed card of that half-suit. Two consequences matter.
An unclaimed half-suit is worth its majority holding rather than nothing, which
is the premise \S\ref{app:dialect-corrections} corrects; and the cap is a
backstop against unbounded play rather than part of the rules, so termination is
an empirical property of the fitted policies and is reported as one, at
\vsixLimitRate{}\% incidence.

\paragraph{The cardless-player rule.} A seat with no cards leaves the asking
cycle: it cannot ask, and it is an illegal target, because ask legality requires
the target to hold cards. It remains in the game in two ways. It can be named in
a teammate's declaration as the holder of a card --- impossible while it is
genuinely empty, so its emptiness constrains the allocations a teammate may
name, which is a hard constraint the exact count law of
\S\ref{sec:belief-count} consumes. And, with \texttt{cardlessMayDeclare}
enabled, it may itself declare. The turn can reach a cardless seat only
immediately after a declaration removed its last cards, and that seat then
chooses which live teammate receives it.

\paragraph{An opponent void is permanent.} Cards move in exactly one direction:
from a target to an asker, and only on a successful ask. Asking in half-suit $h$
requires the asker to hold another card of $h$. A seat holding no card of $h$
therefore can neither ask in $h$ nor receive a card of $h$, and no sequence of
legal actions restores it. Voids are monotone, they are public, and a void
created against an opponent removes that opponent's right to ask in the
half-suit for the rest of the game. This is a rules-level fact rather than a
tactical one, and it is what the offensive-void term of \S\ref{sec:mech-cheap}
prices.

\paragraph{The forced endgame.} When one team holds no cards at all, no legal
ask remains for either side and the live team must declare every remaining
half-suit. Teammates may negotiate openly but may exchange nothing except
willingness to declare, so the driver implements the selection as a threshold
sweep over the eight-level ladder of Table~\ref{tab:dialect} rather than as a
comparison of confidences: for each threshold in descending order the driver
scans the live team's seats over the remaining half-suits and asks each whether
it is willing to declare at confidence at least that threshold; the first seat
that says yes declares, the ladder restarts, and the sweep repeats until every
half-suit is resolved. Declarations here resolve exactly as elsewhere, including
the possibility of misdeclaring and handing a half-suit to the cardless team,
and are recorded separately in the event history. If every seat refuses at every
level, the remainder is awarded by physical majority as above.

\subsection{Two clauses corrected}
\label{app:dialect-corrections}

The v0.4 appendix is accurate on every clause but two. Both errors are stated in
that report and repeated in the source; neither was retracted by the v0.5 paper,
which measured them and did not print the result. They are restated correctly
here so that a reader who reaches this appendix first is not misled, and
\S\ref{sec:corr-dialect} gives them as corrections with the file and line of the
text being corrected.

\paragraph{An undeclared half-suit does not score nothing.} The v0.4 report's
declaration section justifies the unconditional second stage of its forcing rule
--- declare the best candidate whatever its estimated probability, past a second
horizon --- on the ground that an undeclared half-suit scores nothing. Inside
this harness it does not. The action cap calls the residue adjudicator described
above, which awards each unresolved half-suit by physical majority, and on a
wrong post-horizon declaration the declaring team holds a mean
\numResidueHeld{} of the six cards, which is the majority, which is what
adjudication would have given it. The stage does not convert nothing into a coin
flip; it converts a probable award into a certain concession. The premise is
true only of a dialect that awards nothing at the cap, and the v0.4 study
changed the harness away from that dialect after its parameter vector was
frozen, without updating the justification. The v0.4 rules section states the
dialect correctly, so that report contradicts itself within two sections.

\paragraph{Restarting the ladder sharpens nothing in observed play.} The v0.4
appendix states that a consequence of restarting the willingness ladder is that
early, safer declarations resolve cards and thereby sharpen the allocations
available for the later ones. The reasoning is correct and the situation does
not arise. Over every occasion on which a team went cardless with live
half-suits remaining --- \numLadderOccasions{} of them at seed
\numLadderSeed{} --- the number of live half-suits at that moment was
\numLadderLiveSets{} in every case, so there is nothing to order and every
ordering counterfactual is identical to the non-ordering one by construction.
The structure is not impossible: a cardless team requires only that the other
team hold a multiple of six cards. It does not occur because these policies cash
locked half-suits promptly. A policy made more patient about declaring would
begin to produce multi-half-suit forced endgames, at which point the restart
would start doing the work the v0.4 appendix credits it with, and the machinery
should be built alongside the patience change rather than after it.
%% --- end sections_v06/A-dialect ---
%% --- begin sections_v06/B-inference ---
\section{Inference: derivations and implementation notes}
\label{app:inference}

This appendix supplies the derivations \S\ref{sec:belief} states without proof
and the implementation detail a reimplementer needs. It also carries the
exactness audit in full, because the headline result of \S\ref{sec:belief} is a
comparison against an object claimed to be exact, and a claim of that kind is
worth exactly what its audit is worth.

The source files are \filepath{engine/src/belief.hpp} (constraint bookkeeping and
the approximate path), \filepath{engine/src/blockdp.hpp} (the exact count law),
and \filepath{engine/src/oracle.hpp} (the brute-force enumerator). Notation is
that of \S\ref{sec:belief-state}.

\subsection{The constraint system}
\label{app:inf-constraints}

The observer's state is a \texttt{Knowledge} object: a resolved-owner array, a
per-card surviving-support bitmask $M_c$, the public hand counts, per-seat and
per-half-suit ask and miss tallies, and a vector of live ask-legality
certificates. (C1)--(C3) are maintained incrementally as bit operations on the
support masks, and a support that collapses to a single seat resolves the card,
which in turn removes any certificate the resolution satisfies.

(C4) is derived rather than stored. Capacities are recomputed from the public
hand counts by \eqref{eq:capacity} after every event, and a seat with no spare
capacity is removed from the support of every unresolved card, which can cascade.
This is the only place the two constraint families interact, and it is why a
declaration --- correct or incorrect --- needs no special handling anywhere in
the system. A declaration removes six cards from play and changes the public
counts of the seats that held them; the six cards leave $U$ with their holders
known, the half-suit leaves $\mathcal{H}$, and each $h_p$ in
\eqref{eq:capacity} drops by the number of those cards that seat held. The same
identity that absorbs a successful ask absorbs a declaration. An incorrect
declaration is not a special case either: it reveals the true holders of all six
cards, which is strictly more information than a correct one gives the opposing
team, and it arrives through the same channel.

(C5) is stored as a list. Each certificate is a pair --- the asking seat and the
bitmask $D$ of \eqref{eq:disj} --- and a certificate whose card set collapses to
a single card resolves that card immediately rather than being stored. Two
properties of the store matter downstream. First, every certificate is confined
to one half-suit, which is what \S\ref{app:inf-blocks} exploits. Second, the
store admits duplicates: a repeated ask in the same half-suit under the same
support produces an identical certificate, and \texttt{onEvent} appended it
without checking. The de-duplication repair of \S\ref{sec:engine-repairs} is
opt-in for a reason that belongs in this appendix rather than in the main text:
the approximate conditioning step of \S\ref{app:inf-sinkhorn} applies each stored
certificate once, so two copies of one certificate condition twice and the
duplicate is not inert. De-duplicating is therefore a change of posterior, not a
change of representation, and it is off by default so that the \vfour{} and
\vfive{} numbers this report compares against are unchanged.

\subsection{The capacity-vector count law}
\label{app:inf-count}

\paragraph{The layer identity.} Order the unresolved cards $c_1,\dots,c_{|U|}$
and let $q$ be the capacity vector. A partial assignment of the first $j$ cards
leaves a vector $n$ of unused capacities, and because every one of those cards
consumes exactly one unit of capacity from exactly one seat,
\[
  \sum_p n_p \;=\; \sum_p q_p - j \;=\; |U| - j ,
\]
so $j$ is recoverable from $n$ and the layer index need not be stored. That is
what collapses the natural $|U|$-layer table into two single arrays. Each is laid
out as one array of size $\prod_p (q_p+1)$, addressed by the mixed-radix offset
$\mathrm{off}(n) = \sum_p n_p s_p$ with strides $s_p = \prod_{r<p}(q_r+1)$, and a
companion index sorts the reachable states by $\sum_p n_p$ so that a sweep visits
one layer at a time. Count vectors are packed one byte per seat into a single
\verb|uint64_t|, which makes the dominance test $n \ge t$ that guards a block
transition one branch-free word operation.

\paragraph{The bound.} The observer's own cards are resolved by (C1), so $q_i =
0$ and the table is a product over the five other seats. Also $|U| \le 45$: at
the deal the 54 cards lie in nine live half-suits, the observer holds nine and
those are resolved, and thereafter a card leaves $U$ when its holder is revealed
or its half-suit is declared and never re-enters. Subject to $\sum_{p\ne i} q_p =
|U|$, the product $\prod_{p \ne i}(q_p+1)$ is maximised at equal capacities, by
the arithmetic--geometric mean inequality, so
\begin{equation}
  \prod_{p\ne i}(q_p+1) \;\le\; \Bigl(\frac{\sum_{p \ne i}(q_p+1)}{5}\Bigr)^{5}
  \;=\; \Bigl(\frac{|U|}{5}+1\Bigr)^{5} \;\le\; 10^{5}.
  \label{eq:bound}
\end{equation}
Each state has at most six outgoing transitions, so a full forward--backward pass
costs at most $6 \times 10^5$ multiply--adds regardless of the position and the
two arrays together occupy under two megabytes. The implementation refuses to
build above a fixed state cap and reports the failure to the caller rather than
degrading silently.

\paragraph{The per-card marginal.} The marginal is the contraction of the two
tables at card $j$,
\begin{equation}
  \mu_{c_j, p}
    \;=\; \frac{\mathbf{1}[\,p \in M_{c_j}\,]}{Z}
      \sum_{\substack{n \,:\, |n| = |U| - j + 1 \\ n_p \ge 1}}
        F_{j-1}(n)\; B_j(n - e_p),
  \label{eq:marg}
\end{equation}
which is exact: every consistent assignment is counted once, by splitting it at
card $j$ into a prefix counted by $F_{j-1}$ and a suffix counted by $B_j$. Both
tables hold integer-valued counts in double precision, so the arithmetic is exact
while the counts stay below $2^{53}$, which \eqref{eq:bound} guarantees on every
reachable state.

\paragraph{The direct sampler.} The backward table doubles as a sampler. Having
placed $c_1,\dots,c_{j-1}$ and arrived at remaining capacities $n$, card $c_j$
goes to seat $p$ with probability proportional to $B_j(n - e_p)$, restricted to
$p \in M_{c_j}$ with $n_p \ge 1$. A seat receives positive weight only when the
suffix can still be completed, so every prefix is extendable and no draw is
discarded: the procedure is rejection-free for (C1)--(C4) and uniform over the
assignments satisfying them, because the weights are path counts. When a live
(C5) certificate is present the sampler is unchanged and the completed assignment
is tested against the literal disjunction \eqref{eq:disj}, drawing again on
failure, so the combined procedure is exact but not rejection-free. The sampler
is off the deployed decision path and is used as ground truth in the
cross-checks of \S\ref{app:inf-audit}.

\subsection{Half-suit block enumeration}
\label{app:inf-blocks}

$U$ is partitioned by half-suit and each group is classified.

\begin{description}[style=nextline]
  \item[Card group.] No live certificate in that half-suit. The group's cards are
        expanded one at a time by \eqref{eq:fwd}--\eqref{eq:bwd}, branching factor
        $|M_c| \le 5$.
  \item[Block group.] At least one live certificate. The group's assignments are
        enumerated exhaustively by depth-first search over the supports --- at
        most $5^6 = 15{,}625$ candidates for a six-card half-suit --- each tested
        against every certificate of that half-suit, and the survivors aggregated
        by count vector into $g_H(t)$ of \eqref{eq:blocktable}. Alongside
        $g_H(t)$ the engine accumulates the per-card breakdown
        $g_H^{(c\to p)}(t)$ that \eqref{eq:q1} needs, so one enumeration serves
        all three queries.
\end{description}

A group of either kind consumes a known number of cards, so the layer identity
survives and the outer recursion runs over groups: $F_b(n) = \sum_t g_b(t)
F_{b-1}(n+t)$, with $g_b$ the indicator of a single card placement in the
card-group case. Nothing in this treatment is approximate: the enumeration is
exhaustive, the certificate test is the literal disjunction, and the outer
recursion is a counting argument. Inclusion--exclusion over $k$ certificates
would instead cost $2^k$ passes over the whole of $U$, which is the cost the
half-suit locality of (C5) avoids.

\begin{corollary}[Equal probability within a count vector]
\label{cor:equalprob}
Under the uniform posterior of \S\ref{sec:belief-state}, any two allocations of
the same half-suit that satisfy (C1)--(C5) and share a count vector have equal
probability.
\end{corollary}

\begin{proof}
By \eqref{eq:q2} the probability of a surviving allocation $A$ depends on $A$ only
through $t(A)$, since $S_t$ is a function of the count vector and of the tables
outside the block.
\end{proof}

The maximum-probability allocation is therefore decided by the count vector
alone, and the choice among survivors within a count vector is arbitrary. The
caveat is the operative half: an allocation that violates a certificate has
probability zero \emph{even when its count vector is shared by survivors}, so
survival must be checked explicitly and cannot be inferred from the count vector.
The count-vector table records only survivors, which makes the corollary usable
and makes the caveat necessary at the same time. \S\ref{app:inf-latent} records
which of the two consumers of this table actually performs the check.

\subsection{The approximate conditioning step}
\label{app:inf-sinkhorn}

The approximate path maintains a matrix $p_{c,p}$ over unresolved cards and
seats, initialised to the prior weights \eqref{eq:prior} on the support $M_c$ and
zero off it. One Sinkhorn iteration normalises each row to sum to one, then
scales each column $p$ by $q_p / \sum_c p_{c,p}$, enforcing \eqref{eq:capacity}
in expectation. The deployed setting is four outer sweeps of eight such
iterations, with a conditioning step for the certificates applied between
consecutive sweeps and a final row normalisation.

The conditioning step derives as follows. Fix a certificate with asking seat $a$
and card set $D$, and let $E$ be the event $\bigvee_{d\in D}[\sigma(d)=a]$.
Treating $\{p_{d,a}\}_{d \in D}$ as independent Bernoulli indicators,
\[
  \Pr[\neg E] \;=\; \prod_{e \in D} (1 - p_{e,a}),
  \qquad
  \Pr[E] \;=\; 1 - \prod_{e \in D} (1 - p_{e,a}) .
\]
For $d \in D$ the event $\sigma(d) = a$ implies $E$, so $\Pr[\sigma(d) = a \mid
E] = p_{d,a}/\Pr[E]$. For a seat $p \ne a$, the event $\sigma(d) = p$ is
compatible with $E$ only if some other card of $D$ goes to $a$, whose independent
probability is $1 - \prod_{e \in D\setminus\{d\}}(1 - p_{e,a})$. Together,
\begin{equation}
  \Pr[\sigma(d) = a \mid E] = \frac{p_{d,a}}{1 - \prod_{e \in D}(1 - p_{e,a})},
  \qquad
  \Pr[\sigma(d) = p \mid E] = p_{d,p}\,
    \frac{1 - \prod_{e \in D \setminus \{d\}}(1 - p_{e,a})}{1 - \prod_{e \in D}(1 - p_{e,a})} .
  \label{eq:cond}
\end{equation}
Rows are renormalised afterwards and the next sweep restores the capacity
constraints the reweighting disturbed. Two degenerate cases are guarded: when
$\Pr[E]$ underflows the mass is forced onto the certificate rather than dividing
by a near-zero denominator, and when $\Pr[\neg E]$ underflows the certificate is
already implied and the step is skipped.

The independence assumption in \eqref{eq:cond} is false, because the capacity
constraints couple the cards of $D$, so interleaving it with capacity fitting does
not converge to the exact marginals of \eqref{eq:q1}. That is the sense in which
the deployed path is approximate, and it is the only sense: the bookkeeping of
(C1)--(C4) it runs on is exact.

\paragraph{The deployed declaration estimator.} The declaration query
\eqref{eq:q2} asks for a joint probability and the approximate path has only
marginals, so it evaluates the joint by sequential conditioning. For a candidate
allocation assigning cards $c_1,\dots,c_6$ to seats $p_1,\dots,p_6$,
\[
  \hat\alpha(A) \;=\; \prod_{j=1}^{6} \hat\mu^{(j)}_{c_j, p_j},
\]
where $\hat\mu^{(j)}$ is recomputed on the state in which $c_1,\dots,c_{j-1}$
have already been resolved to $p_1,\dots,p_{j-1}$. Each step fixes one card,
decrements that seat's capacity, propagates the support reductions and refits. A
card already resolved contributes one if it matches and zero otherwise, a card
whose support excludes its assigned seat returns zero immediately, and the
product short-circuits once it falls below a floor, since candidates that weak
are rejected regardless. The chain rule is exact --- the joint is the product of
successive conditionals in any fixed order --- so the estimator is exact in
structure and approximate only through the marginal solver. A product of
\emph{unconditioned} marginals would not be: it ignores the capacity coupling
that makes \eqref{eq:q2} necessary in the first place.

\subsection{The policy prior, derived}
\label{app:inf-prior}

Equation \eqref{eq:prior} is a likelihood correction expressed as prior weights.
The full Bayesian posterior over deals weights each deal $x$ by $L(\mathcal{I}
\mid x) = \prod_t \pi_{p_t}(a_t \mid \mathcal{I}_{<t}, x)$, the probability that
the observed actions would have been chosen given $x$. That object is not
available: it requires a model of every opponent's policy and a product over the
whole history. What \eqref{eq:prior} does instead is retain the one sufficient
statistic of the history that any plausible policy makes informative --- how many
of a seat's public asks fell in a given half-suit, against how many fell
elsewhere --- and expose its two coefficients to the same optimiser that fits the
policy. $\theta$ rewards a seat's asks in the card's own half-suit and $\phi$
penalises its asks elsewhere.

The weights enter as the initialisation of the Sinkhorn fit rather than as a
post-hoc multiplication, which is what keeps the capacity constraints exact:
$\theta = \phi = 0$ recovers the uniform initialisation and hence the
policy-agnostic posterior. It also means the prior has no route into the exact
path. \texttt{BlockDP::build} takes the \texttt{Knowledge} object alone; there is
no argument through which a coefficient could reach it, and the exact path never
calls the weight function at all. Any belief-mode ablation in this corpus is
therefore confounded with a prior ablation, and that confound is not a defect of
the experiment but a property of the architecture.

\paragraph{The exact/approximate predictive comparison in full.} The comparison
of \S\ref{sec:belief-exactworse} is run by \texttt{fish v6probe --mode=belief}
and its design is as follows. At every ask decision of the measured team, on the
acting seat's own \texttt{Knowledge} object, each candidate posterior is computed
on the \emph{same} state: the exact count law of \S\ref{sec:belief-count}, and
the approximate path at a list of $(\theta,\phi)$ settings that includes $\theta =
\phi = 0$ and the shipped values. Two quantities are then accumulated per
posterior.

\begin{enumerate}
  \item \textbf{Predictive log loss.} For every unresolved card of the state, the
        true holder is read from the ground-truth deal and $-\log \mu_{c,
        \text{truth}}$ is accumulated, floored to avoid an infinite contribution
        from a posterior that assigns exact zero. The reported figure is the mean
        over cards. The unit is one card, not one decision, so a state with a
        large unresolved pool contributes proportionally more --- which is the
        right weighting, because those are the states inference is for.
  \item \textbf{Argmax hit rate.} The legal asks are enumerated at that state and
        the candidates the acting seat's own knowledge already proves dead are
        discarded. Among the remainder the posterior's own maximiser of
        $\mu_{c,t}$ is selected, and the accumulator records whether the target
        really holds the card. The unit is one decision, and each posterior
        chooses its own ask, so the measure is the quality of the posterior's own
        recommendation rather than of a shared one.
\end{enumerate}

Two properties of this design carry the argument. It is played on identical
states, so no posterior is scored on a trajectory it induced. And it involves no
fitted ask weights on either side --- the argmax is over the posterior itself,
not over the linear utility of \S\ref{sec:mechanisms} --- so nothing in the
comparison depends on which posterior the policy was tuned against. That is what
makes it a verdict on the posteriors rather than a substitution result, and it is
why \S\ref{sec:belief-exactworse} can close a question that the v0.4 study's
belief-substitution ablation could not.

The result is Table~\ref{tab:belief}. The exact count law under a uniform deal
prior scores \vsixBeliefExactNLL{} nats per card and \vsixBeliefExactHit{}\%
argmax hit rate; the approximate path with the prior deleted scores
\vsixBeliefNoPriorNLL{} and \vsixBeliefNoPriorHit{}\%; the shipped prior scores
\vsixBeliefShippedNLL{} and \vsixBeliefShippedHit{}\%. The exact object is last
of the three on both measures. The approximate path beats the exact one even at
$\theta = \phi = 0$, where the two encode the same rules-only information; the
remaining gap to the shipped row is the prior alone, worth
\numThetaPriorHitGain{} points of hit rate and \numThetaPriorNatGain{} nats.

The interpretation is the one \S\ref{sec:belief-exactworse} gives and is repeated
here because the appendix is where a sceptical reader arrives: exactness is
relative to a prior, the uniform deal prior is a modelling assumption and not a
rule, and on this evidence it is a worse assumption than a crude count of who has
asked where.

\subsection{The exactness audit and its coverage}
\label{app:inf-audit}

Agreement between the count law and the block construction is weak evidence,
because they share a derivation. The audit of record is therefore a brute-force
enumerator that factorises nothing: for one observer's constraint state the
unresolved cards are searched by depth-first assignment to seats, candidates are
pruned against the surviving support (carrying (C1)--(C3)) and against the
capacity (carrying (C4)), and (C5) is tested at each complete leaf against the
literal disjunction. Surviving leaves are counted into $Z$, tallied into per-card
ownership counts, and tallied into a per-half-suit table keyed by the packed
assignment. Nothing is factorised, so the count behind every named allocation is
an explicit count of leaves.

That audit was run and reported by the v0.4 study, and this report does not
re-run it \cite{fishbot04}. Its two properties matter here and are restated
rather than re-measured. The deterministic comparisons agree \emph{exactly} ---
not to a tolerance --- because both sides compute integer counts over the same
finite set of assignments and hold them in double precision below $2^{53}$, form
the same ratio from identical operands, and obtain identical doubles. And its
coverage is a subsample: the search is exponential in $|U|$, so a state is
attempted only when the capacity dynamic program reports a small enough number of
consistent deals, and a state that exceeds the budget is counted as skipped and
reported rather than dropped. Skipping is correlated with exactly the property
that makes a state hard --- a large unresolved pool --- so the audit is evidence
about small reachable states and is silent about the high-entropy early game
where inference matters most. A pass is exact on a subsample, and that is all it
is.

Two v0.6 facts attach to this. The first is that the audit's premise did not hold
when it was run. \S\ref{sec:engine-repairs} shows that a \texttt{BlockDP}
instance's tables could be repointed by an unrelated build on the same thread;
the oracle comparison was unaffected because its driver builds one instance and
queries it before anything else can rebuild, which is a property of the driver
rather than of the engine. After the repair it is a property of the engine. The
second is that the exactness audit answers a narrower question than the phrase
suggests: it establishes that the block construction computes the posterior of
(C1)--(C5) correctly, and \S\ref{app:inf-prior} establishes that this posterior
is not the best available one.

The v0.6-era checks that were run are the reachability measurements of
\S\ref{app:inf-latent}.

\subsection{Latent approximations inside functions documented as exact}
\label{app:inf-latent}

Three routines in the inference stack are documented as exact and are not, in
regimes no current caller reaches. Each is recorded here with its reachability
argument, because unreachable-today is a property of the call sites and not of
the code, and every one of these is a trap for the next study rather than an
error in this one. The source is
\filepath{research/v06/notes/R2-inference-stack.md}, \S8.

\paragraph{A greedy branch in the count routine.} The masked counting routine in
\filepath{engine/src/belief.hpp} has a fast path for masks that are all equal and
one for singletons. For masks that are neither --- heterogeneous and
non-singleton --- it falls through to a greedy sequential assignment that picks
the seat with the most remaining capacity, which is an approximation inside a
function whose comment says exact. No current caller reaches it: over
\numCountUniformQueries{} uniform and \numCountGeneralQueries{} general queries
the routine agreed with the exact posterior to
$\numCountAgree{}$. A future caller asking a heterogeneous per-card question ---
``is each of these cards on my team?'' is the natural one --- would get a silent
approximation.

\paragraph{The allocation search does not re-check survival.} The block table's
maximum-probability allocation search materialises an assignment from a count
vector and returns it without re-testing it against the certificates of
\eqref{eq:disj}, whereas the allocation-probability query on the same table does
re-test and says why. The search is on the deployed decision path. It was
measured at \numAllocSurvivalFailures{} failures in \numAllocSurvivalChecks{}
checks, \numAllocSurvivalCert{} of them on states carrying a live certificate,
over \numAllocSurvivalDeals{} deals at seed \numAllocSurvivalSeed{}. That is
unfalsified and it is not proved. Corollary~\ref{cor:equalprob} does not supply
the missing step: it says survivors sharing a count vector are equiprobable, and
says nothing about whether the assignment the search materialises is one of them.

\paragraph{Silent truncation of a group table.} The group enumeration returns a
failure flag when the number of distinct count vectors exceeds
\numMaxEnt{}, but the caller that populates the table ignores the flag and
returns, leaving a \emph{partial} table which the team-ownership query then
normalises as if it were complete. The regime is currently unreachable by
arithmetic: the number of count vectors for six cards over six seats is
$\binom{11}{5} = \numMaxCountVectors{} < \numMaxEnt{}$, and no truncation was
observed in \numAllocSurvivalChecks{} checks. It becomes reachable the moment the
half-suit size or the table capacity moves. Two further caps have the same shape
and the same status: a \numCertCap{}-certificate limit per group, and
\numNodeGuard{}-node guards on the two depth-first searches.
%% --- end sections_v06/B-inference ---
%% --- begin sections_v06/C-parameters ---
\section{Configuration and fitted parameters}
\label{app:params}

This appendix records the frozen configuration in enough detail to reconstruct it
exactly. It is a reproducibility artifact and not a description of strategy: the
ask features are mutually correlated and normalised to different effective
ranges, so the sign and magnitude of a single coefficient do not identify the
contribution of the corresponding feature, and \S\ref{sec:fitting} explains why
the point applies with more force to a vector fitted by a repaired optimiser than
to one fitted by a broken one.

The appendix is arranged so that the freeze/parse round trip of
\S\ref{sec:fit-schedule} is checkable by a reader: the layout of the flat vector,
the printed values, the specification string that reproduces them, the bounds
applied on parse, and the provenance string that names the run the values came
from.

\subsection{The ask features}
\label{app:params-features}

The ask utility is a fitted linear score over \numAskFeats{} features, unchanged
in definition from the v0.4 and v0.5 policies. Table~\ref{tab:features} defines
them. All are normalised to roughly $[0,1]$ so the fitted weights are comparable
in scale; $H$ denotes the named card's half-suit and $p = \hat\mu_{c,t}$ the
approximate posterior that the target holds the named card, so every feature is
computed from the deployed inference path of \S\ref{sec:belief-count} rather than
from the exact one.

\begin{table}[ht]
\centering
\caption{The \numAskFeats{} ask features $\varphi_k$. Probabilities written
$\Pr[\cdot]$ are approximate per-card marginals. Features $3$, $5$, $7$ and $15$
are the ownership terms, and in \vfive{} and \vsix{} each is multiplied by an
ownership gate that is zero when the ask cannot succeed; that gate is the
mechanism the v0.5 study introduced to stop the ownership terms paying for an ask
of probability zero, and \S\ref{sec:mechanisms} measures what it costs.}
\label{tab:features}
\small
\begin{tabular}{rlp{0.58\linewidth}}
\toprule
$k$ & Name & Definition \\
\midrule
0  & hit probability     & $p = \hat\mu_{c,t}$, the posterior that the target holds the named card \\
1  & squared hit         & $p^2$; lets the fit curve the trade-off between certainty and value \\
2  & certain hit         & the indicator that $p$ exceeds the near-one threshold \\
3  & own progress        & cards of $H$ in hand, over six, gated \\
4  & team control        & expected fraction of $H$ held by our team \\
5  & lock completion     & $\prod_{d \in H \setminus \{c\}} \Pr[d \text{ is ours}]$: the probability this ask alone completes team ownership of $H$, gated \\
6  & continuation        & best posterior on any other missing card of $H$ \\
7  & completion bonus    & a step at four or more cards of $H$ in hand and a smaller one at three, gated \\
8  & reply threat        & $(1-p)\cdot$ the best ask the target could make against us if handed the turn \\
9  & information leak    & the indicator that our team has never publicly revealed an interest in $H$ \\
10 & target hand size    & target's public card count over nine \\
11 & empties the target  & $p$ if the target holds exactly one card, else zero \\
12 & repeats our half-suit & the indicator that this repeats our own previous half-suit \\
13 & known team cards    & publicly located cards of $H$ on our team, over six \\
14 & location entropy    & the binary entropy of $p$ \\
15 & team owns $H$       & $\prod_{d \in H}\Pr[d \text{ is ours}]$ before the ask, gated; an independence proxy, not the exact query \eqref{eq:q3} \\
16 & exposure on a miss  & $(1-p)\cdot$ our cards visible in half-suits the target has asked in \\
17 & trailing pressure   & $p$ when behind by two or more half-suits \\
18 & runway              & expected length of the run this ask begins, from the best remaining per-card posteriors \\
19 & leak magnitude      & feature 9 scaled by how much of $H$ we hold \\
\bottomrule
\end{tabular}
\end{table}

\vsix{} adds \numVsixAskTerms{} further ask terms. They are the mechanism of
\S\ref{sec:mech-tiebreak} and not a rescaling of the existing score:
\texttt{wVoid} rewards an ask whose hit takes the target's last card of
the half-suit, \texttt{wTeamHas} penalises an ask for a card the actor's own team
already holds, and \texttt{wLastLive} separates the case where the target is the
last live opponent. They enter the same linear term and are fitted as three
further coordinates rather than set by hand, and what each is worth is
reported in \S\ref{sec:results-ablations}.

\subsection{The frozen parameter vector}
\label{app:params-vector}

The v0.6 flat vector has \numVsixParams{} coordinates in one fixed order:

\begin{center}
\begin{tabular}{ll}
\toprule
Coordinates & Contents \\
\midrule
$0 \ldots \numAskFeats{}-1$ & the ask weights $w_0,\dots,w_{19}$ of Table~\ref{tab:features} \\
next \numVfiveKnobs{} & the v0.5 decision knobs, in the order of Table~\ref{tab:bounds} \\
last \numVsixAskTerms{} & \texttt{wVoid}, \texttt{wTeamHas}, \texttt{wLastLive} \\
\bottomrule
\end{tabular}
\end{center}

\noindent The offset at which each block begins is derived from the ask-feature
count and the v0.5 knob count in both the parser and the freeze step; it is not
hard-coded anywhere. That is not a stylistic preference. A hard-coded offset in
exactly this place aliased two ask weights onto the leading two knobs in the v0.4
study when the ask-feature count grew, and a hard-coded \texttt{18} survived in
the optimiser until \S\ref{sec:engine-repairs} removed it.

\begin{table}[ht]
\centering
\caption{The frozen \vsix{} configuration: \numVsixParams{} coordinates in the
order above. Generated by \texttt{engine/build\_tables\_v06.py} from the run
named in the provenance string of \S\ref{app:params-bounds}. This table is a
reproducibility record, not a description of strategy.}
\label{tab:params}
\footnotesize
%% --- begin tables_v06/params ---
\begin{tabular}{lrr}
\toprule
Coordinate & \vfive{} & \vsix{} \\
\midrule
\multicolumn{3}{c}{\textit{--- artifact not regenerated ---}} \\
\bottomrule
\end{tabular}
%% --- end tables_v06/params ---
\end{table}

The same vector can be supplied to any build of the engine as a single
\texttt{allparams=} string of \numVsixParams{} pipe-separated coordinates in that
order, which is the reconstruction key: a specification
\texttt{v06:allparams=\ldots} rebuilds the configuration from the numbers alone.
A shorter vector is accepted and the missing coordinates take their compiled
defaults, which is how a v0.5-length vector of \numVfiveParams{} coordinates
parses into a v0.6 build with the three new terms at zero --- and that, together
with \texttt{legacy=1}, is what makes the identity control of
\S\ref{sec:engine-repairs} expressible.

\subsection{Bounds and per-coordinate step sizes}
\label{app:params-bounds}

Bounds are applied when a sampled vector is turned into a configuration, by the
engine's parameter application in \filepath{engine/src/v06.hpp} and by the freeze
step in \filepath{engine/freeze_config_v06.py}, so a coordinate outside its bound
is clamped rather than rejected and the vector scored during fitting is the
vector compiled into the binary. The two integer coordinates are rounded to the
nearest integer after clamping.

\begin{table}[ht]
\centering
\caption{Bounds applied to each coordinate of the \numVsixParams{}-coordinate
vector on parse. The ask weights and the three v0.6 ask terms are unbounded; the
\numVfiveKnobs{} knobs are listed in the vector's own order. The per-coordinate
step size used in fitting is $\sigma_{\mathrm{rel}}$ times the width of the bound
(\S\ref{sec:fit-sigma}), so an unbounded coordinate takes the step of its
nominal search range rather than an unbounded one.}
\label{tab:bounds}
\small
\begin{tabular}{rll}
\toprule
Coordinate & Quantity & Bound \\
\midrule
ask weights & $w_0,\dots,w_{19}$ & unbounded \\
\midrule
1 & declaration threshold & $[0.5,\,0.9999]$ \\
2 & locked-half-suit threshold & $[0.5,\,0.99999]$ \\
3 & ask floor & $[0,\,0.9]$ \\
4 & patience pool & $[0,\,45]$, rounded to an integer \\
5 & opponent-card floor & $[0,\,20]$ \\
6 & value weight $\lambda_{\mathrm{val}}$ & $[0,\,\infty)$ \\
7 & linear weight $\lambda_{\mathrm{lin}}$ & $[0,\,\infty)$ \\
8 & minimum team probability & $[0.05,\,0.99]$ \\
9 & declaration margin & unbounded \\
10 & prior $\theta$ & $[0,\,2]$ \\
11 & prior $\phi$ & $[0,\,1]$ \\
12 & lookahead width $K$ & $[0,\,24]$, rounded to an integer \\
13 & chain weight $\lambda_{\mathrm{chain}}$ & $[0,\,\infty)$ \\
14 & threat weight $\lambda_{\mathrm{threat}}$ & $[0,\,\infty)$ \\
\midrule
v0.6 terms & \texttt{wVoid}, \texttt{wTeamHas}, \texttt{wLastLive} & unbounded \\
\bottomrule
\end{tabular}
\end{table}

The clip column is the evidence for \S\ref{sec:fit-sigma} and belongs in this
appendix per coordinate rather than as a single summary, because the defect it
diagnoses is a \emph{distribution} over coordinates and not a maximum: the
previous optimiser was bang-bang on the bounded knobs and frozen on the unbounded
ones at the same time, and a single worst-case figure hides the second half of
that. Every generation record of a fit trace carries the full per-coordinate clip
vector, so the table below is read directly out of the run file rather than
recomputed. At generation zero the worst coordinate of run~A clips at
\vsixFitAClipGenZero{}\% and of run~C at \vsixFitCClipGenZero{}\%.

\paragraph{Provenance.} \filepath{engine/freeze_config_v06.py} writes the
\emph{whole} flat vector into the compiled defaults and, next to it, a provenance
string recording the run file it came from, whether the values are that run's
final weights record or a generation's distribution mean, the first sixteen hex
digits of the run file's SHA-256, the generation count, the objective, the
pairing flag, the opponent panel, the fitting seed, the deals and rotations per
cell, and the relative step size. The field is never empty: an unfrozen build
carries a string that says so. This is the general form of a defect the v0.4
study shipped --- a configuration whose coefficients were not the ones any
recorded fit produced, because the freeze step wrote only part of the vector ---
and the provenance string closes it for everything the freeze step writes. What
it does not cover is named in \S\ref{app:params-value}.

\subsection{Value-function coefficients}
\label{app:params-value}

The value function that supplies the one-ply lookahead term and the stopping rule
is a linear score over sixteen state features: a bias and the score differential;
three control terms built from the expected fraction of each half-suit our team
holds, together with a sharpened version and the contested mass; the signed count
of half-suits already owned each way; counts of what is left to play for, from
the card differential and the size of the unresolved pool to the near-complete
half-suits leaning each way; and side to move with its interactions with control
and with the unresolved pool. Features are collected from all six seats at every
decision point rather than from the mover alone, because under mover-only
collection the side-to-move feature is constant and its coefficient would be
unidentified.

These sixteen coefficients are \emph{not} part of the \numVsixParams{}-coordinate
vector, are not searched by the optimiser, and are not written by the freeze
step. They are compiled defaults carried forward unchanged from the v0.4 study,
and the v0.4 study's own provenance gap --- that the compiled coefficients differ
numerically from the ones in its recorded ridge fit, so its reported
goodness-of-fit describes an artifact rather than the deployed vector ---
therefore stands in this report as well. It is recorded in
\S\ref{app:repro-gaps} rather than repaired, because repairing it would change
the configuration that produced every number here and would break the identity
control of \S\ref{sec:engine-repairs}.

\subsection{The policy-specification grammar}
\label{app:params-grammar}

Every agent in the engine is named by a specification string, parsed by
\texttt{makeAgent} in \filepath{engine/src/factory.hpp}. The grammar is

\begin{center}
\texttt{name} \quad or \quad \texttt{name:key=value,key=value,\dots}
\end{center}

\noindent where \texttt{name} selects the policy family and each
\texttt{key=value} pair overrides one field of that family's configuration. Keys
are order-independent, there is no whitespace, and \textbf{unrecognised keys are
ignored}. The last property is worth stating rather than tolerating: a misspelt
key fails silently and the default is used, so a specification always parses and
a typo is indistinguishable from an intentional default. This is how a v0.5
specification could carry \texttt{gateaudit=1} for an entire study while the
option was parsed only in the v0.4 branch (\S\ref{sec:engine-repairs}), and it is
why the run scripts of \S\ref{app:repro-commands} write every key explicitly even
when it equals the compiled default.

The keys used in this report fall into four groups. The inference keys are
\texttt{belief}, which selects the posterior path and accepts \texttt{fast} (the
deployed default), \texttt{block} (the exact count law of
\S\ref{sec:belief-count}), \texttt{exact}, \texttt{exactdisj}, \texttt{sinkhorn},
\texttt{indep} and \texttt{hybrid}, together with \texttt{souter},
\texttt{sinner} and \texttt{particles} for the solver. A caution attaches to
these: only \texttt{fast} calls the policy prior of \S\ref{sec:belief-prior} at
all, so every belief-mode row is confounded with a prior ablation, as
\S\ref{app:inf-prior} explains. The policy keys --- \texttt{decl},
\texttt{lockthr}, \texttt{minteam}, \texttt{askfloor}, \texttt{pool},
\texttt{oppfloor}, \texttt{vweight}, \texttt{lweight}, \texttt{vmargin},
\texttt{ptheta}, \texttt{pphi}, \texttt{topk}, \texttt{chain},
\texttt{threat} --- set the corresponding coordinates of
Table~\ref{tab:bounds}, and individual coefficients are addressable as
\texttt{w0}--\texttt{w19}, or in bulk as \texttt{weights=},
\texttt{vweights=} and \texttt{allparams=}. The instrumentation keys are
\texttt{gate} and \texttt{mgate} for the declaration pre-gates and
\texttt{gateaudit} for the audit of \S\ref{sec:engine-repairs}. The v0.6 keys are
listed in \S\ref{app:params-switches}.

\subsection{Mechanism switches and their shipped states}
\label{app:params-switches}

Table~\ref{tab:switches} lists one row per v0.6 mechanism: the switch that
enables it, its state in the shipped configuration, and the section that reports
the measurement which decided that state. A mechanism that is off is off because
a measurement said so, and the measurement is named; nothing here is off because
it was not tried.

\begin{table}[ht]
\centering
\caption{The v0.6 mechanism switches. ``Shipped'' is the state in the frozen
configuration of \S\ref{app:params-vector}. ``Fitted'' distinguishes a mechanism
whose parameters were coordinates of the fit from one run at a hand-set value.
Every row's evidence is a paired comparison at the budget
\S\ref{sec:protocol-resolution} requires.}
\label{tab:switches}
\small
\begin{tabular}{llll}
\toprule
Switch & Mechanism & Shipped & Evidence \\
\midrule
\texttt{legacy} & restore the v0.5 vector exactly & off & \S\ref{sec:engine-repairs} \\
\texttt{xf} & the \numVsixAskTerms{} v0.6 ask terms & on, fitted & \S\ref{sec:mech-tiebreak} \\
\texttt{wvoid} & void creation & fitted & \S\ref{sec:results-ablations} \\
\texttt{wteam} & team-ownership discount & fitted & \S\ref{sec:results-ablations} \\
\texttt{wlast} & last-live-opponent split & fitted & \S\ref{sec:results-ablations} \\
\texttt{extie} & exact-posterior tie resolution & off & \S\ref{sec:diag-irreducible} \\
\texttt{exactp} & exact posterior as the ask marginal & off & \S\ref{sec:belief-exactworse} \\
\texttt{s1} & determinized information-set search & off & \S\ref{sec:search-curse} \\
\texttt{kappa} & the paired deviation threshold & --- & \S\ref{sec:search-falsifier} \\
\texttt{maxq} & endgame-restricted search & --- & \S\ref{sec:search-falsifier} \\
\texttt{dead} & the deliberate miss, rationed & off & \S\ref{sec:search-deadask} \\
\texttt{deadsearch} & the deliberate miss as a search candidate & off & \S\ref{sec:search-deadask} \\
\texttt{gateaudit} & declaration pre-gate instrumentation & off & \S\ref{sec:engine-repairs} \\
\bottomrule
\end{tabular}
\end{table}

Two conventions in that table are deliberate. A switch shown as \texttt{---} is a
parameter of a mechanism that is itself off, so it has no shipped state; its row
exists because the sweep over it is a result even though the mechanism is not
deployed. And the certificate de-duplication flag of
\S\ref{sec:engine-repairs} is not in the table, because it is not a mechanism: it
is an engine repair whose default is off precisely so that it changes nothing.
%% --- end sections_v06/C-parameters ---
%% --- begin sections_v06/D-ablations ---
\section{Extended ablations}
\label{app:ablations}

\subsection{Design}
\label{app:abl-design}

Every ablation in this study is paired. A variant plays the \emph{same} deals against the \emph{same}
panel as the reference, and the statistic is the per-deal margin, bootstrapped resampling deals as
clusters because the rotations of one deal are a single correlated observation. This is the
`ablate' command, and it is the only estimator used for a mechanism claim anywhere in this paper.

The reason is in \S\ref{sec:protocol-paired} and is worth restating here, because it is the single
most consequential methodological fact this study produced. Five mechanisms cleared a $95\%$ paired
interval at four hundred to sixteen hundred games per cell --- the budget at which most published
ablation rows in this project's three prior studies were collected --- and returned \emph{exactly}
zero at three thousand, on two disjoint seed banks. Table~\ref{tab:abl-null} is that table.

\subsection{The mechanisms that vanished}
\label{app:abl-null}

\begin{table}[h]\centering
\caption{Five mechanisms, at the corpus's usual per-cell budget and at this study's. A negative
delta means the variant beats the reference. Every large-$n$ row is the result of record.}
\label{tab:abl-null}
\small
\begin{tabular}{lll}
\toprule
mechanism & small-$n$ paired delta & large-$n$, two disjoint banks \\
\midrule
void creation, $w=1.5$ & $-0.0292$ $[-0.0537, -0.0037]$ & $+0.0007$ $[-0.0213, +0.0227]$; $+0.0200$ $[-0.0017, +0.0413]$ \\
void creation, $w=3$ & $-0.0313$ $[-0.0593, -0.0033]$ & as above \\
delete the perseveration bonus & $-0.0200$ $[-0.0356, -0.0044]$ & $-0.0012$ $[-0.0137, +0.0112]$; $+0.0035$ $[-0.0092, +0.0158]$ \\
declaration margin $-0.06$ & $+1.75$ points unpaired & $+0.0133$ $[+0.0013, +0.0247]$ \\
delete the negative value-of-information term & $-0.0053$ $[-0.0244, +0.0141]$ & not separated \\
\bottomrule
\end{tabular}
\end{table}

Two readings, and the second is the one that generalises. The first is that none of these mechanisms
works. The second is that the evidence standard this project has used since v0.3 cannot separate a
real one-point effect from a seed draw, which means the design register the three prior studies
built --- and which this study inherited and set out to work through --- was partly assembled from
noise.

\subsection{The switched-off mechanisms}
\label{app:abl-off}

\vsix{} ships several mechanisms switched off. Each is listed with why, so that a later study does
not rebuild it:

\begin{itemize}
\item \textbf{Test-time search} (\S\ref{sec:search}). Off by default and shipped as the separate
  \vsixsearch{} configuration: it is the one mechanism here that beats a static rule, and it costs
  three orders of magnitude in throughput.
\item \textbf{Ownership features scaled by hit probability.} Measured at $-1.33$ points by the v0.5
  study once M1 removes the case the scaling was for.
\item \textbf{A blanket $(\text{card},\text{target})$ repetition guard.} Measured at $-6.13$ points
  --- the largest single-switch effect in this codebase --- because cards move, so a repeat after a
  public transfer is frequently the best ask on the board. The \emph{dead}-ask repeat guard v0.6
  uses is a different object and is provably free (\S\ref{sec:search-deadask}).
\item \textbf{Certificate de-duplication.} Correct, and it changes the approximate posterior, so it
  is opt-in and off by default to keep v0.4 and v0.5 bit-identical.
\end{itemize}

\subsection{Per-opponent detail}
\label{app:abl-peropponent}

\begin{table}[h]\centering
\caption{Paired ablations against \vsix{} over the six-opponent panel.}
\label{tab:abl-full}
%% --- begin tables_v06/ablations ---
\begin{tabular}{lrrr}
\toprule
Variant & Pooled & \vsix{} $-$ variant & 95\% paired CI \\
\midrule
\texttt{v05} & 65.73 & -2.69 & [-4.35, -0.98] \\
\texttt{v06:legacy=1} & 65.73 & -2.69 & [-4.35, -0.98] \\
\texttt{v06:xf=0} & 68.54 & +0.12 & [-1.67, +1.92] \\
\texttt{v06:s1=1,det=16,cand=6,kappa=2.0,maxq=26} & 69.77 & +1.35 & [-0.04, +2.79] \\
\bottomrule
\end{tabular}
%% --- end tables_v06/ablations ---
\end{table}

\subsection{The four-way that separates two confounded switches}
\label{app:abl-fourway}

A non-zero extra ask weight is what sets the flag that selects the v0.6 scoring path, and the v0.6
scoring path is the one that does \emph{not} run v0.5's top-$K$ chain and threat re-scoring.
Switching the extra terms off therefore switches that pass back on, so the obvious ablation moves
two things at once. Table~\ref{tab:fourway} runs all four cells.

\begin{table}[h]\centering
\caption{The four-way. A negative delta means the variant beats \vsix{}. All four are null.}
\label{tab:fourway}
%% --- begin tables_v06/fourway ---
\begin{tabular}{llrrr}
\toprule
Arm & Configuration & Pooled & \vsix{} $-$ variant & 95\% paired CI \\
\midrule
\texttt{v06:wvoid=0,wteam=0,wlast=0} & extras off, chain off & 68.88 & -0.46 & [-1.15, +0.23] \\
\texttt{v06:chain2=1} & extras on, chain ON & 69.17 & -0.75 & [-2.44, +0.92] \\
\texttt{v06:xf=0} & extras off, chain ON & 68.54 & -0.12 & [-1.96, +1.69] \\
\texttt{v06:rtie=1} & extras on, chain off, ties at random & 69.10 & -0.69 & [-2.54, +1.12] \\
\bottomrule
\end{tabular}
%% --- end tables_v06/fourway ---
\end{table}

Two things follow. The extra ask terms are worth nothing, and so is the chain pass --- which is why
\vsix{} drops it, and why \vsix{} is the fastest policy in the family
(\vsixThroughputSix{} games per second against \vfive{}'s \vsixThroughputFive{}). And resolving
the tie group at random is worth nothing too, which is the control that makes
\S\ref{sec:diag-irreducible}'s reading of the search result possible.

\subsection{Specification of every arm}
\label{app:abl-specs}

Arms are named by the configuration grammar of Appendix~\ref{app:params}. \texttt{v06:legacy=1} is
\vsix{} carrying v0.5's parameter vector and is bit-identical to \vfive{}; it is the arm that
isolates the refit from every other change. \texttt{v06:xf=0} removes the three extra ask terms.
\texttt{v06:s1=1,...} switches the search on at the stated budget.
%% --- end sections_v06/D-ablations ---
%% --- begin sections_v06/E-search-detail ---
\section{Test-time search: implementation detail}
\label{app:search}

\subsection{Sampling the deal}
\label{app:search-sampling}

The capacity dynamic program of \S\ref{sec:belief-count} is an exact sampler as well as an exact
counter: descending it with the backward table as the branch weights draws uniformly from the deals
satisfying constraints C1--C4. Certificates (C5) are enforced by rejection, exactly as the
reference posterior does, so an accepted draw is an exact sample of the policy-agnostic posterior.

The policy prior is then applied as an \emph{importance weight} rather than being folded into the
sampler. Two reasons. The sampler stays exact and testable against the brute-force enumerator; and
the prior stays a separable, auditable object, which matters because \S\ref{sec:belief-exactworse}
shows the prior is carrying more of the posterior's predictive value than the exactness is. The
weighted statistics use Kish's effective sample size, so the deviation rule's standard error is not
overstated when the weights are uneven.

Measured cost at a decision: the dynamic program builds in tens to hundreds of microseconds
depending on how much of the deal is still unresolved, and one accepted draw costs well under a
microsecond. Determinization is not the expensive part of this search; the rollout is.

\subsection{Reconstructing the six information sets}
\label{app:search-control}

The reconstruction rests on a structural fact about the deduction state: every write it performs is
a function of the public event alone, except the two lines that maintain the observer's own hand.
Six such objects fed the same event stream from an all-possible start therefore agree bit for bit
outside those two fields, so the public deduction state is \emph{one} object, replayed once per
decision, and a seat's information set is that object refined by the hand the determinization gives
it.

This is checked against ground truth rather than argued. Over \vsixReconStates{} states and \vsixReconChecks{} card
checks the reconstruction differs from the seat's real deduction state on \vsixReconMaskPct\% of
cards, and in \vsixReconWider{} of \vsixReconWider{} differing cases the reconstruction is \emph{wider}: it is a strict under-approximation.
No simulated player is ever credited with knowledge the real one lacks, which is the direction the
soundness argument needs.

The cost of reconstructing all six seats is a few microseconds, because the shared object is
replayed once and the six refinements are linear in the deck.

\subsection{The blueprint used inside the rollout}
\label{app:search-blueprint}

A rollout is only as informative as the policy that plays it. Table~\ref{tab:rollout} gives the
frontier. Two axes were available and the corpus had only ever used one of them:

\begin{itemize}
\item \textbf{Dropping the two-ply refinement} costs little strength and removes most of the
  per-decision cost, because that refinement rebuilds the approximate posterior twice per candidate.
\item \textbf{Reducing the Sinkhorn iteration count}, which the corpus had not tried, is the cheap
  axis. The deployed setting is four outer and eight inner sweeps; two and four is statistically
  indistinguishable from the full blueprint at a fraction of the cost, and one and one costs a few
  points.
\item \textbf{Replacing the belief with independent marginals} is $62\times$ cheaper per event and
  $38$ points weaker. It was the first blueprint tried and it is why the first search build scored
  in single digits.
\end{itemize}

\subsection{The deviation rule}
\label{app:search-guard}

Candidate $0$ is the blueprint's own choice. For each other candidate $r$ the search forms the
paired difference $d_r = \sum_d w_d\,(v_{r,d} - v_{0,d})$ over the shared determinizations, its
weighted variance, and a standard error using the effective sample size; it deviates to
$\arg\max_r\, (d_r - \kappa\,\mathrm{se}_r)$ only when that lower bound is positive.

Inside the blueprint's own tie group the penalty is waived, because there the blueprint expresses no
preference and ``deviating'' from whichever candidate the enumerator emitted first is not a
deviation. That waiver is load-bearing: applying the penalty inside the tie group as well takes the
win rate to \vsixSearchTieGuardedWin\% --- the advantage goes with it.

\subsection{The deviation rate, and why it has to be decomposed}
\label{app:search-kappa}

Table~\ref{tab:searchdev} is the diagnostic, and read naively it says the search is mis-scaled: it
moves \vsixSearchDevTwoFive\% of the decisions it searches at the best setting, and the rate never
falls below about $31\%$ however far $\kappa$ is raised.

The floor is the waived tie group. Applying the same penalty there collapses the rate to
\vsixSearchDevGuardedTwoFive\% at $\kappa = 2.5$ and \vsixSearchDevGuardedFour\% at
$\kappa = 4$, so the search's \emph{evidence-backed} deviation rate is inside the band the
literature associates with a refinement, and the $31\%$ floor is re-choosing among candidates the
blueprint could not separate. Since the win goes away when that waiver is removed
(\vsixSearchTieGuardedWin\%), the tie group is where the search is doing its work --- though
\S\ref{sec:search-falsifier} states why the attribution is not resolved at this sample size.
%% --- end sections_v06/E-search-detail ---
%% --- begin sections_v06/F-calibration ---
\section{Calibration}
\label{app:calibration}

\subsection{What is forecast}
\label{app:calib-ask}

Two probabilities are emitted and can be scored against outcomes. The ask forecast is the
posterior probability that the named target holds the named card, recorded at every ask by the
measured team. The declaration forecast is the policy's own probability that the allocation it is
about to name is correct, recorded at every voluntary declaration.

\subsection{Reliability}
\label{app:calib-decl}

See \filepath{research/v06/results/E6-calibration.txt} for the raw ten-bin reliability diagrams and
the Brier, log-loss and expected-calibration-error summaries for both forecasts.

\subsection{The calibration/resolution distinction}
\label{app:calib-versions}

\S\ref{sec:belief-exactworse} turns on a distinction that a single calibration number hides. A
forecaster can be well calibrated and poorly resolved: if it says $0.46$ and is right $46\%$ of the
time, it is calibrated, however uninformative $0.46$ is. The exact posterior and the deployed
approximation are \emph{both} well calibrated in this game --- the mean forecast of each tracks its
realised frequency closely --- and they differ in resolution, in the approximation's favour, because
the approximation carries a soft model of how opponents choose their asks and the exact object does
not. Any comparison of inference paths that reports only calibration will therefore report that they
are equally good, and be wrong.

\subsection{Caveats}
\label{app:calib-caveats}

Forecasts are collected from one team only, so the sample is not independent across seats within a
game; the intervals in the reliability table are indicative rather than inferential. The declaration
forecast is a probability under the policy's own posterior, not under the exact one, and
\S\ref{sec:belief} gives the size of the gap between them.
%% --- end sections_v06/F-calibration ---
%% --- begin sections_v06/G-reproducibility ---
\section{Reproducibility and artifact manifest}
\label{app:repro}

Every number in this report must be reachable from a command in this appendix.
That is the standard the v0.4 report set and this one matches it, with one
addition the earlier studies did not have: the numbers pipeline is arranged so
that a figure which has \emph{not} been regenerated is visibly wrong on the page
rather than plausibly stale (\S\ref{app:repro-build}).

\subsection{Source layout}
\label{app:repro-source}

The engine is a single C++20 translation unit assembled from headers in
\filepath{engine/src/}. The rules, the card and half-suit types and the action
encoding are in \filepath{engine/src/fish.hpp}. Constraint tracking (C1--C5) and
the approximate posterior are in \filepath{engine/src/belief.hpp}; the exact
count law of \S\ref{sec:belief-count} --- its partition function, its marginals,
its allocation query and the table-aliasing guard of
\S\ref{sec:engine-repairs} --- is in \filepath{engine/src/blockdp.hpp}; the
brute-force enumerator that audits it is in \filepath{engine/src/oracle.hpp}.

The policies are layered. \filepath{engine/src/v04.hpp} and
\filepath{engine/src/v05.hpp} hold the two preceding generations, and
\filepath{engine/src/v06.hpp} holds \texttt{V06Agent}, which derives from
\texttt{V05Agent} and defers to it when no v0.6 mechanism is enabled; the
determinized search's rollout machinery is in
\filepath{engine/src/v06_rollout.hpp}. The scripted opponent population is in
\filepath{engine/src/baselines.hpp}. The game driver, which owns turn order,
declaration rounds and the information-safety audit, is in
\filepath{engine/src/game.hpp}; the match runner and both bootstrap procedures
are in \filepath{engine/src/arena.hpp}; the optimiser of \S\ref{sec:fitting} is
in \filepath{engine/src/tuner.hpp}. The v0.6 diagnostics are in
\filepath{engine/src/probe_v06.hpp} and the aliasing probe in
\filepath{engine/src/probe_forcedendgame.hpp}. The specification parser is in
\filepath{engine/src/factory.hpp} and the entry point in
\filepath{engine/src/main.cpp}.

Build with \verb|make| in \filepath{engine/}. The build produces a single binary,
\filepath{engine/fish}.

\subsection{Commands}
\label{app:repro-commands}

The battery is \filepath{engine/experiments_v06.sh}, which runs sixteen blocks
labelled \texttt{E0}--\texttt{E15} and writes one artifact per block into
\filepath{research/v06/results/}. It is presented below in the order the script
executes them, which is not numerical order: the two gates run first by
construction. \texttt{-{}-games} counts \emph{deals}, not games, so every game
count in this report is the product of the deal count and the rotation count.

\begin{verbatim}
# ---- build -----------------------------------------------------------
cd engine && make

# ---- E0  GATE: identity controls -------------------------------------
#      md5 of the whole pathology block, v06 with every switch off
#      against v05; then the BlockDP aliasing probe on both policies.
#      Nothing downstream is collected if this fails.
./fish pathology --a=v06:legacy=1 --b=v06:legacy=1 --games=60 --seed=31
./fish pathology --a=v05         --b=v05          --games=60 --seed=31
./fish blockalias --a=v04 --games=25 --seed=31
./fish blockalias --a=v05 --games=25 --seed=31

# ---- E1  rules, information safety and belief soundness --------------
./fish verify --games=600

# ---- E2  GATE: the pathology KPIs ------------------------------------
./fish pathology --a=v06 --b=v06 --games=300 --rotations=2 --seed=31
./fish pathology --a=v05 --b=v05 --games=300 --rotations=2 --seed=31
./fish pathology --a=v04 --b=v04 --games=300 --rotations=2 --seed=31
./fish pathology --a=v06 --b=v05 --games=300 --rotations=2 --seed=90210

# ---- E3  head-to-head, five held-out banks ---------------------------
for S in 90210 31337 515151 777001 424242; do
  for OPP in v05 v04; do
    ./fish match --a=v06 --b=$OPP --games=300 --rotations=6 \
                 --seed=$S --json
  done
done

# ---- E4  per-opponent profile; the headline is the MINIMUM -----------
for OPP in v05 v04 v03 v02 lockout detective diversifier hunter \
           bluffer random silent feint withholder; do
  for A in v06 v05 v04; do
    ./fish match --a=$A --b=$OPP --games=300 --rotations=6 \
                 --seed=515253 --json
  done
done

# ---- E5  paired mechanism ablations ----------------------------------
./fish ablate --ref=v06 --panel=v05,v04,v03,lockout,detective,withholder \
    --games=400 --rotations=2 --seed=606060 \
    --variants="v05;v06:legacy=1;v06:xf=0;v06:s1=1,det=16,cand=6,kappa=2.0,maxq=26"

# ---- E6  calibration --------------------------------------------------
./fish calibrate --a=v06 --b=v05 --games=400 --seed=717171

# ---- E7  rule dialects ------------------------------------------------
for FLAG in "" --no-out-of-turn --no-cardless-declare --legacy; do
  ./fish match --a=v06 --b=v05 --games=250 --rotations=6 \
               --seed=828282 $FLAG --json
done

# ---- E8  exchangeable ties, and belief as a predictor -----------------
./fish v6probe --mode=ties --a=v06 --b=v06 --games=150 --seed=31
./fish v6probe --mode=ties --a=v05 --b=v05 --games=150 --seed=31
./fish v6probe --mode=ties --a=v05 --b=v03 --games=150 --seed=90210
./fish v6probe --mode=belief --a=v05 --b=v05 --games=120 --seed=31

# ---- E12 the search ladder: the curse, and what a guard recovers ------
#      row 1 is the control: the same code with the blueprint forced.
for CFG in s1=1,det=8,cand=6,blend=1000000 \
           s1=1,det=8,cand=6,kappa=0 \
           s1=1,det=12,cand=4,kappa=0 \
           s1=1,det=12,cand=4,kappa=1 \
           s1=1,det=12,cand=4,kappa=2.5 \
           s1=1,det=12,cand=4,kappa=4 \
           s1=1,det=16,cand=6,kappa=2,maxq=26 \
           s1=1,det=16,cand=6,kappa=2,maxq=26,deadsearch=2; do
  ./fish match --a="v06:legacy=1,$CFG" --b=v05 --games=120 \
               --rotations=6 --seed=90210 --json
done

# ---- E15 the deliberate miss: the gate's negative control -------------
./fish pathology --a=v05:m1=0,m1p=1 --b=v05:m1=0,m1p=1 \
                 --games=200 --rotations=2 --seed=31
./fish match --a=v05:m1=0,m1p=1 --b=v05 --games=200 --rotations=6 \
             --seed=90210

# ---- E14 search deviation rate ----------------------------------------
for K in 0 1 2.5 4 6; do
  ./fish v6probe --mode=search --a="v06:legacy=1,s1=1,det=12,cand=4,kappa=$K" \
                 --b=v05 --games=20 --seed=90210
done

# ---- E13 rollout-blueprint fidelity -----------------------------------
for R in "v05:belief=indep,topk=0" "v05:topk=0,souter=1,sinner=1" \
         "v05:topk=0,souter=1,sinner=3" "v05:topk=0,souter=2,sinner=4" \
         "v05:topk=0"; do
  ./fish match --a="$R" --b=v05 --games=200 --rotations=6 \
               --seed=90210 --json
  ./fish bench --a="$R" --b="$R" --games=100 --threads=1
done

# ---- E11 partner regimes ----------------------------------------------
for P in "" v03 detective withholder; do
  for A in v06 v05; do
    ./fish match --a=$A --b=v05 --partners="$P" --games=400 \
                 --rotations=2 --seed=313131 --json
  done
done

# ---- E9  throughput ----------------------------------------------------
for S in v04 v05 v06 "v05:topk=0" "v05:belief=indep,topk=0"; do
  ./fish bench --a="$S" --b="$S" --games=200
done

# ---- E10 declaration pre-gate audit ------------------------------------
#      dead code for v0.5 until this study; see sec:engine-repairs.
./fish gateaudit --a=v06:gateaudit=1 --games=400 --rotations=6 --seed=90210

# ---- the whole battery -------------------------------------------------
./experiments_v06.sh

# ---- exploitability: fit an exploiter, re-measure it on a fresh bank ----
#      v0.4 is probed first as a POSITIVE CONTROL, because its published
#      figure is known and the probe must reproduce it before any new
#      number from the same machinery is believed.
BASE=v05 TARGETS="v04 v05" ./exploitability_v06.sh

# ---- fitting: the two runs of record ------------------------------------
./fish tune --base=v05 --panel=v05,v03,lockout,withholder --full --paired \
            --obj=minimaxregret --sigmarel=0.05 --pop=20 --elite=5 \
            --games=150 --gens=30 --seed=20260823 \
            --out=../research/v06/runs/fitA.jsonl
./fish tune --base=v06 --panel=v05,v03,withholder,feint --full --paired \
            --obj=minimaxregret --sigmarel=0.04 --pop=14 --elite=4 \
            --games=200 --gens=14 --seed=20260824 \
            --out=../research/v06/runs/fitC.jsonl

# ---- harness acceptance test: recovery from a handicapped base ----------
#      NOT part of experiments_v06.sh.  The exact panel and seed of the
#      run reported in sec:fit-schedule are not recorded; see
#      app:repro-gaps.
./fish match --a=v05:w0=0 --b=v05 --games=300 --rotations=6 --seed=<bank>
./fish tune  --base=v05:w0=0 --full --paired --obj=minimaxregret \
             --sigmarel=0.05 --gens=10 --seed=<recorded in the trace> \
             --out=<trace>
./fish match --a="v05:allparams=$w" --b=v05 --games=300 --rotations=6 \
             --seed=<bank> --json
./fish match --a="v05:allparams=$w" --b=v03 --games=300 --rotations=6 \
             --seed=<bank> --json

# ---- freeze a fitted vector into the compiled defaults -------------------
python3 freeze_config_v06.py ../research/v06/runs/fitC.jsonl --dry
python3 freeze_config_v06.py ../research/v06/runs/fitC.jsonl
\end{verbatim}

The battery ran in \vsixWallClock{} hours on \vsixCores{} cores. Throughput
figures were measured under that load and are lower bounds for that reason
(\S\ref{sec:protocol-paired}).

\subsection{The build and the numbers pipeline}
\label{app:repro-build}

No number appears as a literal anywhere in \filepath{paper/sections_v06/}. Every
figure is a macro, and the macros come from a deliberately two-file scheme
introduced by the v0.5 study and kept here.

\filepath{paper/numbers_v06.tex} is the placeholder file. It declares every macro
the manuscript uses, with \verb|\providecommand| so that partial regeneration
cannot clash, and it splits them into two classes that are handled differently.
A macro naming a quantity the v0.6 battery produces carries a deliberately
impossible literal --- \texttt{XX.XX}, \texttt{X,XXX} --- so that a forgotten
regeneration is loud on the page rather than silently plausible. A macro naming a
quantity \emph{transcribed} from an artifact that already exists on disk, such as
a figure carried forward from the v0.4 or v0.5 record, carries its real value
under a comment naming the file it was read from, because leaving it as
\texttt{XX.XX} would make the sentence quoting it unreadable rather than visibly
unfinished.

\filepath{engine/build_tables_v06.py} reads
\filepath{research/v06/results/} and writes
\filepath{paper/numbers_v06_generated.tex}, which the manuscript inputs
\emph{after} the placeholder file, so a generated value always wins. Three
disciplines are enforced in the generator rather than by convention: every macro
carries a comment naming the artifact and the field it came from; a macro whose
artifact is missing is not emitted at all, so the placeholder survives and the
omission is visible; and the emission is \verb|\providecommand| followed by
\verb|\renewcommand|, so input order cannot change the result. The same script
writes the row bodies of the tables in \filepath{paper/tables_v06/}, and a table
whose artifact is absent is guarded at the \verb|\input| site so the manuscript
still compiles.

Two checkers back this up. \filepath{paper/check_provenance.py} verifies that
every transcribed macro sits under a comment naming an artifact and that the
artifact exists on disk, and reports the transcribed/generated split that
\S\ref{app:repro-manifest} quotes. \filepath{paper/check.py} enforces the
no-literal rule and a manuscript-specific list of banned phrasings; a v0.6
profile still has to be added to its profile table, and until it is, the
no-literal rule for this manuscript is enforced by reading rather than by
tooling. That is a gap and is listed as one.

The phrases banned for this manuscript are:

\begin{itemize}
  \item any wording that calls the deployed approximate posterior \emph{exact}.
        It is exact in its bookkeeping of (C1)--(C4) and approximate in its
        treatment of (C5), and \S\ref{app:inf-sinkhorn} says exactly where.
  \item any wording that treats the \emph{substitution} result --- swapping the
        exact posterior into a policy fitted for the approximate one --- as a
        verdict on exact inference. The verdict in this report rests on the
        predictive comparison of \S\ref{app:inf-prior}, which carries no fitted
        weights on either side; the substitution result is confounded and is
        reported as confounded.
  \item any claim of a first.
  \item any headline aggregate over the opponent panel. The headline is the worst
        case and the minimax regret (\S\ref{sec:protocol-metrics}).
\end{itemize}

The manuscript is built by \verb|npm run paper:v06|, which flattens the
\verb|\input| tree into \filepath{paper/fishbot_v06_standalone.tex} for
submission and then compiles \filepath{paper/fishbot_v06.tex} with
\texttt{tectonic} into \filepath{output/pdf/}.

\subsection{Artifact manifest}
\label{app:repro-manifest}

Table~\ref{tab:manifest} is generated by \filepath{engine/build_tables_v06.py} from
the artifact files themselves: the script walks
\filepath{research/v06/results/}, records each file's size and SHA-256, and pairs
it with the command that produced it. It is not maintained by hand.

\phantomsection\label{tab:manifest}\paragraph{Artifact manifest.} Artifact manifest for the experiments reported here. ``ID'' is the
experiment identifier used throughout the text, ``Artifact'' the file under
\texttt{research/v06/results/}, ``Design'' the population and experimental unit
recorded by the generator, and the last column the leading hex digits of the
file's SHA-256. The full digests, byte counts and exact command lines are in
\texttt{research/v06/results/MANIFEST.json}.

%% --- begin tables_v06/manifest ---
\begin{longtable}{lrl}
\toprule
Artifact & Bytes & sha256 (first 12) \\
\midrule
\endfirsthead
\toprule
Artifact & Bytes & sha256 (first 12) \\
\midrule
\endhead
\texttt{E0-identity.txt} & 594 & \texttt{d4228f6ac7fc} \\
\texttt{E1-verify.txt} & 119 & \texttt{1cf3b1e6f994} \\
\texttt{E10-gateaudit.txt} & 372 & \texttt{22a9cbda2a33} \\
\texttt{E11-partners.jsonl} & 4,374 & \texttt{b752dad1df7c} \\
\texttt{E12-search.jsonl} & 4,858 & \texttt{63b1e2954c4f} \\
\texttt{E13-rollout.jsonl} & 3,162 & \texttt{cb0e9e57e64d} \\
\texttt{E14-searchdev.txt} & 2,840 & \texttt{b7c28fc98676} \\
\texttt{E15-deliberate-miss.txt} & 2,127 & \texttt{1dbff1355e7f} \\
\texttt{E2-pathology.txt} & 2,928 & \texttt{4b0aec8137d4} \\
\texttt{E3-headtohead.jsonl} & 5,286 & \texttt{47cf9d834d4f} \\
\texttt{E4-perstyle.jsonl} & 20,486 & \texttt{154ad97702f1} \\
\texttt{E5-ablations.json} & 818 & \texttt{6761baedb9ce} \\
\texttt{E6-calibration.txt} & 891 & \texttt{8fb4ca32941c} \\
\texttt{E7-rules.jsonl} & 2,185 & \texttt{f9056ed0672e} \\
\texttt{E8-belief.txt} & 916 & \texttt{d892d62d8555} \\
\texttt{E8-ties.txt} & 3,686 & \texttt{adaf2d51e124} \\
\texttt{E9-throughput.txt} & 330 & \texttt{32a244f942a0} \\
\texttt{F0-search-confirm.jsonl} & 5,062 & \texttt{2aa98d49ff21} \\
\texttt{F1-chain2x2.json} & 816 & \texttt{d663f80c0f8a} \\
\texttt{F10-panel-banks.json} & 730 & \texttt{2663c7fd665d} \\
\texttt{F2-belief-noprior.txt} & 846 & \texttt{f5c2579cbd9b} \\
\texttt{F3-oracle.txt} & 700 & \texttt{94fab4303d6f} \\
\texttt{F3-selftest.txt} & 289 & \texttt{35936f9f0376} \\
\texttt{F4-perstyle-banks.jsonl} & 17,129 & \texttt{0c812e296d1d} \\
\texttt{F5-exploitability.log} & 1,962 & \texttt{3565de5dd586} \\
\texttt{F6-reconstruction.txt} & 588 & \texttt{0a6d8611301c} \\
\texttt{F7-tieguard.txt} & 1,229 & \texttt{000af9d5b4f9} \\
\texttt{F7-winrate.txt} & 235 & \texttt{2b70dcb2e7c9} \\
\texttt{F8-tiesearch.txt} & 530 & \texttt{684396fe47c5} \\
\texttt{F9-tieonly.txt} & 507 & \texttt{fae275343f67} \\
\texttt{MANIFEST.json} & 5,079 & \texttt{b7a30d0edc6a} \\
\texttt{X1-lbr-v04.jsonl} & 7,498 & \texttt{008b9f2625ac} \\
\texttt{X1-lbr-v05.jsonl} & 7,474 & \texttt{8a54fb672b99} \\
\texttt{X1-lbr-v06.jsonl} & 7,496 & \texttt{da5e5a201c9f} \\
\texttt{X1-lbr.jsonl} & 2,536 & \texttt{624d30143c6e} \\
\bottomrule
\end{longtable}
%% --- end tables_v06/manifest ---



\filepath{research/v06/results/MANIFEST.json} additionally records the commit the
manifest was generated at, whether the working tree was clean at that moment, the
shipped policy specification, the belief mode, the fitted-parameter count, the
ask-feature count, the fitting seeds and every evaluation bank of
\S\ref{sec:protocol-seeds}. Fitting traces are in
\filepath{research/v06/runs/}, and the authoritative record of which
configuration was frozen is the provenance string of
\S\ref{app:params-bounds}, which names the run file and its digest inside the
compiled header itself.

The manifest holds \vsixManifestFiles{} artifacts. Of the figures quoted in this
manuscript, \vsixProvGenerated{} are generated straight from an artifact and
\vsixProvTranscribed{} are transcribed from a report that already existed;
\S\ref{sec:limitations} states what that split means for how much of this report
is machine-checked.

\subsection{Known reproducibility gaps}
\label{app:repro-gaps}

Recorded rather than repaired, because repairing any of them would change a
configuration this report measured.

\begin{enumerate}
  \item \textbf{The harness acceptance run is not scripted.} The recovery test of
        \S\ref{sec:fit-schedule} --- a fit from the handicapped base
        \texttt{v05:w0=0} --- is not one of the blocks of
        \filepath{engine/experiments_v06.sh}, its panel and seed are not recorded
        in any committed script, and \filepath{engine/build_tables_v06.py} does
        not emit its figures. Those figures are therefore transcribed, not
        generated, and the run can be repeated but not reproduced exactly.
  \item \textbf{The value-function coefficients are outside the freeze.} The
        sixteen coefficients of \S\ref{app:params-value} are compiled defaults
        carried forward from the v0.4 study, are not part of the fitted vector,
        and are not written by \filepath{engine/freeze_config_v06.py}. The v0.4
        study's own gap --- that its compiled coefficients are not the ones its
        recorded ridge fit produced --- is inherited here unchanged.
  \item \textbf{The exactness audit was not re-run.} The brute-force oracle
        comparison of \S\ref{app:inf-audit} is carried forward from the v0.4
        study. Its coverage was a subsample of small reachable states then and no
        larger now, and the v0.6 battery adds no block that re-runs it.
  \item \textbf{Replay identity is tested at one thread count.} The random
        streams are keyed so that thread count cannot affect a game's outcome,
        but \texttt{fish verify} checks replay identity only at a fixed thread
        count and does not test that claim directly
        (\S\ref{sec:engine-verify}).
  \item \textbf{The no-literal rule is not yet tooled for this manuscript.}
        \filepath{paper/check.py} has profiles for the v0.4 and v0.5 manuscripts
        and not for this one, so the rule of \S\ref{app:repro-build} is currently
        enforced by review.
  \item \textbf{A raw shared-pool read still differs across an unrelated build.}
        The aliasing guard of \S\ref{sec:engine-repairs} makes every public query
        correct and leaves the underlying buffer pool shared, so
        \vsixAliasRawMismatch{} raw field reads in the probe still differ. Code
        that reaches into a \texttt{BlockDP}'s fields rather than calling its
        queries is unsupported, and nothing in the engine does.
\end{enumerate}
%% --- end sections_v06/G-reproducibility ---
%% --- begin sections_v06/H-human-play ---
\section{Implications for human play}
\label{app:human}

This project exists because a person plays this game and wanted a better opponent. This appendix
answers the questions that came from live play, in the order they were asked, and separates what
the evidence supports from what it does not.

\subsection{Does FishBot play for itself or for the team?}
\label{app:human-labels}

For the team, unambiguously, and this is a property of the objective rather than a tuning choice.
Every quantity the value function reads is team-relative: the score differential is the team's, the
card counts are the team's totals, and the per-half-suit control term is the expected fraction of
that half-suit held by the \emph{team}, not by the seat. There is no per-seat score in this game and
none in the policy. A configuration that maximised the acting seat's own holdings would be a
different program, and it is not the one that ships.

What is true, and is the substantive half of the question, is that the \emph{policy} is unilateral.
Through v0.5, FishBot maintained a deduction state for itself and for nobody else --- the only two
clones of that object anywhere in the policy are clones of its own --- so no term could ask what a
partner knows, or which of two teammates is better placed to act. Teammates entered only as
aggregate probability mass. v0.6 adds the machinery for per-seat modelling, uses it inside the
search, and reports that at the resolution of \S\ref{sec:protocol} the terms built on it are worth
nothing. That is a finding about this game rather than a concession: with every card movement
public, most of what a partner knows, you already know.

\subsection{Why a well-trained human still out-reasons it}
\label{app:human-ask}

The honest answer from this study is that the gap is in the \emph{action space}, not in the
arithmetic. M1, the mechanism that fixed v0.4's deadlock, restricts the candidate set to asks the
actor cannot prove will miss. That is the right fix for the deadlock and it deletes the move class
that most named human tactics are built from: blackballing, deliberately handing the turn to the
opponent who can do least with it, and paying for a signal with a safe ask are all a guaranteed miss
at a \emph{chosen} seat. Such a move is available at $79.2\%$ of decisions, and at two or more
distinct chosen opponents at $50.1\%$.

\S\ref{sec:search-deadask} reports what happened when it was given back. Unbanned wholesale, it
raises the win rate and destroys the policy. Rationed against the fitted linear score, it is never
chosen --- and the reason is structural rather than a matter of weights: the hit-probability term is
zero on a deliberate miss by construction, and every other term of the score is a penalty, so no
setting of the admission margin changes the output. Its value is the value of the position the
donated turn produces, which is a forward-looking quantity that no static feature computes. Offered
to the search, which is the only evaluator in the system that can price it, it produces the best
search configuration measured, and the search as a whole is the one mechanism in this study that
beats a static rule (\S\ref{sec:search}).

So: a human who blackballs is doing something the bot structurally cannot express, and the bot is
not obviously losing points by failing to. Both halves of that sentence are measured.

\subsection{Deception}
\label{app:human-bluff}

The manoeuvre the project owner reported --- holding cards of a half-suit you were asked for, then
declining to ask back in it --- is implemented as the \emph{withholder} archetype and is the one
opponent against which v0.6 is decisively better. The gain replicates in sign and size on two
disjoint seed banks and is the largest replicated effect in this study. It comes from the refit
rather than from an opponent model: the fitting panel contains the archetype and the objective is
minimax regret, so the optimiser buys robustness against it directly.

The complementary archetype, the \emph{feint}, manufactures a misleading ask-legality certificate
instead, and it remains the worst cell in the profile. v0.5's exposure there was created by its own
fit and could not be tuned away at fixed weights; v0.6's fit includes it in the panel and improves
it on one bank while giving a little back on the other.

\subsection{Playing against it}
\label{app:human-declare}

Three things a person should know when sitting down against this program.

\begin{itemize}
\item \textbf{It never asks a question it can prove will fail.} That is a hard deduction from your
  own public record, so a question it does ask is always one it cannot rule out --- which means an
  ask carries less information about its hand than a human's ask does.
\item \textbf{It is deterministic.} Given the same public history and the same hand it plays the
  same move. A person who has played many games against it can, in principle, read it. Adding
  stochasticity is the outstanding item: the machinery for a partner-aware temperature is specified
  and the harness now supports the two partner regimes as separate columns, but the temperature
  itself is not fitted here.
\item \textbf{Its declarations are conservative and accurate.} In mirror play it declares wrongly on
  roughly two percent of declarations, and the pre-gates that decide when to evaluate a declaration
  are exact --- \vsixGateFalseNeg{} false negatives in \vsixGateRejected{} rejections. If it
  declares, it is almost always right; the exploitable side is the timing, not the accuracy.
\end{itemize}

\subsection{Where a human should expect to keep an edge}
\label{app:human-ties}

At more than half of its decisions the program's own score cannot separate its leading candidates,
and \S\ref{sec:diag-irreducible} establishes that a correct posterior cannot separate them either.
A human facing the same information is in the same position. This is worth stating plainly because
it is the opposite of the usual intuition about card-game programs: the residual uncertainty here is
not a modelling deficiency that better software will remove --- it is the game.

Where a human should expect an edge is in the places this study found and could not close: choosing
\emph{which} opponent to hand the turn to when nothing good is available, paying a small material
price for a signal a partner will read, and varying behaviour so as not to be read. All three are
moves in the class M1 deletes.
%% --- end sections_v06/H-human-play ---
%% --- begin sections_v06/I-related-extended ---
\section{Extended related work}
\label{app:related}

This appendix records background that informed the design but that the main
argument does not depend on. Nothing here supports a result in
\S\ref{sec:results}; a reader who wants only the load-bearing literature should
read \S\ref{sec:related} and stop there. The organising convention is that each
entry says what the work does and then says which of three things this study
did with it: \textbf{uses} it as method, \textbf{cites} it as context, or
\textbf{declines} it. \S\ref{app:rel-declined} collects the declined set,
because in a study whose thesis is that the binding constraint is information
rather than machinery, the machinery left out is part of the argument.

\subsection{Further informal sources on the game}
\label{app:rel-informal}

Beyond Pagat \cite{pagat} and Develin \cite{develin} the descriptive corpus is
thin. Pagat's tactics list covers deliberate misses that inform a partner,
locking out a dangerous opponent by never asking them anything, and exhausting
one rank across all opponents before switching; it also records a
partner-signalling convention and a ``Forced Claims'' rule, both attributed to
named contributors. Develin's chapter contains the corpus's only quantitative
endgame discussion, in which the alternative to a coin-flip declaration would
need a success probability greater than one. The Wikipedia entry
\cite{wikiliterature} names the ``stalemate breaker'', and secondary strategy
pages \cite{dorsa,strategy-site} restate blackballing, asking the asker, and an
endgame delegation heuristic. These are used as evidence about how the game is
played and talked about, not as sources of quantitative claims. The two
substantial sources also disagree with each other on the point that matters most
to a study of conventions: Develin's chapter forbids pre-agreed conventions
outright, where Pagat records a partner-signalling convention as an accepted
tactic. That disagreement is why the partner regime is a declared parameter in
\S\ref{sec:mech-partner} rather than a fixed assumption.

On Somani's agent \cite{somani,somani-server}: it is a Python implementation
with a complete rules engine and a sound deduction fixed point, and the
four-player, 48-card variant it implements makes the declaration problem
structurally easier than the one studied here. With a single teammate, naming
the allocation of a six-card half-suit ranges over $2^6$ possibilities and is
frequently forced by the deduction fixed point; with two teammates it ranges
over $3^6 = 729$, and the allocation question stops being a formality. No
strength numbers have been published for it, so it is not usable as a baseline
in any case. The v0.4 report records a reading of its training code that
suggests the released model never received a win/loss gradient; that reading is
repeated here without re-verification, because nothing in this study depends on
it.

Valet \cite{goadrich2026valet} is the one piece of 2026 infrastructure that
could change the situation. It encodes twenty-one traditional
imperfect-information card games in a common description language and publishes
measured branching factors and game durations, which is exactly the external
calibration a claim like ``this game is hard'' needs. Whether any of the
twenty-one is a team game is not stated on its abstract page, and the full text
was not fetched, so no comparison is drawn anywhere in this report. Recorded so
that the next study fetches it rather than re-deriving the question.

\subsection{Determinization and perfect-information sampling}
\label{app:rel-determinization}

\S\ref{sec:related-determinization} states that averaging perfect-information
values over deals sampled from the public history collapses to a hit-probability
heuristic in this game. The collapse is the following. Almost every reachable
position carries the same perfect-information value, because the team on turn
takes everything it can reach, so with $t$ ranging over opponents, $c$ over
askable cards and $h$ the public history,
\[
  \arg\max_{(t,c)} \; \mathbb{E}\bigl[V^{\mathrm{PI}}(t,c)\bigr]
  \;\approx\;
  \arg\max_{(t,c)} \; \Pr[\,t \text{ holds } c \mid h\,].
\]
The right-hand side is one ply deep and blind to the information an ask leaks,
the information a refusal supplies and the option value of postponing a
declaration. This is a degeneracy of the perfect-information game itself, and it
is why the determinized search of \S\ref{sec:search} scores candidates by
rolling out under a belief-limited blueprint rather than by solving the sampled
world.

Determinization for bridge was proposed well before it was made to work
\cite{levy}, and constrained deal generation in practice is done by rejection:
deal against restrictions, then discard deals the program's own bidding module
would not have produced \cite{ginsberg-gib}. Frank and Basin's diagnosis
--- strategy fusion and non-locality --- is what makes it unsound however many
worlds are drawn \cite{frank-basin-ai}, and the repairs that work are those
attacking non-locality as well as fusion \cite{frank-basin-aaai}. Long,
Sturtevant, Buro and Furtak's leaf correlation, bias and disambiguation factor
explain when it nevertheless performs well \cite{long-pimc}. Information Set
Monte Carlo Tree Search replaces the independent trees with a single tree over
information sets, redeterminizing each iteration \cite{cowling-ismcts}, and has
been extended by re-determinizing at non-root seats \cite{goodman-ris}.
\textbf{Cited as context}; none is used as an engine here.

$\alpha\mu$ \cite{cazenave2019alphamu} deserves a separate note because it is
the closest thing in this literature to a repair this study could have adopted.
It backs up Pareto fronts of boolean world-vectors instead of scalar averages,
which represents ``this plan works in world set A, that plan works in world set
B'' without committing to a scalar mean --- structurally the same move as the
multi-valued leaf of \S\ref{sec:mech-leaves}, one ply down --- and its published
timings are an existence proof that vector-valued search over roughly twenty
worlds is affordable in about a second per move on a CPU. Its reported gain in
Bridge comes from differing with plain sampling on only a small fraction of
decisions, which is the shape a good search should have and the calibration
\S\ref{sec:search} uses when judging whether its own search is mis-scaled.
\textbf{Cited as context, and its diagnostic used}: the algorithm is declined
because it repairs unsoundness rather than degeneracy, and this game suffers the
latter.

Knowledge-based paranoia search \cite{edelkamp2021paranoia} is the closest
methodological neighbour this study has: it runs on a CPU with no network, is
built around deduced knowledge rather than sampled worlds, and reasons about the
worst case over consistent worlds rather than the average. All three match this
engine, whose deduction state carries hard exclusions, capacity constraints and
ask-legality certificates, and whose count law already computes the probability
of a named allocation. Worst-case-over-consistent-worlds is also the right
criterion for the one irreversible action in the game. \textbf{Its criterion is
used} in the declaration path of \S\ref{sec:mech-declare}; its game-tree search
is declined, because Skat's horizon and branching factor are not this game's. No
compute cost is attributed to it anywhere, because its runtime is not stated on
its abstract page.

\subsection{Depth-limited and subgame solving}
\label{app:rel-depthlimited}

Brown, Sandholm and Amos \cite{brown2018depthlimited} is the one result in this
group that is \textbf{used as method}. Its content is a single observation with
a large consequence: a search that stops at a depth limit and evaluates the leaf
with one value function has silently assumed the opponent plays one fixed
continuation, and that assumption is what makes depth-limited search in an
imperfect-information game unsound. Letting the opponent choose among $k$
continuations at the limit, each yielding its own leaf values, forces robustness
to all of them, and the demonstration ran on a four-core CPU and 16\,GB rather
than a supercomputer. The prerequisite it identifies as expensive --- a
portfolio of distinct continuation strategies --- is already paid for in this
project, because a set of compiled baseline policies of different styles exists
and is used for evaluation. The caveat is stated in \S\ref{sec:mech-leaves} and
repeated here: the soundness argument is two-player zero-sum, this game's leaf
has three uncorrelated opponent seats, and the construction is therefore adopted
as a robustification heuristic. A second caveat is that the method assumes a
blueprint approximating equilibrium, and this policy is a fitted heuristic, so
what the leaves buy is robustness against a style portfolio and not robustness
against equilibrium deviations.

Zhang and Sandholm's knowledge-limited subgame solving
\cite{zhang2021subgamesolvingwithoutck} is \textbf{cited as context and its
warning honoured}. Their motivation is precisely this game's blocker: current
techniques analyse the entire common-knowledge closure, which is acceptable when
that closure is small and impossible when it is not. Their remedy is to work
with low-order knowledge, pruning nodes no longer reachable on arrival at an
information set, and their honest result is that this can increase
exploitability, after which they develop the conditions that restore safety.
This engine already ships a crude instance of the same move --- a gate that
restricts candidate asks to those with strictly positive hard-consistent
probability --- whose measured gains were large and whose exploitability cost
has never been measured. \S\ref{sec:mech-exploit} exists because of this
paragraph.

Subgame solving in adversarial team games \cite{zhang2022teamsubgame} builds a
gadget from the team belief DAG and refines a blueprint by column generation,
running in polynomial time given a best-response oracle when the blueprint is
sparse. \textbf{Declined}, and recorded as a verified non-transfer: the sparsity
condition is the escape hatch, and a policy with support over roughly seventy
legal asks at each of roughly a hundred decisions per game is not sparse. It is
named explicitly because it is the most natural thing a reader will ask about.

The counter-strategy line is \textbf{cited for one warning}. Continual
depth-limited responses \cite{milec2021cdr} state plainly that existing robust
response methods do not compose with limited-look-ahead solving off the shelf,
and supply a fix; adapting beyond the depth limit \cite{milec2025abd} extends
the idea to counter-strategies in large games while staying robust against
rational opponents. Any future version of this policy that carries both a
depth-limited search and an online opponent model needs that composition
argument, and this report does not have one, because it builds only the first.
Look-ahead reasoning with a learned model \cite{kubicek2025lamir} is the right
direction and requires a learned abstraction, so it is deferred. Equilibrium
refinements in gadget games \cite{kubicek2026refinements} are not applicable
here at all --- no gadget game is solved anywhere in this study --- but the
finding transfers as a meta-lesson: when a subproblem has many optima that are
theoretically equivalent, the choice among them is worth more than half of the
exploitability, and that is a strong statement of the tie-breaking argument in
\S\ref{sec:diag-ties}.

The full-solver line is \textbf{cited as context and declined as
implementation}. DeepStack \cite{moravcik2017deepstack} introduced continual
re-solving from a range and a set of opponent counterfactual values; its
re-solve gadget requires one value per opponent information set, which at three
opponents holding nine cards each is not reachable on any CPU budget. ReBeL
\cite{brown2020rebel} treats the public belief state as the state and proves
that running the same procedure at test time is safe, which is the licence under
which \S\ref{sec:search} operates at all; its algorithm iterates over the belief
support and is therefore blocked by cardinality. Student of Games
\cite{schmid2023sog} grows the public tree during solving and identifies
Scotland Yard --- a public-observation, hidden-position game --- as the closest
published analogue to this information structure; the transferable piece is
architectural and cheap, namely that a candidate set should be grown where the
search is uncertain rather than fixed at a constant, and even that is not
implemented here.

\subsection{Policy-aware inference in trick-taking games}
\label{app:rel-inference}

The Skat line is the direct ancestor of this study's belief construction and is
\textbf{used}: reweighting the deals consistent with the public record by a
model of how each player would have acted, whether through a supervised
per-card location model \cite{solinas-supervised} or through reach probabilities
taken from a policy \cite{rebstock-inference}, and reported as worth $+217$
tournament points per 36 games as a standalone module \cite{buro-kermit}. The
policy prior of \S\ref{sec:belief-prior} is a two-parameter version of the same
idea. What \S\ref{sec:belief-exactworse} contributes to this line is a case
where the trade-off inverts far enough to be visible in win rate: exactness is
affordable here, and the exact object still loses, because it is the one that
cannot see behaviour.

History filtering --- constructing a history consistent with a public state ---
is FNP-complete in the worst case, and enumeration is polynomial only for sparse
public trees \cite{solinas-history}. \textbf{Cited as context}, and this game
escapes the problem entirely: its public record determines every card movement,
so the only object to be reconstructed is the initial deal, which is what makes
the exact count law of \S\ref{sec:belief-count} possible.

Belief modelling in cooperative games is \textbf{declined}. The Bayesian Action
Decoder made the public belief an explicit object of the policy
\cite{foerster2019bad}, and its independent-marginals factorisation is the part
that does not survive contact with fixed hand sizes and card conservation; the
repair is the constrained scaling and block enumeration this engine already
implements. Approximate policy iteration over common-payoff games
\cite{sokota2021capi} formalises factorised prescriptions and is a formalism
rather than a tool here. Learned Belief Search \cite{hu2021lbs} learns an
approximate belief because Hanabi's exact belief is intractable; this game's is
not, so learning one would be strictly worse than what exists. That is a real
structural advantage of this domain and is stated plainly in
\S\ref{sec:related-inference}: the belief half of the problem is solved, and
only the search half is open.

\subsection{Cooperative play, conventions, and deviations}
\label{app:rel-cooperative}

Hanabi is the reference benchmark for coordination through observable actions
\cite{bard-hanabi}, and is relevant because this game likewise has no side
channel. The Simplified Action Decoder showed that letting teammates condition
on the greedy rather than the exploratory action is a prerequisite for
conventions to survive exploratory training \cite{hu-sad}; Other-Play
demonstrated how much of a self-play score is convention rather than skill, with
the same agents scoring far worse when paired across independent training runs
\cite{hu-otherplay}; and Off-Belief Learning provides an explicit dial on
convention depth together with a uniqueness result that makes independent runs
interoperable \cite{hu2021obl}. Test-time search over a fixed blueprint provides
a further improvement without any training at all \cite{lerer2020sparta}. Joint
policy search, which scores simultaneous changes at several information sets ---
which is what inventing a convention requires --- improved a bridge bidding
system by $+0.63$ IMPs per board against a strong commercial program
\cite{tian2020jps}. All \textbf{cited as context}; none is trained here.

Two structural asymmetries separate this game from Hanabi and are worth stating
because they cut in opposite directions. The first is the audience. Hanabi has
no adversary, so every bit of signal is pure profit; here a leaked card location
is directly executable, because once it is public that a player holds a card,
any opponent holding another card of that half-suit can take it and keep the
turn. The Hanabi construction that most clearly fails to port is the modular
hat-guessing code, which works because each receiver sees every other receiver's
hand and can subtract it \cite{cox-hanabi}; here no player sees any hand. The
second asymmetry runs the other way. Off-Belief Learning's grounded level-one
play, which refuses to read conventional meaning into a partner's action, is an
approximation in Hanabi and is nearly exact here, because ask legality is a
rules-level certificate: a legal ask proves its asker holds another card of the
named half-suit, with no convention involved. The prediction that grounded play
is much closer to optimal in this game than in Hanabi is testable with a single
flag and no training, and it is the ablation of \S\ref{sec:results-ablations}
that separates hard certificate reading from the soft prior weights.

Self-explaining deviations \cite{hu2022sed} are the most on-topic piece of this
literature for the question of finding non-obvious plans, and are
\textbf{declined for this version with the reason recorded}. A self-explaining
deviation is an action that departs from what common understanding calls
reasonable, taken so that a teammate reasoning by theory of mind concludes that
circumstances must be abnormal; the associated algorithm produced the first
algorithmic finesse plays in Hanabi. The analogue here is exact and currently
unexploited: an ask certifies a holding, so an ask with low hit probability that
certifies something a teammate could not otherwise deduce is precisely such a
deviation, and the belief machinery to score one already exists. It is declined
in this study for two reasons. The information features of the shipped policy
are all negative --- it charges for information rather than paying for it --- so
the mechanism would have to be built rather than tuned; and a deviation of this
kind is a convention by another name, which means it is the mechanism most
likely to look excellent in self-play and fail against a human or an earlier
version. Any future attempt must be reported with its cross-play cell and gated
behind the partner regime of \S\ref{sec:mech-partner}.

\subsection{Human-compatible play and regularised search}
\label{app:rel-human}

piKL \cite{jacob2022pikl} regularises a search result toward a reference policy
with a divergence bound that shrinks as the regularisation weight grows, and
reports improved top-one agreement with human moves in chess and Go alongside
strong play. \textbf{Cited as context, and its shape borrowed.} The framing that
matters here is not human-likeness: it is that a KL anchor is a tunable,
measurable interpolation between a reference policy and an unconstrained search,
with the reference recovered exactly in the limit, so an ablation along that
axis has a guaranteed floor. In a team game the anchor also does real strategic
work, because a search that finds a clever line its own teammates will misread
is a net loss.

Human-regularised search and learning \cite{hu2022humanregularized} extends this
to coordination with real partners in three steps, of which the third ---
regularised search at test time to absorb partner distribution shift --- is the
published architecture for exactly the requirement that this policy behave
differently with a bot partner and a human one. \textbf{Cited as context};
there is no human corpus for this game, so it is not implemented. Cicero
\cite{fair2022cicero} is cited for the family rather than the architecture: a
seven-player mixed cooperative-competitive game whose planning core is an
anchored search over actions, and whose lesson is that in games with allies the
planner's job is to act well against a model of what the other players will
actually do rather than against an equilibrium. The ad-hoc human-AI coordination
challenge \cite{dizdarevic2025ah2ac2} supplies the negative instruction this
study follows: an agent trained with no human data at all outperformed a
best-response-to-behavioural-cloning baseline against human proxies, so a small
human corpus is not the way in.

\subsection{Team games and common information}
\label{app:rel-team}

The common-information approach \cite{nayyar2013common} reformulates a
decentralised control problem as a single-agent problem for a coordinator who
knows the common information and selects prescriptions --- maps from each
controller's local information to its action --- and shows that the resulting
problem is a POMDP. \textbf{Used as a correctness lens.} Two consequences are
load-bearing in this report. First, it names the design object: the public
deduction state is one shared thing, and what distinguishes seats is a
refinement of it, which is the mechanism of \S\ref{sec:mech-knowledge} and the
reason that mechanism is $O(1)$ rather than six-fold. Second, it explains why
teammates must run identical procedures: a prescription is agreed in advance,
and a seat that deviates by searching invalidates its teammates' belief updates.
The coordinator's own problem is not solvable at this scale --- a prescription
maps every nine-card hand to an ask --- so nothing here computes one.

The equilibrium-computation literature for teams is \textbf{cited to delimit
scope and declined in full}. Celli and Gatti establish the complexity of the
team-maxmin equilibrium with correlation and the worst-case cost of forgoing
correlation \cite{celli-gatti}; Farina and coauthors connect ex-ante team
collusion to extensive-form correlation \cite{farina-collusion}; the team belief
DAG and the two-player conversion give a zero-sum game whose equilibria are
realization-equivalent to a team-maxmin equilibrium with correlation
\cite{zhang-tbdag,carminati-tpi}, at a representation cost exponential in how
much private information distinguishes teammates inside a public state. That
exponent is the blocker: three teammates hold unrestricted nine-card hands, and
one seat's information set at deal time contains
$45!/(9!)^5 \approx 1.9 \times 10^{28}$ deals. Hidden-role games
\cite{carminati2024hiddenrole} are cited only to delimit scope: teams here are
public from the deal, so there is no role uncertainty and none of that machinery
applies. Nothing in this report is an equilibrium result, and no claim in it
should be read as one.

\subsection{Counting and constrained sampling}
\label{app:rel-counting}

Counting the deals consistent with a public history is a permanent computation.
Ryser's inclusion--exclusion formula is exact in $O(2^n n)$ time \cite{ryser},
which is unusable at $n = 54$, and Glynn's sign-vector formula is the other
exact method \cite{glynn}; the Jerrum--Sinclair--Vigoda chain is a fully
polynomial randomised approximation scheme \cite{jsv} whose constants make it
impractical at this scale. Matrix scaling is the usual practical approximation
and carries a deterministic strongly polynomial guarantee \cite{lsw}; a Sinkhorn
iteration of this kind is what the deployed approximate path uses. Sequential
importance sampling is the standard contingency-table alternative
\cite{chen-sis}, and is not used, because it can underestimate badly while
appearing to have converged. \textbf{Cited as context, with one used.}

The reason none of this is on the critical path is structural. The hidden state
is exactly the initial deal, each seat's share of it is a known constant number
of cards, and every ask-legality certificate is confined to a single half-suit,
so the counting problem factorises into blocks that can be enumerated exactly
within a half-suit and combined by a capacity dynamic programme. That is the
construction of \S\ref{sec:belief-count}, and it is why this study can ask
whether exactness helps rather than whether it is affordable.

\subsection{Evaluation methodology and variance reduction}
\label{app:rel-eval}

The v0.4 study declined variance reduction outright. This study takes the same
decision about fitted control variates and a different one about pairing, and
the distinction is worth stating because the two are often filed together.
Advantage-sum estimators \cite{zinkevich-divat} and AIVAT \cite{burch-aivat}
learn or hand-build a value function and subtract it, and their reported gains
are large in two-player settings and smaller in six-player ones, which is the
closer analogue here. They remain \textbf{declined}. AIVAT is unbiased for
\emph{any} value function, which is a constraint on its expectation rather than
on any realised estimate, and a control variate fitted by the same author on the
same agent using the same evaluation data is the pathology that literature
itself warns about.

What this study adopts instead is pairing on common random numbers: the same
deals played in both arms, with treatment and control alternated across seat
rotations, and the difference taken per deal
(\S\ref{sec:protocol-paired}, \S\ref{sec:fit-paired}). It costs nothing to
implement, introduces no fitted object, and removes the largest part of deal
variance --- the report quotes the realised width ratio against the unpaired
estimator rather than asserting a benefit. It is also the standard remedy rather
than a novel one: duplicate bridge replays each board with the same cards in the
same seats, and the computer poker competition added common seeds across
pairings for the same reason \cite{acpc}. The cost of the alternative is that a
fitted estimator would have to be validated, and validating it would consume
more games than pairing saves.

One further piece of methodology is named because it is frequently applied here
by analogy and should not be. Refining an abstraction does not monotonically
reduce exploitability \cite{waugh-abstraction}, which sounds like the general
form of \S\ref{sec:belief-exactworse}'s finding that substituting a more
accurate belief into a policy fitted around a less accurate one does not improve
it. The analogy is suggestive and the result does not transfer: this policy is
not an abstraction of a solved game, so neither the pathologies that literature
describes nor its guarantees apply. The finding here is measured on its own
terms and defended by a control, not inherited.

Rating and stopping are treated briefly. Bradley--Terry-style rating models
\cite{coulom-whr,herbrich-trueskill} are not invariant to redundancy in the
rated population and cannot represent cyclic dominance, both of which are real
failure modes when the rated
population is one training run's own checkpoints; maximum-entropy Nash averaging
over the antisymmetric logit matrix is the proposed alternative
\cite{balduzzi-reeval}. Ratings appear in this report only as a compact summary
of a fixed round-robin, never as a claim about a strategy space. And monitoring
an interval computed for a fixed sample size and stopping when it looks
favourable is not a test, which is why sequential designs on normalised Elo
exist \cite{nelo} and why the seed banks here are fixed in advance and reported
at their planned size (\S\ref{sec:protocol-seeds}).

\subsection{Work this study deliberately does not use}
\label{app:rel-declined}

The declined set, with the reason for each. These are design decisions, and
stating them is the point of this appendix.

\begin{itemize}
  \item \textbf{Full solvers over public belief states}
        \cite{moravcik2017deepstack,brown2020rebel,schmid2023sog}. Every one
        requires iterating or learning over a common-knowledge closure whose
        cardinality here is on the order of $10^{28}$. The closure is used as a
        sufficient statistic for evaluation and never as an index set for
        iteration. ReBeL's safety result is retained as the licence for
        test-time search; its algorithm is not.
  \item \textbf{The team belief DAG and subgame solving over it}
        \cite{zhang-tbdag,carminati-tpi,zhang2022teamsubgame}. Representation
        cost is exponential in the private information distinguishing teammates
        inside a public state, and the polynomial escape hatch is conditional on
        blueprint sparsity that this policy does not have.
  \item \textbf{Learned belief models} \cite{hu2021lbs,foerster2019bad}. The
        exact belief is computed and oracle-validated here, so an approximation
        would be strictly worse. This is the one place where this domain is
        easier than Hanabi rather than harder.
  \item \textbf{Neural policies and decision-time gradient updates}
        \cite{fickinger2021rlfinetuning}. There is no network training
        infrastructure in this project and none is warranted by the results:
        the two largest effects measured in this report are a tie-break and a
        statistical guard, neither of which is a capacity problem. The 2026
        lightweight-agent study reaches the same conclusion independently
        \cite{kelidari2026goldstandard}.
  \item \textbf{Hidden-role machinery} \cite{carminati2024hiddenrole}. Teams are
        public from the deal; there is no role uncertainty to model.
  \item \textbf{Gadget-game equilibrium refinement}
        \cite{kubicek2026refinements} and \textbf{learned-model look-ahead}
        \cite{kubicek2025lamir}. No gadget game is constructed here and no model
        is learned; both are recorded for their transferable lessons only.
  \item \textbf{Fitted control variates for evaluation}
        \cite{zinkevich-divat,burch-aivat}. Declined for the reason in
        \S\ref{app:rel-eval}; pairing on common random numbers is adopted
        instead.
  \item \textbf{Self-explaining deviations} \cite{hu2022sed}. Declined for this
        version rather than permanently, with the construction sketched in
        \S\ref{app:rel-cooperative} so that the next study can build it against
        a cross-play measurement rather than a self-play one.
\end{itemize}

Two exclusions are not in this list because they are not choices. There is no
equilibrium computation anywhere in this report, and no claim in it is an
equilibrium claim. And there is no human dataset, so every statement about human
play in \S\ref{app:human} is a statement about a small number of recorded games
and is scoped as such.
%% --- end sections_v06/I-related-extended ---

\end{document}
